\documentclass[11pt]{article}
\usepackage[margin=1.05in]{geometry}
\usepackage{amsmath,amssymb,amsthm}
\usepackage{mathrsfs}
\usepackage{booktabs}
\usepackage{graphicx}
\usepackage{microtype}
\usepackage[section]{placeins}
\usepackage[colorlinks,citecolor=blue,linkcolor=blue,urlcolor=blue]{hyperref}

\newtheorem{theorem}{Theorem}[section]
\newtheorem{lemma}[theorem]{Lemma}
\newtheorem{proposition}[theorem]{Proposition}
\theoremstyle{definition}
\newtheorem{conjecture}[theorem]{Conjecture}
\newtheorem{remark}[theorem]{Remark}

\DeclareMathOperator{\Tr}{Tr}
\DeclareMathOperator{\Var}{Var}
\newcommand{\R}{\mathbb{R}}
\newcommand{\E}{\mathbb{E}}
\newcommand{\Sph}{S^{d-1}}
\newcommand{\Kt}{K_{\mathrm{T}}}

% Every \num... macro below is written by refresh_p9.py from the run records.
% BEGIN AUTOTABLE macros
\newcommand{\numAdvAdvantage}{-0.88}
\newcommand{\numAdvBestLam}{10^{-5}}
\newcommand{\numAdvDeltaEff}{-0.10}
\newcommand{\numAdvRiskBest}{0.015}
\newcommand{\numAdvRiskGated}{0.018}
\newcommand{\numAdvRiskInterp}{0.0248}
\newcommand{\numAdvRiskMean}{1.00}
\newcommand{\numAdvRiskRidge}{0.0132}
\newcommand{\numBetaBernFour}{0.0572\pm0.0020}
\newcommand{\numBetaBernSixteen}{0.1354\pm0.0016}
\newcommand{\numBetaBernTwoFiveSix}{0.8104\pm0.0047}
\newcommand{\numBetaRadFour}{0.0573\pm0.0011}
\newcommand{\numBetaRadSixteen}{0.1356\pm0.0015}
\newcommand{\numBetaRadTwoFiveSix}{0.8155\pm0.0037}
\newcommand{\numBmomWorst}{1.4}
\newcommand{\numBurgAlpha}{2.02}
\newcommand{\numBurgExp}{0.78}
\newcommand{\numBurgFactor}{6}
\newcommand{\numBurgR}{0.63}
\newcommand{\numBurgRate}{0.72}
\newcommand{\numBurgTakeover}{100}
\newcommand{\numBurgTakeoverExp}{2}
\newcommand{\numCertMax}{0.76}
\newcommand{\numCertMin}{0.21}
\newcommand{\numCone}{0.1389}
\newcommand{\numControlWorstSe}{2.3}
\newcommand{\numCorrLanes}{31}
\newcommand{\numCorrPairs}{248}
\newcommand{\numCten}{0.3106}
\newcommand{\numCtenLong}{0.31060759756232}
\newcommand{\numDratioMax}{3.9}
\newcommand{\numDratioMin}{3.7}
\newcommand{\numEdgeMax}{2.036}
\newcommand{\numEdgeMin}{2.007}
\newcommand{\numEmitGainLanes}{15}
\newcommand{\numEmitPhysGain}{25.0}
\newcommand{\numEmitSqGain}{21.0}
\newcommand{\numEtaMax}{1.61}
\newcommand{\numEtaMin}{0.01}
\newcommand{\numExpBernA}{0.62}
\newcommand{\numExpBernB}{0.61}
\newcommand{\numExpBernC}{0.68}
\newcommand{\numExpBernD}{0.65}
\newcommand{\numExpRadA}{0.62}
\newcommand{\numExpRadB}{0.61}
\newcommand{\numExpRadC}{0.68}
\newcommand{\numExpRadD}{0.65}
\newcommand{\numGainemitYFour}{0.5}
\newcommand{\numGainemitYOne}{25.5}
\newcommand{\numGainemitYThree}{15.3}
\newcommand{\numGainemitYTwo}{22.1}
\newcommand{\numGainocoTwooTwo}{0.3}
\newcommand{\numGainocoTwoscoTwo}{1.3}
\newcommand{\numGainocoTwowcoTwo}{0.3}
\newcommand{\numGammaMax}{0.15}
\newcommand{\numGapBlocks}{20}
\newcommand{\numGapDims}{7}
\newcommand{\numGapInv}{8}
\newcommand{\numGapInvDiffMax}{0.012}
\newcommand{\numGapInvNmax}{1287}
\newcommand{\numGapInvRiskMax}{0.39}
\newcommand{\numGapPairs}{211}
\newcommand{\numGateCaptured}{99}
\newcommand{\numGateRiskBest}{3.3\times10^{-6}}
\newcommand{\numGateRiskGated}{3.5\times10^{-6}}
\newcommand{\numGateRiskMean}{1.5\times10^{-5}}
\newcommand{\numGateTheta}{0.99}
\newcommand{\numGateTrainMse}{1.1\times10^{-5}}
\newcommand{\numHardDmax}{120}
\newcommand{\numHardDmin}{20}
\newcommand{\numHermCone}{0.63166}
\newcommand{\numHermCtwelve}{0.33574}
\newcommand{\numHermErr}{2\times10^{-17}}
\newcommand{\numHermPr}{0.019483}
\newcommand{\numHermSh}{0.016546}
\newcommand{\numHermV}{0.43503}
\newcommand{\numHurtemitYFour}{32}
\newcommand{\numHurtemitYOne}{0}
\newcommand{\numHurtemitYThree}{6}
\newcommand{\numHurtemitYTwo}{0}
\newcommand{\numHurtocoTwooTwo}{31}
\newcommand{\numHurtocoTwoscoTwo}{4}
\newcommand{\numHurtocoTwowcoTwo}{28}
\newcommand{\numIdentWorst}{4\times10^{-12}}
\newcommand{\numInvMeanCvMax}{0.84}
\newcommand{\numInvMeanCvMin}{0.35}
\newcommand{\numInvMeanSpreadMax}{7}
\newcommand{\numInvMeanSpreadMin}{2}
\newcommand{\numInvMedCvMax}{0.07}
\newcommand{\numInvMedCvMin}{0.04}
\newcommand{\numInvMedSpreadMax}{26}
\newcommand{\numInvMedSpreadMin}{7}
\newcommand{\numKsRatioMax}{1.10}
\newcommand{\numKsRatioMin}{1.06}
\newcommand{\numKurtDtwentyfourNmax}{71{,}636}
\newcommand{\numKurtDtwentyfourNmin}{17{,}752}
\newcommand{\numLanesemitYFour}{5}
\newcommand{\numLanesemitYOne}{5}
\newcommand{\numLanesemitYThree}{5}
\newcommand{\numLanesemitYTwo}{5}
\newcommand{\numLanesocoTwooTwo}{5}
\newcommand{\numLanesocoTwoscoTwo}{1}
\newcommand{\numLanesocoTwowcoTwo}{5}
\newcommand{\numLbChi}{0.06}
\newcommand{\numLbEndsSe}{24}
\newcommand{\numLbExcA}{6.05}
\newcommand{\numLbExcB}{14.9}
\newcommand{\numLbExcC}{6}
\newcommand{\numLbExcPull}{0.07}
\newcommand{\numLbFitA}{1.990}
\newcommand{\numLbFitAse}{0.010}
\newcommand{\numLbFitB}{7.65}
\newcommand{\numLbFitBse}{0.32}
\newcommand{\numLbFitPull}{1.0}
\newcommand{\numLbMidPull}{0.30}
\newcommand{\numLbTwoA}{1.990}
\newcommand{\numLbTwoB}{7.66}
\newcommand{\numMomChiDof}{7}
\newcommand{\numMomChiTwoLast}{9.6}
\newcommand{\numMomDmax}{200}
\newcommand{\numMomDmin}{20}
\newcommand{\numMomFitB}{7.2}
\newcommand{\numMomFitBse}{0.7}
\newcommand{\numMomFitWorst}{0.2}
\newcommand{\numMomLastMfour}{2.037}
\newcommand{\numMomLastPull}{7}
\newcommand{\numNBernFour}{20}
\newcommand{\numNBernSixteen}{20}
\newcommand{\numNBernTwoFiveSix}{20}
\newcommand{\numNRadFour}{40}
\newcommand{\numNRadSixteen}{60}
\newcommand{\numNRadTwoFiveSix}{40}
\newcommand{\numOcoLanes}{11}
\newcommand{\numOcoMarginMaxFour}{2.2}
\newcommand{\numOcoNegFive}{10}
\newcommand{\numOofAtEmitRiseMax}{74}
\newcommand{\numOofAtEmitRiseMin}{11}
\newcommand{\numOofAtOcoLanes}{11}
\newcommand{\numOofAtOcoNeg}{10}
\newcommand{\numOofBestEmitMax}{1.0}
\newcommand{\numOofBestOcoMax}{0.3}
\newcommand{\numOofBestOcoMin}{-0.4}
\newcommand{\numOofBestYfourMax}{3.1}
\newcommand{\numOofBestYfourNeg}{4}
\newcommand{\numOofPhysNeg}{30}
\newcommand{\numOofPhysSignAgree}{31}
\newcommand{\numOofPhysWorseMax}{29}
\newcommand{\numOpRatioFirst}{1.75}
\newcommand{\numOpRatioLast}{1.99}
\newcommand{\numOverlapFirst}{0.71}
\newcommand{\numOverlapLast}{0.80}
\newcommand{\numPGainemitYFour}{0.5}
\newcommand{\numPGainemitYOne}{25.9}
\newcommand{\numPGainemitYThree}{22.2}
\newcommand{\numPGainemitYTwo}{27.1}
\newcommand{\numPGainocoTwooTwo}{0.2}
\newcommand{\numPGainocoTwoscoTwo}{0.6}
\newcommand{\numPGainocoTwowcoTwo}{0.2}
\newcommand{\numRatioMaxHalf}{2.1}
\newcommand{\numRatioMaxOne}{6.1}
\newcommand{\numRatioMaxTwo}{12.5}
\newcommand{\numRatioMinHalf}{1.7}
\newcommand{\numRatioMinOne}{5.5}
\newcommand{\numRatioMinTwo}{8.1}
\newcommand{\numRayC}{1.14}
\newcommand{\numRayMeasA}{0.1441}
\newcommand{\numRayMeasB}{0.3254}
\newcommand{\numRayMeasC}{0.0654}
\newcommand{\numRayOffA}{0.5}
\newcommand{\numRayOffB}{0.2}
\newcommand{\numRayOffC}{3.7}
\newcommand{\numRayPredA}{0.1448}
\newcommand{\numRayPredB}{0.3261}
\newcommand{\numRayPredC}{0.0678}
\newcommand{\numRecords}{223}
\newcommand{\numRedMedFirst}{1.0\times10^{-11}}
\newcommand{\numRedMedLast}{6.3\times10^{-9}}
\newcommand{\numShareAMax}{0.39}
\newcommand{\numShareAMin}{0.33}
\newcommand{\numShareBMax}{0.46}
\newcommand{\numShareBMin}{0.39}
\newcommand{\numShareDMax}{0.22}
\newcommand{\numShareDMin}{0.20}
\newcommand{\numSignAgree}{31}
\newcommand{\numSignTot}{31}
\newcommand{\numSlackMax}{3358}
\newcommand{\numSlackMin}{60}
\newcommand{\numSminRatioFirst}{0.815}
\newcommand{\numSminRatioLast}{0.967}
\newcommand{\numSminSdFirst}{0.04}
\newcommand{\numSminSdLast}{0.002}
\newcommand{\numSwRatioMaxSixtyFour}{2.8}
\newcommand{\numSwRatioMinSixtyFour}{2.4}
\newcommand{\numSwRatioOneTwoEight}{1.9}
\newcommand{\numTauHalfMax}{0.9994}
\newcommand{\numTauHalfMin}{0.9980}
\newcommand{\numTransferLanes}{0}
\newcommand{\numTransferPairs}{0}
\newcommand{\numTransferSlackMin}{4.1}
\newcommand{\numTthreeBlockDesigns}{36}
\newcommand{\numTthreeIdLfirst}{1}
\newcommand{\numTthreeIdLlast}{6}
\newcommand{\numTthreeIdSdFirst}{0.18}
\newcommand{\numTthreeIdSdLast}{0.004}
\newcommand{\numTthreeSdMax}{0.0012}
\newcommand{\numVgapMax}{1.2\times10^{-2}}
\newcommand{\numVhDmax}{44}
\newcommand{\numVhDmin}{10}
\newcommand{\numVhMaxHalf}{0.64}
\newcommand{\numVhMaxOne}{1.88}
\newcommand{\numVhMaxTwo}{3.88}
\newcommand{\numVhMinHalf}{0.54}
\newcommand{\numVhMinOne}{1.70}
\newcommand{\numVhMinTwo}{2.51}
\newcommand{\numVhRmaxMax}{12}
\newcommand{\numVhRmaxMin}{4.1}
\newcommand{\numVhRqFirst}{3.1}
\newcommand{\numVhRqLast}{2.1}
\newcommand{\numVmeanGapMax}{0.13}
\newcommand{\numYthreeNegA}{3}
\newcommand{\numYthreeNegB}{2}
\newcommand{\numYthreeNegC}{1}
% END AUTOTABLE macros

\title{Three problems in kernel methods: proofs and finite-size measurements}

\author{Yitzchak Shmalo\thanks{\raggedright Einstein Institute of Mathematics, The Hebrew University
of Jerusalem, Jerusalem, Israel. \texttt{yitzchak.shmalo@gmail.com}.}}

\date{16 September 2026}

\begin{document}
\maketitle

\begin{abstract}
Three questions about kernel methods are answered here, and a fourth is measured. For the Gaussian
kernel with a fixed bandwidth on the sphere, the minimum-norm interpolant of noisy labels is benign
along the sequence $d_\ell=2^{2^\ell}$, $n_\ell=2^{2^{2\ell}}$: its excess risk tends to zero in
probability, and almost surely along the sequence, for every fixed-degree spherical-polynomial target
(Theorem~\ref{thm:T2}). At the exponential sample size $n=2^d$ the same interpolant is obstructed: the
expected excess risk on pure noise is at least $\frac5{36}\sigma^2$ in the limit, for every fixed
bandwidth (Theorem~\ref{thm:T3}). For ridge regression on two augmented views of each sample through a
random feature map with an odd activation, the optimised objective and the clean prediction risk have
the same limits in law as those of a correlated Gaussian surrogate that keeps the shared latent
variable of the two views (Theorem~\ref{thm:T1}); the proof rests on an exact Gaussianisation of each
block conditional on the feature matrix. The hard edge of the degree-two Veronese matrix remains
open; the two lemmas that stand between its reduction and the hard edge are measured, and the
measurements point to a sharper statement, that the leave-one-out normal has a semicircular
spectrum. The proofs are accompanied by finite-size measurements of the quantities they turn on: the
leverage certificate of the lower bound, the harmonic dimension gap that separates the benign regime
from the obstructed one, and the leverage concentration behind the equivalence. Two further results
are evaluated on data: an exact same-sample risk identity for a kernel correction fitted on top of a
learned mean, on the EMIT and OCO-2 radiative-transfer emulation corpora, where its sign depends on the
component being emulated and the correction fitted to out-of-fold residuals gains nothing; and the
compressed-resolvent hypothesis behind growing-rank spiked singular vectors, whose compression error
along rays of fixed rank-to-size ratio approaches $\sqrt{r/n}$ plus a term of order $r/n$.
\end{abstract}

\section{Introduction}

Kernel ridge regression and its ridgeless limit are among the few learning methods whose
high-dimensional theory is close to complete, and several natural questions about them have stayed
open. This paper settles three of them and measures a fourth. The questions are stated precisely in
Section~\ref{sec:four}; they concern the minimum-norm interpolant of the Gaussian kernel on the
sphere at two very different sample sizes, and ridge regression on augmented random features.

Along the doubly exponential sequence $d_\ell=2^{2^\ell}$, $n_\ell=2^{2^{2\ell}}$ the interpolant is
benign for every fixed-degree spherical-polynomial target and every fixed bandwidth
(Theorem~\ref{thm:T2}, Section~\ref{sec:proofT2}); the mechanism is a dimension gap, polynomial in $d$
on both sides of the critical harmonic degree. At the exponential sample size $n=2^d$ the interpolant
of pure noise is obstructed: its expected excess risk is at least $\frac5{36}\sigma^2$ in the limit
(Theorem~\ref{thm:T3}, Section~\ref{sec:proofT3}); the proof is an expectation argument built on the
leverage of a block of harmonic degrees, with no concentration step. For two augmented views of each
sample passed through a random feature map with an odd activation, the optimised ridge objective and
the clean prediction risk have the same limits in law as those of a correlated Gaussian surrogate
(Theorem~\ref{thm:T1}, Section~\ref{sec:proofT1}); conditional on the feature matrix there is a Gaussian
block whose second moments match the actual block exactly, and the replacement argument then only
has to pay for the fluctuation of a $2\times2$ leverage matrix. The fourth question, the hard edge of
the degree-two Veronese matrix, is open; Section~\ref{sec:t4} states the reduction that isolates it
and measures the two lemmas that the reduction names.

Part~II measures, at sizes a computer can reach, the objects the proofs are made of. A proved constant
may be far from the truth: between $d=9$ and $d=13$ the quantity bounded in Theorem~\ref{thm:T3}
measures between $\numVhMinOne$ and $\numVhMaxOne$ times $\sigma^2$ at $h=1$ and between
$\numVhMinHalf$ and $\numVhMaxHalf$ at $h=1/2$ (Section~\ref{sec:t3}), so the theorem is not tight, and
the measurement says by roughly how much. A lemma in a reduction can be checked before it is proved:
both lemmas that the reduction of Target~4 names look true at finite size, one of them in a stronger
form than the reduction needs, and the other with its constant identified as the fourth moment of the
semicircle law (Section~\ref{sec:t4}). An exact identity can be confronted with data directly:
Section~\ref{sec:corr} does this for a same-sample risk identity on two radiative-transfer emulation
corpora, where it separates the components on which a residual kernel correction helps from the one
on which it hurts. Finite $d$ cannot settle a $\liminf_{d\to\infty}$, and the sample sizes of
Theorem~\ref{thm:T2} leave the reachable world at $\ell=3$; where a table rests on a single seed or a
single corpus, the text says so.

\part{Theorems and proofs}

\section{The problems}\label{sec:four}

Throughout, $X_1,\dots,X_n$ are i.i.d.\ uniform on $\Sph$, $K$ is the Gaussian kernel
\[
K(x,y) = \exp\!\Big(-\frac{\|x-y\|^2}{2h^2}\Big) = e^{-\gamma}e^{\gamma\langle x,y\rangle},
\qquad \gamma = h^{-2},
\]
with $h>0$ fixed, $N_{d,s}$ is the dimension of the degree-$s$ spherical harmonics, $P_s$ is the
zonal harmonic normalized so $P_s(1)=1$ \cite{AtkinsonHan}, and $\lambda_{d,s}$ are the Mercer
eigenvalues \cite{MinhNiyogiYao}, so that
$K(x,y)=\sum_s \lambda_{d,s}N_{d,s}P_s(\langle x,y\rangle)$ and $\tau_s = N_{d,s}\lambda_{d,s}$
satisfies $\sum_s \tau_s = 1$.

\begin{description}
\item[Target 1 (correlated Gaussian equivalence).] Two augmented views of each sample,
$x_{ij}=\sigma(F^\top(g_i+\Delta z_{ij}))$, $j=1,2$, with a linear teacher $t_i=\xi^\top g_i$ and the
ridge objective $L(w)=\frac1{4n}\sum_i\sum_j(t_i-x_{ij}^\top w)^2+\frac\lambda2\|w\|^2$. Does the
optimized objective, and the clean test risk of its minimizer, agree with those of a matched
Gaussian surrogate? This is a correlated form of the Gaussian equivalence of random-feature models
\cite{Goldt,HuLu,MeiMontanari}.
\item[Target 2 (benign overfitting).] Along $d_\ell=2^{2^\ell}$, $n_\ell=2^{2^{2\ell}}$, with fixed
$h$, Gaussian noise and any fixed-degree spherical-polynomial target of bounded $L^2$ norm, does the
conditional excess risk of the minimum-norm interpolant tend to zero in probability? Benign
interpolation with kernels is a high-dimensional phenomenon
\cite{LiangRakhlin,RakhlinZhai,Bartlett,Ghorbani,Mallinar}; the question is whether it survives a
bandwidth that does not scale with the dimension.
\item[Target 3 (the obstruction).] With $n=2^d$ and zero signal, is
$\liminf_{d\to\infty}\E\,\mathcal{E}_{2^d,d}\ge c(h)\sigma^2$ for an explicit $c(h)>0$?
\item[Target 4 (the Veronese hard edge).] With $p=d(d+1)/2$ and
$z(x)=\big((x_a^2-1)/\sqrt2\big)_a\oplus(x_ax_b)_{a<b}$, $x\sim N(0,I_d)$, does $p\,s_{\min}(Z)^2$
for the square $p\times p$ matrix $Z$ with rows $z(X_i)$ obey the real square Wishart hard edge, whose
limit law is $1-e^{-\sqrt t-t/2}$ \cite{Edelman,TracyWidom}?
\end{description}

Theorems~\ref{thm:T2}, \ref{thm:T3} and~\ref{thm:T1} below answer Targets~2, 3 and~1. The proof of
Target~2 also covers every target in the empirical low block below the critical degree, and this
enlargement is sharp as a uniform bounded-$L^2$ class statement. Target~4 is open; Section~\ref{sec:t4}
gives its reduction. Two further results established in \cite{PN24} are evaluated on data in
Part~II: an exact identity for the change in population risk when a residual kernel correction is
fitted on the same sample as the mean it corrects (Section~\ref{sec:corr}), and the
compressed-resolvent hypothesis behind the theory of growing-rank spiked singular values and vectors
\cite{BGN,BBP} and its consequences for pruning (Section~\ref{sec:prune}).

% ===== Part I, section 1: setting and harmonic preliminaries; section 2: Target 3 =====
\section{Harmonic preliminaries}\label{sec:prelim}

Throughout, $X,X_1,\dots,X_n$ are independent and uniform on $\Sph$, $d\ge3$, and
\[
K(x,y)=\exp\Big(-\frac{\|x-y\|^2}{2h^2}\Big)=e^{-\gamma}e^{\gamma\langle x,y\rangle},\qquad \gamma=h^{-2},
\]
with $h>0$ fixed. We write $\mathcal H_{d,s}$ for the space of spherical harmonics of degree $s$ on
$\Sph$, $N_s=N_{d,s}=\dim\mathcal H_{d,s}$, and $\Pi_s(x,y)=\sum_{a=1}^{N_s}\psi_{s,a}(x)\psi_{s,a}(y)$
for its reproducing kernel with respect to the uniform probability measure, so that
$\Pi_s(x,y)=N_sP_s(\langle x,y\rangle)$ with $P_s$ the zonal harmonic normalised by $P_s(1)=1$
\cite{AtkinsonHan}. The addition theorem gives $\Pi_s(x,x)=N_s$, and orthogonality of distinct degrees
gives, for every $x,z\in\Sph$,
\begin{equation}
\E_Y\big[\Pi_s(x,Y)\Pi_t(Y,z)\big]=\delta_{st}\,\Pi_s(x,z),\qquad \E_Y\Pi_s(x,Y)=0\quad(s\ge1).
\label{eq:repro}
\end{equation}
The Mercer expansion of the kernel is $K(x,y)=\sum_{s\ge0}\lambda_s\Pi_s(x,y)$ with $\lambda_s=\lambda_{d,s}>0$
\cite{MinhNiyogiYao}; since $K(x,x)=1$, the trace masses $\tau_s=N_s\lambda_s$ satisfy $\sum_s\tau_s=1$.
For a sample we write $K=(K(X_i,X_j))_{ij}$ for the Gram matrix, $\Pi_s(X,X)=(\Pi_s(X_i,X_j))_{ij}$, and
$K_s=\lambda_s\Pi_s(X,X)$, so that $K=\sum_sK_s$. Each $\Pi_s(X,X)=\Phi_s\Phi_s^\top$ with
$\Phi_s=(\psi_{s,a}(X_i))_{i,a}\in\R^{n\times N_s}$, and $\|\Phi_s\|_F^2=\sum_i\Pi_s(X_i,X_i)=nN_s$.

\begin{lemma}[harmonic multiplicities]\label{lem:mult}
For $d\ge3$ and $s\ge0$,
\[
N_{d,s}=\binom{d+s-1}{s}-\binom{d+s-3}{s-2}=\frac{(2s+d-2)\,(s+d-3)!}{s!\,(d-2)!},
\]
\[
\rho_s:=\frac{N_{d,s+1}}{N_{d,s}}=\frac{(d+2s)(d+s-2)}{(s+1)(d+2s-2)} .
\]
The ratio $\rho_s$ is at least one and is non-increasing in $s$. For $s\le d$ one has
$\tfrac34\binom{d+s-1}{s}\le N_{d,s}\le\binom{d+s-1}{s}$, and for $s^2=o(d)$,
$N_{d,s}=\frac{d^s}{s!}\,e^{O(s^2/d)}$.
\end{lemma}

\begin{proof}
The closed forms are classical \cite{AtkinsonHan}. From the product form,
$\rho_s=\frac{d+s-2}{s+1}\cdot\frac{d+2s}{d+2s-2}$; both factors are non-increasing in $s$, and
$\frac{d+s-2}{s+1}\ge1$ for $d\ge3$. The ratio of the two binomial coefficients is
$\frac{s(s-1)}{(d+s-1)(d+s-2)}\le\big(\frac{s}{d+s}\big)^2\le\frac14$ for $s\le d$, which gives the two-sided
bound. For $s^2=o(d)$, $\binom{d+s-1}{s}=\frac{d^s}{s!}\prod_{j=1}^{s-1}(1+j/d)=\frac{d^s}{s!}e^{O(s^2/d)}$
and the subtracted term is smaller by the factor $O(s^2/d^2)$.
\end{proof}

\begin{lemma}[Mercer eigenvalues of the Gaussian kernel]\label{lem:mercer}
With $(d)_{m}:=d(d+2)\cdots(d+2m-2)$ and $(d)_0=1$,
\begin{equation}
\lambda_{d,s}=e^{-\gamma}\sum_{r\ge0}\frac{\gamma^{s+2r}}{2^r\,r!\,(d)_{s+r}}
=e^{-\gamma}\,\Gamma\!\big(\tfrac d2\big)\Big(\frac2\gamma\Big)^{d/2-1}I_{s+d/2-1}(\gamma),
\label{eq:mercer}
\end{equation}
where $I_\nu$ is the modified Bessel function of the first kind. Consequently, for all $s\ge0$,
\begin{equation}
\frac{\lambda_{d,s+1}}{\lambda_{d,s}}\le\frac{\gamma}{d+2s},\qquad
\lambda_{d,s}\ge\frac{e^{-\gamma}\gamma^s}{(d)_s},\qquad
\tau_s\le e^{-\gamma}\sum_{m\ge s}\frac{\gamma^m}{m!}.
\label{eq:mercerbounds}
\end{equation}
In particular the eigenvalues are strictly decreasing in $s$ as soon as $d>\gamma$.
\end{lemma}

\begin{proof}
Write $K(x,y)=e^{-\gamma}\sum_m\frac{\gamma^m}{m!}\langle x,y\rangle^m$. The Funk--Hecke formula
\cite{AtkinsonHan} gives the degree-$s$ coefficient of $t\mapsto t^m$ on $\Sph$; equivalently, for
$G\sim N(0,I_d)$ and the homogeneous harmonic extension $H_s(ru)=r^sY_s(u)$ of $Y_s\in\mathcal H_{d,s}$,
the identity $\E[H_s(G)e^{t\langle x,G\rangle-t^2/2}]=t^sH_s(x)$ and the comparison of the coefficient of
$t^{s+2r}$ give $\E[\langle x,G\rangle^{s+2r}H_s(G)]=\frac{(s+2r)!}{2^rr!}H_s(x)$, while
$\E\|G\|^{2(s+r)}=(d)_{s+r}$. Dividing out the radial moments yields the series in \eqref{eq:mercer};
the Bessel form is the series $I_\nu(\gamma)=\sum_r\frac{(\gamma/2)^{\nu+2r}}{r!\,\Gamma(\nu+r+1)}$
\cite[10.25.2]{DLMF} rewritten with $\nu=s+d/2-1$. The first two bounds in \eqref{eq:mercerbounds} are
termwise: the $r$-th term of $\lambda_{s+1}$ equals $\frac{\gamma}{d+2(s+r)}$ times the $r$-th term of
$\lambda_s$, and $\frac{\gamma}{d+2(s+r)}\le\frac{\gamma}{d+2s}$; the second is the $r=0$ term. For the
third, expand $\langle x,y\rangle^m=\sum_{s\le m}c_{m,s}P_s(\langle x,y\rangle)$ with $c_{m,s}\ge0$
(the Gegenbauer coefficients of a monomial are non-negative) and $\sum_sc_{m,s}=1$ by evaluating at
$\langle x,y\rangle=1$; hence $\tau_s=e^{-\gamma}\sum_{m\ge s}\frac{\gamma^m}{m!}c_{m,s}\le
e^{-\gamma}\sum_{m\ge s}\frac{\gamma^m}{m!}$.
\end{proof}

\begin{lemma}[hypercontractivity at growing degree]\label{lem:hyper}
For every $Y_s\in\mathcal H_{d,s}$, every $s\ge0$ and every $p\ge2$,
$\|Y_s\|_{L^p(\Sph)}\le(p-1)^{s/2}\|Y_s\|_{L^2(\Sph)}$. In particular, with $r_s:=\Pi_s/N_s$ and $z\in\Sph$
fixed, $\|r_s(X,z)\|_{L^{2q}}\le(2q-1)^{s/2}N_s^{-1/2}$ for every $q\ge1$.
\end{lemma}

\begin{proof}
Let $G\sim N(0,I_d)$ and $H_s(G)=\|G\|^sY_s(G/\|G\|)$, a homogeneous harmonic polynomial of degree $s$
and hence an element of the $s$-th Gaussian chaos. Nelson's hypercontractive inequality
\cite{Nelson,Janson} gives $\|H_s(G)\|_p\le(p-1)^{s/2}\|H_s(G)\|_2$. The radius and the direction of $G$
are independent, so $\|H_s(G)\|_p=\|\,\|G\|^s\|_p\,\|Y_s\|_{L^p}$ and $\|H_s(G)\|_2=\|\,\|G\|^s\|_2\,\|Y_s\|_{L^2}$,
and $\|\,\|G\|^s\|_p\ge\|\,\|G\|^s\|_2$ for $p\ge2$. Dividing gives the first claim; the second follows
from $\|r_s(\cdot,z)\|_{L^2}^2=\Pi_s(z,z)/N_s^2=1/N_s$ by \eqref{eq:repro}.
\end{proof}

\section{Target 3: the exponential-sample lower bound}\label{sec:proofT3}

Here $n=2^d$, the signal is zero, $Y_i=\varepsilon_i$ with $\varepsilon\sim N(0,\sigma^2I_n)$, and
$\hat f(x)=k(x)^\top K^{-1}\varepsilon$ is the minimum-RKHS-norm interpolant, $k(x)=(K(x,X_i))_i$. The Gaussian
kernel is strictly positive definite, so $K$ is invertible almost surely. The excess risk is
$\mathcal E_{n,d}=\E_{X_0}[\hat f(X_0)^2]$, and we write $V(X)=\E_\varepsilon[\mathcal E_{n,d}\mid X]$.

\begin{theorem}[Target 3]\label{thm:T3}
For every fixed $h>0$,
\[
\liminf_{d\to\infty}\ \frac{\E\,\mathcal E_{2^d,d}}{\sigma^2}\ \ge\ \frac5{36},
\]
the expectation being over the design and the noise.
\end{theorem}

The proof has three steps: an exact reduction of $V(X)$ to the leverage of a block of harmonic degrees,
a lower bound on the expected leverage of a block, and the location of the last block whose dimension
does not exceed $n$.

\begin{lemma}[reduction to a block]\label{lem:reduction}
Let $G=\E_{X_0}[k(X_0)k(X_0)^\top]=\sum_s\lambda_sK_s$, $t_s=\Tr(K^{-1}K_s)$ and $q_s=\Tr(K^{-1}K_sK^{-1})$.
Then $V(X)/\sigma^2=\Tr(K^{-1}GK^{-1})=\sum_s\lambda_sq_s$, $\sum_st_s=n$, $t_s^2\le nN_s\lambda_sq_s$ for
every $s$, and for every finite set $B$ of degrees, with $T_B=\sum_{s\in B}t_s$ and $D_B=\sum_{s\in B}N_s$,
\begin{equation}
\frac{V(X)}{\sigma^2}\ \ge\ \frac{T_B^2}{nD_B}.
\label{eq:blockreduction}
\end{equation}
\end{lemma}

\begin{proof}
$\E_\varepsilon\hat f(x)^2=\sigma^2k(x)^\top K^{-2}k(x)$, so $V(X)/\sigma^2=\Tr(K^{-1}GK^{-1})$, and by
\eqref{eq:repro},
\[
G_{ij}=\sum_{s,t}\lambda_s\lambda_t\,\E_{X_0}\big[\Pi_s(X_i,X_0)\Pi_t(X_0,X_j)\big]=\sum_s\lambda_s^2\Pi_s(X_i,X_j),
\]
i.e.\ $G=\sum_s\lambda_sK_s$; hence $V/\sigma^2=\sum_s\lambda_sq_s$, and
$\sum_st_s=\Tr(K^{-1}K)=n$. Since $K_s=\lambda_s\Phi_s\Phi_s^\top$, we have
$t_s=\lambda_s\langle K^{-1}\Phi_s,\Phi_s\rangle_F$ and $q_s=\lambda_s\|K^{-1}\Phi_s\|_F^2$, so by
Cauchy--Schwarz
\[
t_s\le\lambda_s\|K^{-1}\Phi_s\|_F\,\|\Phi_s\|_F=\sqrt{\lambda_sq_s}\,\sqrt{nN_s}.
\]
Summing over $B$ and applying Cauchy--Schwarz once more,
$T_B\le\sqrt{nD_B}\big(\sum_{s\in B}\lambda_sq_s\big)^{1/2}\le\sqrt{nD_B\,V(X)/\sigma^2}$.
\end{proof}

\begin{lemma}[expected leverage of a block]\label{lem:leverage}
Let $d>\gamma$ and $L\ge0$, and put $D=\sum_{s\le L}N_s$, $\lambda=\lambda_{d,L}$, $\tau_{>L}=\sum_{s>L}\tau_s$
and $T_{\le L}=\sum_{s\le L}t_s$. Then
\begin{equation}
\E\,T_{\le L}\ \ge\ \frac{nD}{D+n-1+\tau_{>L}/\lambda}.
\label{eq:leverage}
\end{equation}
\end{lemma}

\begin{proof}
Put $A=\sum_{s\le L}K_s$, $C=K-A=\sum_{s>L}K_s\succeq0$, and $A_w=\lambda\sum_{s\le L}\Pi_s(X,X)$. By
Lemma~\ref{lem:mercer} the eigenvalues decrease in $s$, so $\lambda_s\ge\lambda$ for $s\le L$ and
$0\preceq A_w\preceq A$. For $\epsilon>0$ let $K_\epsilon=K+\epsilon I$ and $C_\epsilon=A_w+C+\epsilon I$; then
$K_\epsilon\succeq C_\epsilon\succ0$, so $K_\epsilon^{-1}\preceq C_\epsilon^{-1}$ by operator monotonicity of the
inverse. Since $A+C+\epsilon I=K_\epsilon$,
\[
\Tr(K_\epsilon^{-1}A)=n-\Tr\big(K_\epsilon^{-1}(C+\epsilon I)\big)\ \ge\ n-\Tr\big(C_\epsilon^{-1}(C+\epsilon I)\big)
=\Tr(C_\epsilon^{-1}A_w),
\]
where the inequality uses $C+\epsilon I\succeq0$. Cauchy--Schwarz for the Hilbert--Schmidt inner product,
applied to $C_\epsilon^{-1/2}A_w^{1/2}$ and $C_\epsilon^{1/2}A_w^{1/2}$, gives
$(\Tr A_w)^2\le\Tr(C_\epsilon^{-1}A_w)\Tr(C_\epsilon A_w)$. By the addition theorem $\Tr A_w=n\lambda D$
pathwise. Letting $\epsilon\downarrow0$ (the Gram matrix is invertible almost surely),
\begin{equation}
T_{\le L}=\Tr(K^{-1}A)\ \ge\ \frac{(n\lambda D)^2}{\Tr(A_w^2)+\Tr(CA_w)}.
\label{eq:pathwise}
\end{equation}
The denominator has exact expectations. By \eqref{eq:repro}, for independent $X,Y$,
$\E[\Pi_{\le L}(X,Y)^2]=\sum_{s\le L}\E_Y\Pi_s(X,Y)^2=\sum_{s\le L}\Pi_s(X,X)=D$ where
$\Pi_{\le L}=\sum_{s\le L}\Pi_s$, while $\Pi_{\le L}(X,X)=D$; hence
$\E\Tr(A_w^2)=\lambda^2\big(nD^2+n(n-1)D\big)$. Likewise $\E[C(X,Y)A_w(X,Y)]=0$ for $X\ne Y$ because the
two kernels expand over disjoint sets of degrees, while $C(X,X)=\tau_{>L}$ and $A_w(X,X)=\lambda D$; hence
$\E\Tr(CA_w)=n\tau_{>L}\lambda D$. No independence between $C$ and $A_w$ is used here: both are functions of
the same design, and the off-diagonal expectation vanishes by orthogonality alone. Since $z\mapsto1/z$ is
convex on $(0,\infty)$, Jensen's inequality applied to \eqref{eq:pathwise} gives
\[
\E\,T_{\le L}\ \ge\ \frac{(n\lambda D)^2}{\lambda^2(nD^2+n(n-1)D)+n\tau_{>L}\lambda D}
=\frac{nD}{D+n-1+\tau_{>L}/\lambda}. \qedhere
\]
\end{proof}

\begin{lemma}[the last subcritical block]\label{lem:window}
Let $n=2^d$ and $L_-=L_-(d)=\max\{L:\sum_{s\le L}N_{d,s}\le n\}$, $D_-=\sum_{s\le L_-}N_{d,s}$ and
$y_d=D_-/n\in(0,1]$. Let $c_*\in(0,1)$ be the root of $H(c)=\log2$, where
$H(c)=(1+c)\log(1+c)-c\log c$. Then $L_-/d\to c_*$, $c_*>1/4$, and
\begin{equation}
\liminf_{d\to\infty}y_d\ \ge\ \frac{c_*}{1+c_*}\ >\ \frac15,
\qquad
\frac{\tau_{>L_-}}{n\,\lambda_{d,L_-}}\ \longrightarrow\ 0 .
\label{eq:lastblock}
\end{equation}
\end{lemma}

\begin{proof}
$H$ is continuous and strictly increasing on $(0,\infty)$ ($H'(c)=\log\frac{1+c}c>0$), with $H(0^+)=0$ and
$H(1/4)=\frac54\log5-2\log2<\log2$ (numerically $0.6255<0.6931$), so $c_*$ exists and $c_*>1/4$.
By Stirling's formula, $\frac1d\log\binom{d+cd}{cd}\to H(c)$ uniformly for $c$ in compact subsets of
$(0,\infty)$. Fix $c<c_*$: by Lemma~\ref{lem:mult}, $\sum_{s\le cd}N_{d,s}\le(cd+1)\binom{d+cd}{cd}
=e^{d(H(c)+o(1))}<2^d$ for large $d$, so $L_-\ge cd$. Fix $c\in(c_*,1)$: $N_{d,cd}\ge\frac34\binom{d+cd-1}{cd}
=e^{d(H(c)+o(1))}>2^d$ for large $d$, so $L_-<cd$. Hence $L_-/d\to c_*$, and by Lemma~\ref{lem:mult}
$\rho_{L_-}=\frac{(d+2L_-)(d+L_--2)}{(L_-+1)(d+2L_--2)}\to\frac{1+c_*}{c_*}=:r_*$, with $r_*<5$ because
$c_*>1/4$.

For the first claim in \eqref{eq:lastblock}: $n<D_-+N_{L_-+1}$ by maximality of $L_-$, so
$y_d>\big(1+N_{L_-+1}/D_-\big)^{-1}$. Fix an integer $J$. For $1\le j\le J$ and $d$ large,
$\rho_{L_--j}=\rho_{L_-}\prod_{i=1}^{j}\frac{\rho_{L_--i}}{\rho_{L_--i+1}}\le\rho_{L_-}(1+1/(L_--J))^{2J}$,
because, from the product form of $\rho_s$ in Lemma~\ref{lem:mult},
$\frac{\rho_{s-1}}{\rho_s}=\frac{d+s-3}{d+s-2}\cdot\frac{s+1}{s}\cdot\frac{(d+2s-2)^2}{(d+2s)(d+2s-4)}
\le\big(1+\tfrac1s\big)\big(1+\tfrac{4}{(d+2s)(d+2s-4)}\big)\le\big(1+\tfrac1s\big)^2$ for $s\ge2$.
Writing $\bar\rho=\rho_{L_-}(1+1/(L_--J))^{2J}$,
$D_-\ge\sum_{j=0}^{J}N_{L_--j}\ge N_{L_-}\sum_{j=0}^{J}\bar\rho^{-j}$, while $N_{L_-+1}=\rho_{L_-}N_{L_-}$.
Since $L_-\to\infty$, $\bar\rho\to r_*$, and therefore
$\liminf_dy_d\ge\big(1+r_*\big/\textstyle\sum_{j=0}^J r_*^{-j}\big)^{-1}$ for every $J$; letting
$J\to\infty$, $\liminf_dy_d\ge\big(1+r_*(1-1/r_*)\big)^{-1}=1/r_*=c_*/(1+c_*)>1/5$.

For the second claim, by Lemma~\ref{lem:mult}, for $s\ge L_-\ge d/10$ (large $d$),
$\rho_s\le2\,\frac{d+s-2}{s+1}\le22$, and by Lemma~\ref{lem:mercer}, $\lambda_{s+1}/\lambda_s\le\gamma/d$.
Hence for $j\ge1$, with $L=L_-$,
$\frac{\tau_{L+j}}{n\lambda_L}=\frac{N_{L+j}}{n}\cdot\frac{\lambda_{L+j}}{\lambda_L}
\le\frac{N_L}{n}\,(22\gamma/d)^j\le(22\gamma/d)^j$, using $N_L\le D_-\le n$. Summing over $j\ge1$ gives
$\tau_{>L}/(n\lambda_L)\le 44\gamma/d$ for $d>44\gamma$.
\end{proof}

\begin{proof}[Proof of Theorem~\ref{thm:T3}]
Since $\gamma=h^{-2}$ is fixed, $d>\gamma$ for all large $d$, so Lemmas~\ref{lem:mercer}
and~\ref{lem:leverage} apply. Take $B=\{0,\dots,L_-\}$ in Lemma~\ref{lem:reduction}, so $D_B=D_-$ and
$T_B=T_{\le L_-}$. Taking
expectations in \eqref{eq:blockreduction} and applying Jensen's inequality, $\E[T^2]\ge(\E T)^2$, then
Lemma~\ref{lem:leverage},
\[
\frac{\E\,\mathcal E_{2^d,d}}{\sigma^2}=\E\frac{V(X)}{\sigma^2}\ \ge\ \frac{(\E T_{\le L_-})^2}{nD_-}
\ \ge\ \frac{nD_-}{\big(D_-+n-1+\tau_{>L_-}/\lambda_{d,L_-}\big)^2}
=\frac{y_d}{\big(1+y_d-\tfrac1n+\eta_d\big)^2},
\]
with $\eta_d=\tau_{>L_-}/(n\lambda_{d,L_-})\to0$ by Lemma~\ref{lem:window}. The function
$g(y)=y/(1+y)^2$ is increasing on $(0,1]$, since $g'(y)=(1-y)/(1+y)^3\ge0$. Given $\epsilon>0$,
Lemma~\ref{lem:window} gives $y_d\ge\frac15-\epsilon$ for all large $d$, and $y_d\le1$ always, so for large
$d$ the right side is at least $g(\tfrac15-\epsilon)-o(1)$. Letting $d\to\infty$ and then $\epsilon\to0$,
$\liminf_d\E\,\mathcal E_{2^d,d}/\sigma^2\ge g(1/5)=\frac{1/5}{(6/5)^2}=\frac5{36}$.
\end{proof}

\begin{remark}\label{rem:sharper}
The argument gives more than the stated constant. Lemma~\ref{lem:window} shows $\liminf y_d\ge1/r_*$
with $1/r_*=c_*/(1+c_*)=0.2271\ldots$, so the same proof yields the constant $g(1/r_*)=0.1509\ldots$;
the constant $5/36$ is the value at the round number $1/5$. The bound is an expectation statement from
start to finish: no concentration of the leverage is used, and none is needed. The certificate that the
proof evaluates, $(T_{\le L}/n)^2/(D_{\le L}/n)$, and the block dimension ratio $D_{\le L}/n$ are exactly the
quantities measured in Section~\ref{sec:t3}, where the one-window bound is compared with the realised
$V(X)/\sigma^2$ at $d\le13$. Running the same argument over several blocks simultaneously improves the
constant; the ten-window certificate $c_{10}=3882594969529/12500000000000=0.3106\ldots$ obtained in
\cite{PN24} is the value against which the measurements of Section~\ref{sec:t3} are also reported.
\end{remark}


% ===== Part I, section 3: Target 2 =====
\section{Target 2: benign overfitting along a doubly exponential sequence}\label{sec:proofT2}

Here $\ell\ge1$ is the level, $R=2^\ell$, $d=2^R$ and $n=d^R$, so that $d_\ell=2^{2^\ell}$ and
$n_\ell=2^{2^{2\ell}}$. The labels are $Y_i=f_d(X_i)+\varepsilon_i$ with $\varepsilon\sim N(0,\sigma^2I_n)$
independent of the design, where $f_d\in\bigoplus_{s\le q_0}\mathcal H_{d,s}$ for a degree bound $q_0$ that
does not depend on $\ell$, and $\|f_d\|_{L^2(\Sph)}\le B$. Let $\hat f$ be the interpolant of the data of
minimal norm in the reproducing kernel Hilbert space $\mathcal H_K$ of $K$, and
\[
\mathcal R_\ell=\|\hat f-f_d\|_{L^2(\Sph)}^2=\E_{X_0}\big[(\hat f(X_0)-f_d(X_0))^2\,\big|\,X,\varepsilon\big].
\]

\begin{theorem}[Target 2]\label{thm:T2}
For every fixed $h>0$, $q_0$, $B$ and $\sigma$, and every $\epsilon>0$, $\sum_\ell\Pr\{\mathcal R_\ell>\epsilon\}<\infty$.
In particular $\mathcal R_\ell\to0$ in probability, and almost surely along the sequence under any coupling of
the levels.
\end{theorem}

The critical degree is $k=R+\ell-1$; degrees $s<k$ form the low block, of dimension
$D=D_{<k}=\sum_{s<k}N_s$, and degrees $s\ge k$ form the tail. We write $\delta_*=\frac1{\log2}-1=0.4426\ldots$,
and throughout ``for large $\ell$'' means for all $\ell$ beyond a threshold depending only on $\gamma$.

\begin{lemma}[scales at the critical degree]\label{lem:scales}
As $\ell\to\infty$,
\[
\frac{N_k}{n}=d^{\,\delta_*+o(1)},\qquad \frac{D}{n}=d^{-(1-\delta_*)+o(1)},
\]
while $\log n=R^2\log2$, $n^{1/R}=d$, and $N_s$ is non-decreasing in $s$.
\end{lemma}

\begin{proof}
Since $k^2=O(R^2)=o(d)$, Lemma~\ref{lem:mult} gives $\log N_k=k\log d-\log k!+o(1)$. With
$\log d=R\log2$, $\ell=\log_2R$ and Stirling's formula $\log k!=k\log k-k+O(\log k)$,
\[
\begin{aligned}
\log\frac{N_k}{n}&=(\ell-1)R\log2-(R+\ell-1)\log(R+\ell-1)+(R+\ell-1)+O(\log R)\\
&=R(1-\log2)-(\ell-1)\log R+O(\ell),
\end{aligned}
\]
because $(\ell-1)\log2=\log(R/2)$ and $\log(R+\ell-1)=\log R+O(\ell/R)$. Dividing by $\log d=R\log2$
gives $\log(N_k/n)/\log d\to(1-\log2)/\log2=\delta_*$. For the low block, Lemma~\ref{lem:mult} gives
$1/\rho_s\le\frac{s+1}{d+s-2}\le\frac{2k}{d}$ for $s<k$ and $d\ge4k$, hence $N_{k-j}\le N_k(2k/d)^j$ and
$D\le N_k\sum_{j\ge1}(2k/d)^j\le4kN_k/d$, so $D/n\le d^{\delta_*-1+o(1)}$; the matching lower bound
$D\ge N_{k-1}\ge N_k\,d^{-1+o(1)}$ follows from $\rho_{k-1}\le\frac{d+2k}{k}\cdot\frac{d+k}{d+2k-2}\le d^{1+o(1)}$.
The identities for $\log n$ and $n^{1/R}$ are immediate from $n=d^R$, $d=2^R$, and monotonicity of $N_s$
is $\rho_s\ge1$ in Lemma~\ref{lem:mult}.
\end{proof}

\begin{lemma}[the tail kernel]\label{lem:tail}
Let $\tau=\sum_{s\ge k}\tau_s$, $w_s=\tau_s/\tau$ for $s\ge k$, $K_T=\sum_{s\ge k}\lambda_s\Pi_s$ and
$r_T=K_T/\tau=\sum_{s\ge k}w_sr_s$ with $r_s=\Pi_s/N_s$. Then $r_T(x,x)=1$, $\E_Yr_T(x,Y)=0$, and for large $\ell$:
\begin{enumerate}
\item[\textup{(i)}] $\tau_{k+r}/\tau_k\le2e^{\gamma}\gamma^rk!/(k+r)!$ for all $r\ge0$; hence $1-w_k\le4e^\gamma\gamma/(k+1)$;
\item[\textup{(ii)}] for every $z\in\Sph$ and $1\le q\le R$, $\ \|r_T(X,z)\|_{L^{2q}}\le4e^{\gamma}(2q-1)^{k/2}N_k^{-1/2}$;
\item[\textup{(iii)}] $\E\,r_T(X,z)^2=\sum_{s\ge k}w_s^2/N_s\in\big[w_k^2/N_k,\,1/N_k\big]$, so $\E\,r_T(X,z)^2=(1+o(1))/N_k$;
\item[\textup{(iv)}] with $G_T=(\E_z[r_T(X_i,z)r_T(z,X_j)])_{ij}$ and $R_T=(r_T(X_i,X_j))_{ij}$,
$\ 0\preceq G_T\preceq N_k^{-1}R_T$ and $(G_T)_{jj}\le1/N_k$.
\end{enumerate}
\end{lemma}

\begin{proof}
$r_T(x,x)=\sum w_s=1$ and $\E_Yr_T(x,Y)=0$ follow from \eqref{eq:repro} since every tail degree is positive.
(i) By \eqref{eq:mercerbounds}, $\tau_{k+r}\le e^{-\gamma}\sum_{m\ge k+r}\gamma^m/m!\le\gamma^{k+r}/(k+r)!$,
while $\tau_k=N_k\lambda_k\ge N_ke^{-\gamma}\gamma^k/(d)_k\ge\frac12e^{-\gamma}\gamma^k/k!$ for large $\ell$, because
$N_kk!/(d)_k=e^{O(k^2/d)}\prod_{j<k}(1+2j/d)^{-1}\to1$ by Lemma~\ref{lem:mult}. The bound on $1-w_k$
follows from $\sum_{r\ge1}\gamma^rk!/(k+r)!\le\sum_{r\ge1}(\gamma/(k+1))^r\le2\gamma/(k+1)$.
(ii) By the triangle inequality, Lemma~\ref{lem:hyper}, monotonicity of $N_s$ and $w_{k+r}\le\tau_{k+r}/\tau_k$,
\[
\|r_T(X,z)\|_{2q}\le\sum_{r\ge0}w_{k+r}\frac{(2q-1)^{(k+r)/2}}{\sqrt{N_{k+r}}}
\le\frac{(2q-1)^{k/2}}{\sqrt{N_k}}\,2e^{\gamma}\sum_{r\ge0}\Big(\frac{\gamma\sqrt{2q-1}}{k+1}\Big)^r
\le\frac{4e^{\gamma}(2q-1)^{k/2}}{\sqrt{N_k}},
\]
since $\gamma\sqrt{2R}/(k+1)\le1/2$ for large $\ell$. (iii) By \eqref{eq:repro}, $\E[r_s(X,z)r_t(X,z)]=\delta_{st}/N_s$;
the bounds use $\sum_sw_s^2\le(\sum_sw_s)^2=1$ and $N_s\ge N_k$. (iv) By \eqref{eq:repro},
$\E_z[r_s(X_i,z)r_t(z,X_j)]=\delta_{st}r_s(X_i,X_j)/N_s$, so $G_T=\sum_{s\ge k}(w_s/N_s)\,w_sr_s(X,X)$; each
$w_sr_s(X,X)$ is positive semidefinite and $w_s/N_s\le1/N_k$, which gives both claims.
\end{proof}

\begin{lemma}[the row-square statistic]\label{lem:rowsquare}
Let $\widetilde X$ be an independent copy of the design and, for $1\le j\le n$,
$S_j=\sum_{i\ne j}r_T(X_i,\widetilde X_j)^2$. Then for large $\ell$,
\[
\E\max_{1\le j\le n}S_j\ \le\ (1+o(1))\,\frac{n}{N_k}.
\]
\end{lemma}

\begin{proof}
Fix $z\in\Sph$ and put $Y_i=r_T(X_i,z)^2-\E r_T(X_i,z)^2$, independent, centred, and identically distributed.
Rosenthal's inequality with the constant of Johnson, Schechtman and Zinn \cite{Rosenthal,JSZ} gives, for $q\ge2$,
$\|\sum_{i\ne j}Y_i\|_q\le C_0q\big(\sqrt n\,\|Y_1\|_2+n^{1/q}\|Y_1\|_q\big)$ with an absolute constant $C_0$.
Take $q=R$. By Lemma~\ref{lem:tail}(ii) with $q=2$ and $q=R$, $\|Y_1\|_2\le\|r_T\|_4^2\le C_\gamma3^k/N_k$ and
$\|Y_1\|_R\le2\|r_T\|_{2R}^2\le C_\gamma(2R)^k/N_k$, and $n^{1/R}=d$. The bounds do not depend on $z$, so
\[
\E\max_j\big|S_j-\E S_j\big|\le\Big(\sum_j\E|S_j-\E S_j|^R\Big)^{1/R}\le n^{1/R}\sup_z\big\|\textstyle\sum_{i\ne j}Y_i\big\|_R
\le C_\gamma\,dR\Big(\frac{\sqrt n\,3^k}{N_k}+\frac{d\,(2R)^k}{N_k}\Big).
\]
Relative to $\E S_j\le n/N_k$, the two terms are $dR3^k/\sqrt n$ and $Rd^2(2R)^k/n$. In base $d$, using
$k=R+\ell-1$, $\log d=R\log2$ and $2R=2^{\ell+1}$: $3^k=d^{(\log_23)(1+(\ell-1)/R)}$ and
$(2R)^k=d^{(\ell+1)(1+(\ell-1)/R)}$, so the two ratios are $R\,d^{\,1+\log_23+o(1)-R/2}$ and
$R\,d^{\,3+\ell+o(\ell)-R}$, both of which tend to zero because $R=2^\ell$. Since $\E S_j=(n-1)\E r_T(X,z)^2\le n/N_k$
by Lemma~\ref{lem:tail}(iii), the claim follows.
\end{proof}

\begin{proposition}[isotropy of the tail Gram matrix]\label{prop:isotropy}
Let $R_T=(r_T(X_i,X_j))_{ij}$ and $A_d=\E\|R_T-I\|_{\mathrm{op}}$. Then $A_d\to0$; more precisely, for every fixed
$\eta<\delta_*$ and large $\ell$, $A_d\le C_{\gamma,\eta}R^2d^{-\eta/2}$.
\end{proposition}

The proof follows the method of Kaushik, Romberg and Muthukumar \cite{KRM}: decoupling, then the
noncommutative Khintchine inequality. It is written out because the quantity that must be controlled at
the growing degree $k$ is the row-square statistic of Lemma~\ref{lem:rowsquare}, and the maximum-entry
relaxation of that statistic which appears in the statement of \cite[Theorem~1]{KRM} is too large here by a
factor of order $d^{\ell}$.

\begin{proof}
Since $r_T(x,x)=1$, $\Delta:=R_T-I=\sum_{i\ne j}r_T(X_i,X_j)e_ie_j^\top$ is a matrix-valued $U$-statistic of
order two. By the decoupling inequality for $U$-statistics with values in a Banach space
\cite{dlPMS,dlPG}, applied to the convex function $\|\cdot\|_{\mathrm{op}}$, there is an absolute constant $C_1$ with
\[
\E\|\Delta\|\le C_1\E\|\widetilde\Delta\|,\qquad
\widetilde\Delta=\sum_{i\ne j}r_T(X_i,\widetilde X_j)e_ie_j^\top=\sum_{j=1}^nc_je_j^\top,
\qquad c_j=\big(r_T(X_i,\widetilde X_j)\mathbf 1_{\{i\ne j\}}\big)_{i\le n},
\]
where $\widetilde X$ is an independent copy of $X$. Condition on $X$. The columns $c_1,\dots,c_n$ are then
independent (each is a function of its own $\widetilde X_j$) and centred, by Lemma~\ref{lem:tail}. Moreover
$\E[c_jc_j^\top\mid X]=P_jG_TP_j=:M_j$ with $P_j=I-e_je_j^\top$ and $G_T$ as in Lemma~\ref{lem:tail}(iv);
hence $\|M_j\|\le\|G_T\|$ and
\begin{equation}
\Big\|\sum_jM_j\Big\|=\big\|(n-2)G_T+\mathrm{diag}(G_T)\big\|\le n\|G_T\|+\frac1{N_k},
\qquad \Big\|\sum_jM_j^2\Big\|\le n\|G_T\|^2 .
\label{eq:Mj}
\end{equation}
All expectations below are conditional on $X$ unless stated otherwise. Write $\|c_j\|^2=S_j$ as in
Lemma~\ref{lem:rowsquare}, $m_2=\E\max_jS_j$, and $U=\E\|\sum_jc_jc_j^\top\|$; thus $m_2$ and $U$ are
$X$-measurable random variables, and it is their $X$-expectations that enter at the end of the proof, where
Lemma~\ref{lem:rowsquare} bounds $\E m_2$.

\emph{First application of Khintchine.} By symmetrisation \cite[Lemma~6.3]{LedouxTalagrand} and the matrix
Khintchine inequality for Rademacher series \cite[Theorem~4.1.1]{Tropp15}, applied to the $n\times n$ matrices
$c_je_j^\top$, for which
\[
\sum_j(c_je_j^\top)(c_je_j^\top)^\top=\sum_jc_jc_j^\top,\qquad
\sum_j(c_je_j^\top)^\top(c_je_j^\top)=\mathrm{diag}(S_1,\dots,S_n),
\]
we get
\begin{equation}
\E\|\widetilde\Delta\|\le2\sqrt{2\log(2n)}\ \E\max\Big(\big\|\textstyle\sum_jc_jc_j^\top\big\|^{1/2},\ \max_j\sqrt{S_j}\Big)
\le2\sqrt{2\log(2n)}\,\big(\sqrt U+\sqrt{m_2}\big),
\label{eq:K1}
\end{equation}
by Jensen's inequality.

\emph{Second application.} The matrices $Z_j=c_jc_j^\top-M_j$ are independent, centred and symmetric, so by
the same two inequalities, $\E\|\sum_jZ_j\|\le2\sqrt{2\log(2n)}\,\E\|\sum_jZ_j^2\|^{1/2}$. Since
$Z_j^2\preceq2(c_jc_j^\top)^2+2M_j^2=2S_jc_jc_j^\top+2M_j^2$, \eqref{eq:Mj} gives
$\|\sum_jZ_j^2\|\le2\max_jS_j\,\|\sum_jc_jc_j^\top\|+2n\|G_T\|^2$, and therefore, by Cauchy--Schwarz,
\[
U\le\Big\|\sum_jM_j\Big\|+\E\Big\|\sum_jZ_j\Big\|
\le n\|G_T\|+\frac1{N_k}+4\sqrt{\log(2n)}\Big(\sqrt{m_2U}+\sqrt n\,\|G_T\|\Big).
\]
Using $4\sqrt{\log(2n)}\sqrt{m_2U}\le U/2+8\log(2n)\,m_2$ and rearranging,
\begin{equation}
U\le2n\|G_T\|\Big(1+4\sqrt{\tfrac{\log(2n)}{n}}\Big)+\frac2{N_k}+16\log(2n)\,m_2 .
\label{eq:U}
\end{equation}

\emph{Closing the recurrence.} Take expectations over $X$. By Lemma~\ref{lem:tail}(iv),
$\E\|G_T\|\le\E\|R_T\|/N_k\le(1+A_d)/N_k$, and by Lemma~\ref{lem:rowsquare}, $\E m_2\le(1+o(1))n/N_k$.
Inserting \eqref{eq:U} into \eqref{eq:K1} (with $\E\sqrt U\le\sqrt{\E U}$) and then into the decoupling bound,
\[
A_d\le C\sqrt{\log n}\Big(\sqrt{\frac{n(1+A_d)}{N_k}}+\sqrt{\frac{\log n\ n}{N_k}}\Big)
\le C\sqrt{\frac{n\log n}{N_k}}\,\big(\sqrt{1+A_d}+\sqrt{\log n}\big)
\]
for large $\ell$, with $C$ absolute. Put $\epsilon_d=n\log n/N_k$; by Lemma~\ref{lem:scales},
$\epsilon_d\log n=d^{-\delta_*+o(1)}R^4\log^22\to0$. If $A_d\ge1$ the inequality gives
\[
\sqrt{A_d}\le C\sqrt{\epsilon_d}\Big(\sqrt2+\frac{\sqrt{\log n}}{\sqrt{A_d}}\Big)\le C\sqrt{\epsilon_d}\big(\sqrt2+\sqrt{\log n}\big),
\]
which is smaller than one for large $\ell$, a contradiction; so $A_d<1$ for large $\ell$, and then
\[
A_d\le C\sqrt{\epsilon_d}\big(\sqrt2+\sqrt{\log n}\big)\le C'\sqrt{\epsilon_d\log n}=C'R^2d^{-\delta_*/2+o(1)},
\]
which is at most $C_{\gamma,\eta}R^2d^{-\eta/2}$ for every fixed $\eta<\delta_*$.
\end{proof}

\begin{lemma}[the interpolant through the low block]\label{lem:lowblock}
Let $\Psi\in\R^{n\times D}$ be the matrix of an $L^2$-orthonormal basis of $\bigoplus_{s<k}\mathcal H_{d,s}$
evaluated at the design, $\Lambda$ the diagonal matrix of the corresponding eigenvalues $\lambda_s$, and
$T=(K_T(X_i,X_j))_{ij}=\tau R_T$. Let $E_1=\{\|\Psi^\top\Psi/n-I_D\|\le1/2\}$ and $E_2=\{\|R_T-I\|\le1/2\}$. Then
$\Pr(E_1^c)\le2De^{-n/(10D)}$ and $\Pr(E_2^c)\le2A_d$. On $E_1\cap E_2$ the minimal-norm interpolant is
$\hat f=\Psi(\cdot)^\top\hat a+k_T(\cdot)^\top T^{-1}(y-\Psi\hat a)$ with $k_T(x)=(K_T(x,X_i))_i$ and
\begin{equation}
\hat a=G_0^{-1}\Psi^\top Wy,\qquad W=\tau T^{-1},\qquad G_0=\tau\Lambda^{-1}+\Psi^\top W\Psi,\qquad
\tfrac23I\preceq W\preceq2I,\qquad \|G_0^{-1}\|\le\frac3n .
\label{eq:normal}
\end{equation}
\end{lemma}

\begin{proof}
The rows $\Psi(X_i)$ are independent, and by \eqref{eq:repro}
\[
\|\Psi(X_i)\|^2=\sum_{s<k}\Pi_s(X_i,X_i)=D,\qquad \E\,\Psi(X_i)\Psi(X_i)^\top=I_D .
\]
The matrix Chernoff inequality of Tropp \cite[Theorem~1.1]{Tropp12}, at deviation $1/2$, gives
\[
\Pr(E_1^c)\le D\big(e^{-1/2}\sqrt2\big)^{n/D}+D\big(e^{1/2}(3/2)^{-3/2}\big)^{n/D}\le2De^{-n/(10D)} .
\]
Markov's inequality gives $\Pr(E_2^c)\le2A_d$. The RKHS $\mathcal H_K$ is the orthogonal sum of
$\mathcal H_{\rm low}=\bigoplus_{s<k}\mathcal H_{d,s}$ with norm $\|\Psi(\cdot)^\top a\|^2_{\mathcal H_K}=a^\top\Lambda^{-1}a$
and of the tail space $\mathcal H_T$ with kernel $K_T$. For fixed low coefficients $a$, the tail function of
minimal norm interpolating $y-\Psi a$ at the design is $k_T(\cdot)^\top T^{-1}(y-\Psi a)$, of squared norm
$(y-\Psi a)^\top T^{-1}(y-\Psi a)$ ($T$ is invertible on $E_2$). The minimal-norm interpolant therefore
minimises $J(a)=a^\top\Lambda^{-1}a+(y-\Psi a)^\top T^{-1}(y-\Psi a)$, a strictly convex quadratic whose
stationarity condition is $(\Lambda^{-1}+\Psi^\top T^{-1}\Psi)a=\Psi^\top T^{-1}y$; multiplying by $\tau$ gives
\eqref{eq:normal}. On $E_2$, $\tfrac12I\preceq T/\tau\preceq\tfrac32I$, hence $\tfrac23I\preceq W\preceq2I$; on
$E_1\cap E_2$, $\Psi^\top W\Psi\succeq\tfrac23\cdot\tfrac n2I=\tfrac n3I$, so $G_0\succeq\tfrac n3I$.
\end{proof}

\begin{proof}[Proof of Theorem~\ref{thm:T2}]
Take $\ell$ large enough that $q_0<k$ and write $f_d=\Psi(\cdot)^\top\theta$ with $\|\theta\|=\|f_d\|_{L^2}\le B$
and $\theta$ supported on degrees $\le q_0$. Work on $E_1\cap E_2$. Since low and tail functions are
orthogonal in $L^2(\Sph)$ and $\|\Psi(\cdot)^\top a\|_{L^2}=\|a\|$,
\[
\mathcal R_\ell=\|\hat a-\theta\|^2+\|f_T\|_{L^2}^2,\qquad f_T=k_T(\cdot)^\top T^{-1}r,\qquad r=y-\Psi\hat a .
\]
\emph{The low coefficients.} With $y=\Psi\theta+\varepsilon$ and \eqref{eq:normal},
$\hat a-\theta=-\tau G_0^{-1}\Lambda^{-1}\theta+G_0^{-1}\Psi^\top W\varepsilon$. For the first term, $\tau\le1$ and, by
\eqref{eq:mercerbounds}, $\lambda_{d,s}\ge e^{-\gamma}\gamma^{s}/(d)_s\ge c_{\gamma,q_0}d^{-q_0}$ for $s\le q_0$, so
$\|\tau G_0^{-1}\Lambda^{-1}\theta\|\le\frac3n\cdot\frac{B}{c_{\gamma,q_0}}\,d^{q_0}$. For the second,
$\E_\varepsilon\|G_0^{-1}\Psi^\top W\varepsilon\|^2=\sigma^2\Tr\big(G_0^{-1}\Psi^\top W^2\Psi G_0^{-1}\big)
\le\sigma^2\|G_0^{-1}\|^2\|W\|^2\Tr(\Psi^\top\Psi)\le\sigma^2\cdot\frac9{n^2}\cdot4\cdot nD=\frac{36\sigma^2D}n$.
Hence
\begin{equation}
\E_\varepsilon\|\hat a-\theta\|^2\le C_{\gamma,q_0}\frac{B^2d^{2q_0}}{n^2}+\frac{72\sigma^2D}{n}.
\label{eq:lowrisk}
\end{equation}
\emph{The tail function.} By \eqref{eq:repro}, $\|f_T\|_{L^2}^2=r^\top T^{-1}H_TT^{-1}r$ with
\[
H_T=\E_{X_0}\big[k_T(X_0)k_T(X_0)^\top\big]=\sum_{s\ge k}\lambda_s^2\Pi_s(X,X)\preceq\lambda_kT,
\]
because $\lambda_s\le\lambda_k$ for $s\ge k$ (Lemma~\ref{lem:mercer}). On $E_2$, $T\succeq\tfrac\tau2I$, so
\[
\|f_T\|_{L^2}^2\le\lambda_k\,r^\top T^{-1}r\le\frac{2\lambda_k}{\tau}\|r\|^2\le\frac2{N_k}\|r\|^2,
\]
using $\lambda_k/\tau\le\lambda_kN_k/(\tau_kN_k)\le1/N_k$. Since $r=\Psi(\theta-\hat a)+\varepsilon$ and $\|\Psi\|^2\le\tfrac32n$
on $E_1$, $\|r\|^2\le3n\|\theta-\hat a\|^2+2\|\varepsilon\|^2$, and with \eqref{eq:lowrisk},
\[
\E_\varepsilon\|f_T\|_{L^2}^2\le\frac{6n}{N_k}\Big(C_{\gamma,q_0}\frac{B^2d^{2q_0}}{n^2}+\frac{72\sigma^2D}n\Big)+\frac{4n\sigma^2}{N_k}.
\]
\emph{The conditional risk bound.} Combining the two displays, on $E_1\cap E_2$,
\begin{equation}
\E_\varepsilon[\mathcal R_\ell\mid X]\le C_{\gamma,q_0}\Big[\sigma^2\Big(\frac Dn+\frac n{N_k}\Big)+B^2\Big(\frac{d^{2q_0}}{n^2}+\frac{d^{2q_0}}{nN_k}\Big)\Big]
=C_{\gamma,q_0,B,\sigma}\,d^{-\delta_*+o(1)},
\label{eq:riskbound}
\end{equation}
where the last step is Lemma~\ref{lem:scales}: $n/N_k=d^{-\delta_*+o(1)}$ is the slowest term, since
$D/n=d^{-(1-\delta_*)+o(1)}$ with $1-\delta_*>\delta_*$, and $d^{2q_0}/n=d^{2q_0-R}$.

\emph{Conclusion.} By Markov's inequality in the noise, $\Pr\{\mathcal R_\ell>\epsilon\mid X\}\le\epsilon^{-1}\E_\varepsilon[\mathcal R_\ell\mid X]$
on $E_1\cap E_2$, so by Lemma~\ref{lem:lowblock} and Proposition~\ref{prop:isotropy},
\[
\Pr\{\mathcal R_\ell>\epsilon\}\le2De^{-n/(10D)}+2C_{\gamma,\eta}R^2d^{-\eta/2}+\frac{C_{\gamma,q_0,B,\sigma}}{\epsilon}\,d^{-\delta_*+o(1)} .
\]
Every term is summable in $\ell$ because $d=d_\ell=2^{2^\ell}$ and $n/D\ge d^{1-\delta_*+o(1)}$. This proves
the summability claim; convergence in probability follows, and the first Borel--Cantelli lemma gives
$\mathcal R_\ell\to0$ almost surely whatever the joint law of the levels, since only the marginal tail
probabilities enter.
\end{proof}

\begin{remark}\label{rem:T2numerics}
The mechanism is a dimension gap: the tail Gram matrix is isotropic because $n\ll N_k$, and the low block is
well conditioned because $D\ll n$; both gaps are polynomial in $d$ with the exponents $\delta_*$ and
$1-\delta_*$ of Lemma~\ref{lem:scales}. Section~\ref{sec:gap} measures the isotropy of $R_T$ and the risk of
the interpolant as functions of the finite-size analogue $\delta=\log(N_{d,k}/n)/\log d$ of these exponents,
on both sides of $\delta=0$, which is where Theorem~\ref{thm:T3} lives.
\end{remark}


% ===== Part I, section 4: Target 1 =====
\section{Target 1: the correlated Gaussian equivalence}\label{sec:proofT1}

\subsection{Model, surrogate and statement}

Let $F=[f_1,\dots,f_m]\in\R^{d\times m}$ have independent $N(0,1/d)$ entries. For each original sample
$i\le n$ let $a_i,z_{i1},z_{i2}$ be independent $N(0,I_d)$ vectors, independent of $F$, and let
$t_i=\xi^\top a_i$ with $\|\xi\|=1$ fixed. The two augmented feature rows of sample $i$ are
\[
x^A_{ij}=\sigma\big(F^\top(a_i+\Delta z_{ij})\big)\in\R^m,\qquad j=1,2,
\]
with $\Delta>0$ fixed and $\sigma:\R\to\R$ odd, bounded, three times differentiable with bounded
derivatives (for instance $\tanh$). With $N=2n$, $y\in\R^N$ the vector of the labels $t_i$ each repeated
twice and $X_A\in\R^{N\times m}$ the matrix of the rows $x^A_{ij}$, the objective is
\begin{equation}
\begin{aligned}
L_A(w)&=\frac1{2N}\|y-X_Aw\|^2+\frac\lambda2\|w\|^2\\
&=\frac1{4n}\sum_{i=1}^n\sum_{j=1}^2\big(t_i-(x^A_{ij})^\top w\big)^2+\frac\lambda2\|w\|^2 ,
\end{aligned}
\label{eq:objective}
\end{equation}
$\lambda>0$ fixed, with minimiser $w_A^\star$ and minimum $L_A^\star$; $n/d\to\alpha>0$ and $m/d\to\beta>0$.
A clean test point is $x^{0,A}=\sigma(F^\top a^0)$ with $t^0=\xi^\top a^0$, $a^0\sim N(0,I_d)$ independent of
everything, and the population clean risk of a vector $w$ is
$R_A(w\mid F)=\E[(t^0-(x^{0,A})^\top w)^2\mid F,w]$.

The surrogate replaces $\sigma$ by its first Hermite coefficients and the rest by independent Gaussian
residuals. Let $s=\sqrt{1+\Delta^2}$, $\rho=s^{-2}$, let $h_k$ be the orthonormal probabilists' Hermite
polynomials, $G\sim N(0,1)$, and for $q>0$
\begin{equation}
\begin{gathered}
a_k(q)=\E\big[\sigma(s\sqrt q\,G)h_k(G)\big],\qquad \kappa(q)=\E\big[\sigma'(s\sqrt q\,G)\big],\\
\eta(q)=\sum_{k\ge3\ \mathrm{odd}}a_k(q)^2,\qquad \gamma(q)=\sum_{k\ge3\ \mathrm{odd}}a_k(q)^2\rho^k .
\end{gathered}
\label{eq:hermitecoef}
\end{equation}
Since $\sigma$ is odd, $a_k(q)=0$ for even $k$, and Gaussian integration by parts gives
$a_1(q)=s\sqrt q\,\kappa(q)$. Put $\mu_e=\kappa(1)$, $\mu_b=\Delta\kappa(1)$, $\mu_2=\sqrt{\gamma(1)}$,
$\mu_3=\sqrt{\eta(1)-\gamma(1)}$ (both radicands are non-negative, see Lemma~\ref{lem:residualcov}), and
define the surrogate rows and clean row
\begin{equation}
x^G_{ij}=\mu_eF^\top a_i+\mu_bF^\top z_{ij}+\mu_2b_i+\mu_3p_{ij},\qquad
x^{0,G}=\nu_1F^\top a^0+\nu_2g^0,
\label{eq:surrogate}
\end{equation}
with $b_i,p_{i1},p_{i2},g^0\sim N(0,I_m)$ independent of everything else, $\nu_1=\E\sigma'(G)$ and
$\nu_2^2=\E\sigma(G)^2-\nu_1^2$. The same $a_i$ and the same $b_i$ enter both views of sample $i$: the
surrogate keeps the shared latent structure of the augmentation. The four training coefficients are those
of Dhifallah and Lu \cite{DhifallahLu}. $L_G$, $w_G^\star$, $L_G^\star$ and $R_G(w\mid F)$ are defined from
\eqref{eq:objective} with $X_G$ in place of $X_A$ and $x^{0,G}$ in place of $x^{0,A}$.

\begin{theorem}[Target 1]\label{thm:T1}
There are constants $C,c>0$ depending only on $\sigma,\Delta,\lambda,\alpha,\beta$ such that for every
$\varphi\in C^2_b(\R)$ and all large $d$,
\begin{align}
\big|\E\varphi(L_A^\star)-\E\varphi(L_G^\star)\big|
&\le C\Big(\|\varphi'\|_\infty\sqrt{\tfrac{\log d}{d}}+\|\varphi''\|_\infty\,\tfrac1d\Big)+C\|\varphi\|_\infty e^{-cd},
\label{eq:T1a}\\
\big|\E\varphi\big(R_A(w_A^\star\mid F)\big)-\E\varphi\big(R_G(w_G^\star\mid F)\big)\big|
&\le C\Big(\|\varphi'\|_\infty\sqrt{\tfrac{\log d}{d}}+\|\varphi''\|_\infty\,\tfrac1d\Big)+C\|\varphi\|_\infty e^{-cd}.
\label{eq:T1b}
\end{align}
In particular the optimised objective and the clean prediction risk of the augmented ridge problem have
the same limits in law as those of the surrogate.
\end{theorem}

The proof replaces the $n$ blocks $(t_i,x_{i1},x_{i2})$ one at a time, in the manner of Lindeberg. Two exact
facts make each step cheap: the change of the optimum under the insertion of one block is an explicit
$2\times2$ quadratic form (Lemma~\ref{lem:ranktwo}), and conditional on $F$ there is a Gaussian block whose
second moments match those of the actual block exactly (Proposition~\ref{prop:gaussianization}), so that
the only mismatch between the actual block and the Gaussian one is the fluctuation of a $2\times2$ leverage
matrix, and the only mismatch between that Gaussian block and the surrogate is the replacement of the
random column norms $q_r=\|f_r\|^2$ by one.

\subsection{Exact identities}

\begin{lemma}[insertion of one block]\label{lem:ranktwo}
Let $X_-\in\R^{(N-2)\times m}$ and $y_-$ be the data with the two rows of sample $i$ removed, let
$L_-(w)=\frac1{2N}\|y_--X_-w\|^2+\frac\lambda2\|w\|^2$ with minimiser $w_-$ and minimum $L_-^\star$, and put
\[
H_-=\lambda I+\tfrac1NX_-^\top X_-,\qquad B=H_-^{-1},\qquad
X_i=\begin{pmatrix}x_{i1}^\top\\x_{i2}^\top\end{pmatrix},\qquad e_i=t_i\mathbf 1_2-X_iw_-,\qquad
Q_i=\tfrac1NX_iBX_i^\top .
\]
Then the minimum $L^\star$ of the full objective satisfies
\begin{equation}
L^\star-L_-^\star=\frac1{2N}\,e_i^\top(I_2+Q_i)^{-1}e_i,\qquad
w^\star-w_-=\frac1N\,BX_i^\top(I_2+Q_i)^{-1}e_i,
\label{eq:ranktwo}
\end{equation}
with $Q_i\succeq0$, $\|(I_2+Q_i)^{-1}\|\le1$ and $\|B\|\le\lambda^{-1}$.
\end{lemma}

\begin{proof}
$L_-$ is a quadratic with Hessian $H_-$, so $L_-(w)=L_-^\star+\frac12(w-w_-)^\top H_-(w-w_-)$, and the full
objective is $L(w)=L_-(w)+\frac1{2N}\|e_i-X_i(w-w_-)\|^2$. Minimising over $u=w-w_-$ gives
$u^\star=(H_-+\frac1NX_i^\top X_i)^{-1}\frac1NX_i^\top e_i$ and
$L^\star-L_-^\star=\frac1{2N}e_i^\top\big(I-\frac1NX_i(H_-+\frac1NX_i^\top X_i)^{-1}X_i^\top\big)e_i$. The
push-through identity $\frac1NX_i(H_-+\frac1NX_i^\top X_i)^{-1}X_i^\top=Q_i(I+Q_i)^{-1}$ gives the first
formula in \eqref{eq:ranktwo}, and $(H_-+\frac1NX_i^\top X_i)^{-1}X_i^\top=BX_i^\top(I+Q_i)^{-1}$ the second.
The remaining claims follow from $H_-\succeq\lambda I$.
\end{proof}

\begin{lemma}[residual covariances]\label{lem:residualcov}
For $q>0$ let $V(q)=\E\sigma(s\sqrt q\,G)^2=\sum_{k\ \mathrm{odd}}a_k(q)^2$ and, for a standard bivariate
Gaussian $(H_1,H_2)$ with correlation $\rho$, $C(q)=\E[\sigma(s\sqrt q H_1)\sigma(s\sqrt qH_2)]=\sum_{k\ \mathrm{odd}}a_k(q)^2\rho^k$.
Then $\eta(q)=V(q)-s^2q\kappa(q)^2$, $\gamma(q)=C(q)-q\kappa(q)^2$, $0\le\gamma(q)\le\eta(q)\le\|\sigma\|_\infty^2$,
and $V,\kappa,C,\eta,\gamma$ are Lipschitz on $[\tfrac12,\tfrac32]$ with constants depending only on
$\sigma$ and $\Delta$.
\end{lemma}

\begin{proof}
The expansions of $V$ and $C$ are Parseval's identity and Mehler's formula \cite[Chapter~1]{Janson} for
the Hermite expansion of $\sigma(s\sqrt q\,\cdot)$; the $k=1$ terms are $a_1(q)^2=s^2q\kappa(q)^2$ and
$a_1(q)^2\rho=q\kappa(q)^2$, which gives the two identities and, since $0<\rho<1$, the inequalities. For the
Lipschitz bounds, differentiation under the Gaussian expectation is justified by the bounded derivatives
of $\sigma$: with $b(q)=s\sqrt q$, $V'(q)=\frac{s}{\sqrt q}\E[G\sigma(b(q)G)\sigma'(b(q)G)]$,
$\kappa'(q)=\frac{s}{2\sqrt q}\E[G\sigma''(b(q)G)]$ and
$C'(q)=\frac{s}{2\sqrt q}\E[H_1\sigma'(b(q)H_1)\sigma(b(q)H_2)+H_2\sigma(b(q)H_1)\sigma'(b(q)H_2)]$, all bounded
on $[\frac12,\frac32]$.
\end{proof}

\begin{proposition}[exact Gaussianisation conditional on $F$]\label{prop:gaussianization}
Let $Q_F=\mathrm{diag}(q_1,\dots,q_m)$, $q_r=\|f_r\|^2$, $G_F=F^\top F$, $\mathscr R=Q_F^{-1/2}G_FQ_F^{-1/2}$ (a
correlation matrix), $D_F=\mathrm{diag}(\kappa(q_r))$, $A_k=\mathrm{diag}(a_k(q_r))$, and
\[
\Sigma_{\rm com}=\sum_{k\ge3\ \mathrm{odd}}\rho^kA_k\mathscr R^{\circ k}A_k,\qquad
\Sigma_{\rm priv}=\sum_{k\ge3\ \mathrm{odd}}(1-\rho^k)A_k\mathscr R^{\circ k}A_k,
\]
where $\circ k$ is the $k$-fold Hadamard power. Both are positive semidefinite. Conditional on $F$ let
$b_i^\#\sim N(0,\Sigma_{\rm com})$ and $p_{i1}^\#,p_{i2}^\#\sim N(0,\Sigma_{\rm priv})$ be independent of each other
and of $(a_i,z_{i1},z_{i2})$, and put
\begin{equation}
x^\#_{ij}=D_FF^\top a_i+\Delta D_FF^\top z_{ij}+b_i^\#+p^\#_{ij}.
\label{eq:sharp}
\end{equation}
Then, conditional on $F$, for $M\in\{A,\#\}$ the vectors $(t_i,x^M_{i1},x^M_{i2})$ have the same means, the same
covariances and the same cross-covariances:
\begin{equation}
\begin{gathered}
\E[t_ix^M_{ij}\mid F]=D_FF^\top\xi,\qquad
\mathrm{Cov}(x^M_{ij}\mid F)=s^2D_FG_FD_F+\Sigma_{\rm com}+\Sigma_{\rm priv},\\
\mathrm{Cov}(x^M_{i1},x^M_{i2}\mid F)=D_FG_FD_F+\Sigma_{\rm com}.
\end{gathered}
\label{eq:matching}
\end{equation}
Likewise, with $\kappa_0(q)=\E\sigma'(\sqrt q\,G)$, $a_{0,k}(q)=\E[\sigma(\sqrt q\,G)h_k(G)]$,
$D_{0,F}=\mathrm{diag}(\kappa_0(q_r))$, $A_{0,k}=\mathrm{diag}(a_{0,k}(q_r))$ and
$\Sigma_{0}=\sum_{k\ge3\ \mathrm{odd}}A_{0,k}\mathscr R^{\circ k}A_{0,k}\succeq0$, the clean row
$x^{0,\#}=D_{0,F}F^\top a^0+r^0$, $r^0\sim N(0,\Sigma_0)$ independent of $a^0$, satisfies
\[
\E[t^0x^{0,\#}\mid F]=\E[t^0x^{0,A}\mid F]=D_{0,F}F^\top\xi,\qquad
\mathrm{Cov}(x^{0,\#}\mid F)=\mathrm{Cov}(x^{0,A}\mid F)=D_{0,F}G_FD_{0,F}+\Sigma_0\,;
\]
consequently $R_A(w\mid F)=R_\#(w\mid F)$ for every $w$ that is independent of $a^0$.
\end{proposition}

\begin{proof}
Fix $F$. For view $j$ and coordinate $r$ the standardised pre-activation
$H_{jr}=f_r^\top(a_i+\Delta z_{ij})/(s\sqrt{q_r})$ is standard Gaussian, with
$\mathrm{Cov}(H_{jr},H_{js})=\mathscr R_{rs}$ within a view and $\mathrm{Cov}(H_{1r},H_{2s})=\rho\mathscr R_{rs}$
across views. Mehler's formula gives, for the actual rows,
\[
\mathrm{Cov}(x^A_{ij}\mid F)=\sum_{k\ \mathrm{odd}}A_k\mathscr R^{\circ k}A_k,\qquad
\mathrm{Cov}(x^A_{i1},x^A_{i2}\mid F)=\sum_{k\ \mathrm{odd}}\rho^kA_k\mathscr R^{\circ k}A_k ,
\]
and the $k=1$ terms are $A_1\mathscr RA_1=s^2D_FG_FD_F$ and $\rho A_1\mathscr RA_1=D_FG_FD_F$ by
$a_1(q)=s\sqrt q\,\kappa(q)$ and $\mathscr R_{rs}\sqrt{q_rq_s}=(G_F)_{rs}$. Since
$(t_i,f_r^\top(a_i+\Delta z_{ij}))$ is centred jointly Gaussian with covariance $\xi^\top f_r$, Stein's identity gives
$\E[t_i\sigma(f_r^\top(a_i+\Delta z_{ij}))\mid F]=\xi^\top f_r\,\E[\sigma'(f_r^\top(a_i+\Delta z_{ij}))\mid F]=\kappa(q_r)\xi^\top f_r$,
which is $D_FF^\top\xi$. Positive semidefiniteness of $\Sigma_{\rm com}$ and $\Sigma_{\rm priv}$ follows from the
Schur product theorem ($\mathscr R^{\circ k}\succeq0$) and congruence by the diagonal matrices $A_k$. The
moments of $x^\#$ are read off \eqref{eq:sharp}: the Gaussian block is centred, its teacher cross-covariance
is $D_FF^\top\E[t_ia_i]=D_FF^\top\xi$, its within-view covariance is
$D_FF^\top F D_F+\Delta^2D_FF^\top FD_F+\Sigma_{\rm com}+\Sigma_{\rm priv}$ and its cross-view covariance drops the
private terms; $1+\Delta^2=s^2$ gives \eqref{eq:matching}. The clean statements are the same computation
with $\Delta=0$, and the risk identity holds because $R_M(w\mid F)=1-2w^\top\E[t^0x^{0,M}\mid F]+w^\top\mathrm{Cov}(x^{0,M}\mid F)w$
depends only on the matched moments.
\end{proof}

\subsection{Estimates}

Let $\Omega_F$ be the event that $\|F\|_{\mathrm{op}}\le3(1+\sqrt\beta)$ and $\max_r|q_r-1|\le\frac12$. Since $F$
has independent Gaussian entries of variance $1/d$ and $q_r\sim\chi^2_d/d$, $\Pr(\Omega_F^c)\le Ce^{-cd}$
\cite[Theorem~4.4.5 and Section~2.8]{Vershynin}. Constants $C$ below depend only on
$\sigma,\Delta,\lambda,\alpha,\beta$, and all statements are for large $d$.

\begin{lemma}[a priori bounds]\label{lem:apriori}
For every hybrid sample in which each block is either actual, of type $\#$, or of type $G$, and every
minimiser $w$ of the corresponding objective (full or with one block removed),
$\lambda\|w\|^2\le\frac1n\sum_it_i^2$; hence $\sup_d\E\|w\|^{2p}<\infty$ for every $p$. On $\Omega_F$, the
within-view and cross-view covariances of every block type have operator norm at most $C$, and for every
$F$-measurable $w$ independent of the block, $\E[|x_{ij}^\top w|^{p}\mid F,w]\le C_p\|w\|^p$ and
$\E[\|e_i\|^p\mid F,w]\le C_p(1+\|w\|^p)$ for every $p$.
\end{lemma}

\begin{proof}
$\frac\lambda2\|w\|^2\le L(w)\le L(0)=\frac1{2N}\|y\|^2=\frac1{2n}\sum_it_i^2$, and the moments of the empirical
mean of $t_i^2\sim\chi^2_1$ are bounded. On $\Omega_F$, $\|G_F\|\le9(1+\sqrt\beta)^2$, $\|D_F\|\le\|\sigma'\|_\infty$,
$\|A_k\|\le\|\sigma\|_\infty$, and $\|\mathscr R\|\le\|Q_F^{-1}\|\,\|G_F\|\le2\|G_F\|$. The Hadamard sums are
bounded because $\sum_k\|A_k\|^2\le\|\sigma\|_\infty^2$ and, for positive semidefinite $A,B$,
$\|A\circ B\|\le\max_rA_{rr}\,\|B\|$, so that $\|\mathscr R^{\circ k}\|\le\|\mathscr R\|$ for every $k$ (a correlation
matrix has unit diagonal). For the actual block, $x_{ij}^\top w$ is a function of the Gaussian vector
$(a_i,z_{ij})$ with Lipschitz constant $\|\sigma'\|_\infty\|Fw\|\le C\|w\|$ and mean zero (oddness), so its
moments are bounded by Gaussian concentration \cite[Theorem~5.6]{BLM}; for the Gaussian block types the
same bound is the Gaussian moment bound with the covariance estimate. $e_i=t_i\mathbf 1_2-X_iw$ is handled
coordinatewise.
\end{proof}

\begin{lemma}[leverage fluctuations]\label{lem:leverage2}
Conditional on the leave-block data (which fix $F$, $B$ and $w_-$), on $\Omega_F$, for the actual block and
for the $\#$ block,
\[
\E\big[\|Q_i-\bar Q_i\|_F^2\,\big|\,\mathcal F_{-i}\big]\le\frac Cd,\qquad \bar Q_i=\E[Q_i\mid\mathcal F_{-i}],
\]
and $\bar Q_i^A=\bar Q_i^\#$. The same holds for $R_i=\frac1NX_iBS_0BX_i^\top$ with any fixed symmetric
$S_0$ with $\|S_0\|\le C$, and $\|(I+Q)^{-1}-(I+\bar Q)^{-1}\|\le\|Q-\bar Q\|$ for positive semidefinite $Q,\bar Q$.
\end{lemma}

\begin{proof}
$(Q_i)_{ab}=\frac1Nx_{ia}^\top Bx_{ib}$. For the actual block it is a function of the Gaussian vector
$(a_i,z_{i1},z_{i2})\in\R^{3d}$ whose gradient with respect to $a_i$ is
$\frac1N\big(F\,\mathrm{diag}(\sigma'_a)Bx_{ib}+F\,\mathrm{diag}(\sigma'_b)Bx_{ia}\big)$ (and similarly for $z_{i1},z_{i2}$),
of norm at most $\frac CN\|F\|\lambda^{-1}\sqrt m\,\|\sigma\|_\infty\le C/\sqrt d$; the Gaussian Poincar\'e
inequality \cite[Theorem~3.20]{BLM} gives the variance bound. For the $\#$ block, $(Q_i)_{ab}$ is a quadratic
form in a Gaussian vector; its gradient has norm at most $\frac CN\|F\|\lambda^{-1}(\|x_{ia}^\#\|+\|x_{ib}^\#\|)$,
whose second moment is $O(m/N^2)=O(1/d)$ by Lemma~\ref{lem:apriori}, and the Poincar\'e inequality
applies again. $\bar Q_i^A=\bar Q_i^\#$ because $\E[x_{ia}^\top Bx_{ib}\mid\mathcal F_{-i}]=\Tr(B\,\mathrm{Cov}(x_{ia},x_{ib}\mid F))$
and the covariances match by Proposition~\ref{prop:gaussianization}. The argument for $R_i$ is identical
with $BS_0B$ in place of $B$. The Lipschitz bound for the inverse is the resolvent identity
$(I+Q)^{-1}-(I+\bar Q)^{-1}=(I+Q)^{-1}(\bar Q-Q)(I+\bar Q)^{-1}$ with both resolvents of norm at most one.
\end{proof}

\begin{lemma}[from the exact Gaussian block to the surrogate]\label{lem:bridge}
Let $C^\#_{ab}$, $C^G_{ab}$ ($a,b\in\{1,2\}$) be the within- and cross-view covariances of the $\#$ and $G$
blocks conditional on $F$, and $c^\#=D_FF^\top\xi$, $c^G=\mu_eF^\top\xi$ their teacher cross-covariances. Then
\[
\E\Big[\mathbf 1_{\Omega_F}\max_{a,b}\|C^\#_{ab}-C^G_{ab}\|_{\mathrm{op}}\Big]+\E\Big[\mathbf 1_{\Omega_F}\|c^\#-c^G\|\Big]
\le C\sqrt{\frac{\log d}{d}} ,
\]
and the same bound holds for the clean covariances and cross-covariances of $x^{0,\#}$ and $x^{0,G}$.
\end{lemma}

\begin{proof}
By \eqref{eq:matching} and \eqref{eq:surrogate}, $C^\#_{11}-C^G_{11}=s^2(D_FG_FD_F-\kappa(1)^2G_F)+
(\Sigma_{\rm com}+\Sigma_{\rm priv}-\eta(1)I)$ and $C^\#_{12}-C^G_{12}=(D_FG_FD_F-\kappa(1)^2G_F)+(\Sigma_{\rm com}-\gamma(1)I)$.
Write $\Sigma_{\rm com}+\Sigma_{\rm priv}=\sum_{k\ge3\ \mathrm{odd}}A_k\mathscr R^{\circ k}A_k=\mathrm{diag}(\eta(q_r))+\Theta$
and $\Sigma_{\rm com}=\mathrm{diag}(\gamma(q_r))+\Theta_\rho$, where $\Theta,\Theta_\rho$ are the off-diagonal parts:
the diagonal of $\mathscr R^{\circ k}$ is one, so the diagonal entries are $\sum_{k\ge3}a_k(q_r)^2=\eta(q_r)$ and
$\sum_{k\ge3}a_k(q_r)^2\rho^k=\gamma(q_r)$. Three estimates:
(i) $\|D_F-\kappa(1)I\|=\max_r|\kappa(q_r)-\kappa(1)|\le C\max_r|q_r-1|$ on $\Omega_F$ by Lemma~\ref{lem:residualcov},
and likewise $\|\mathrm{diag}(\eta(q_r))-\eta(1)I\|$ and $\|\mathrm{diag}(\gamma(q_r))-\gamma(1)I\|$; since
$q_r\sim\chi^2_d/d$ are $m$ sub-exponential variables, $\E\max_r|q_r-1|\le C\sqrt{\log m/d}$
\cite[Section~2.8]{Vershynin}. (ii) $D_FG_FD_F-\kappa(1)^2G_F=(D_F-\kappa(1)I)G_FD_F+\kappa(1)G_F(D_F-\kappa(1)I)$ has
operator norm at most $C\max_r|q_r-1|$ on $\Omega_F$. (iii) For $r\ne s$,
$|\Theta_{rs}|\le\sum_{k\ge3}|a_k(q_r)a_k(q_s)||\mathscr R_{rs}|^k\le|\mathscr R_{rs}|^3\sqrt{V(q_r)V(q_s)}\le\|\sigma\|_\infty^2|\mathscr R_{rs}|^3$,
and $\mathscr R_{rs}=\langle f_r,f_s\rangle/(\|f_r\|\|f_s\|)$ is the cosine of the angle between two independent
Gaussian vectors in $\R^d$, so $\E\mathscr R_{rs}^6\le C/d^3$. Hence
$\E\|\Theta\|_{\mathrm{op}}\le\E\|\Theta\|_F\le\big(\sum_{r\ne s}\E\Theta_{rs}^2\big)^{1/2}\le Cm\,d^{-3/2}\le C d^{-1/2}$,
and the same for $\Theta_\rho$. The teacher term is $(D_F-\kappa(1)I)F^\top\xi$ with $\|F^\top\xi\|\le\|F\|$. The
clean case is identical with $\Delta=0$.
\end{proof}

\subsection{Proof of Theorem~\ref{thm:T1}}

\begin{proof}[Proof of \eqref{eq:T1a}]
For $0\le i\le n$ let $L^{(i)}$ be the objective in which blocks $1,\dots,i$ are of type $G$ and blocks
$i+1,\dots,n$ are actual, so $L^{(0)}=L_A$ and $L^{(n)}=L_G$, and let $L^{(i,\#)}$ be the objective in which
block $i$ is of type $\#$, blocks $<i$ of type $G$ and blocks $>i$ actual. Then
\[
\E\varphi(L_A^\star)-\E\varphi(L_G^\star)=\sum_{i=1}^n\Big(\E\varphi(L^{(i-1)\star})-\E\varphi(L^{(i,\#)\star})\Big)
+\sum_{i=1}^n\Big(\E\varphi(L^{(i,\#)\star})-\E\varphi(L^{(i)\star})\Big).
\]
Fix $i$ and let $\mathcal F_{-i}$ be the $\sigma$-field of $F$ and of all blocks other than $i$; the three
objectives compared share the leave-block minimum $L_-^\star$, the matrix $B$ and the vector $w_-$. By
Lemma~\ref{lem:ranktwo}, $L^\star=L_-^\star+\delta$ with $\delta=\frac1{2N}e_i^\top(I+Q_i)^{-1}e_i$ computed
from the block of the relevant type, and a Taylor expansion gives
\[
\big|\E[\varphi(L_-^\star+\delta)-\varphi(L_-^\star+\delta')\mid\mathcal F_{-i}]\big|
\le\|\varphi'\|_\infty\big|\E[\delta-\delta'\mid\mathcal F_{-i}]\big|+\|\varphi''\|_\infty\,\E[\delta^2+\delta'^2\mid\mathcal F_{-i}].
\]
By Lemma~\ref{lem:apriori}, $\E[\delta^2\mid\mathcal F_{-i}]\le\frac1{4N^2}\E[\|e_i\|^4\mid\mathcal F_{-i}]\le\frac{C(1+\|w_-\|^4)}{N^2}$
on $\Omega_F$, so the second-order terms sum to $C\|\varphi''\|_\infty/d$ after taking expectations. For the
first-order terms, write $Z=(t_i,u_1,u_2)$ with $u_a=x_{ia}^\top w_-$, so $e_i=t_i\mathbf 1_2-(u_1,u_2)^\top$, and freeze the
leverage: $\delta=\frac1{2N}\psi_{\bar Q}(Z)+\frac1{2N}e_i^\top\big((I+Q_i)^{-1}-(I+\bar Q_i)^{-1}\big)e_i$ with
$\psi_{\bar Q}(Z)=e_i^\top(I+\bar Q_i)^{-1}e_i$ a quadratic form in $Z$.

\emph{Actual to $\#$.} The conditional second moments of $Z$ under the actual and the $\#$ block coincide
(Proposition~\ref{prop:gaussianization}, with $w_-$ being $\mathcal F_{-i}$-measurable), and $\bar Q_i^A=\bar Q_i^\#$
(Lemma~\ref{lem:leverage2}); hence $\E[\psi_{\bar Q}(Z^A)\mid\mathcal F_{-i}]=\E[\psi_{\bar Q}(Z^\#)\mid\mathcal F_{-i}]$
exactly. The correction term is bounded by $\frac1{2N}\|e_i\|^2\|Q_i-\bar Q_i\|$, whose conditional expectation
is at most $\frac1{2N}\big(\E\|e_i\|^4\big)^{1/2}\big(\E\|Q_i-\bar Q_i\|_F^2\big)^{1/2}\le\frac{C(1+\|w_-\|^2)}{N\sqrt d}$
on $\Omega_F$, by Lemmas~\ref{lem:apriori} and~\ref{lem:leverage2}, for either block type. Summed over the $n$
steps this is $C\|\varphi'\|_\infty d^{-1/2}$.

\emph{$\#$ to $G$.} Both blocks are Gaussian conditional on $\mathcal F_{-i}$, and
$\E[\psi_{\bar Q}(Z)\mid\mathcal F_{-i}]$ is a fixed linear combination (with coefficients bounded by one) of
$\E t_i^2=1$, $\E[t_iu_a]=w_-^\top c_a$ and $\E[u_au_b]=w_-^\top C_{ab}w_-$, while $(\bar Q_i)_{ab}=\frac1N\Tr(BC_{ab})$.
By Lemma~\ref{lem:bridge}, on $\Omega_F$ the moments differ by at most
$C(1+\|w_-\|^2)\max_{ab}\|C^\#_{ab}-C^G_{ab}\|+C\|w_-\|\|c^\#-c^G\|$, and the frozen leverages by at most
$\frac1N\|B\|_{\mathrm{tr}}\max_{ab}\|C^\#_{ab}-C^G_{ab}\|\le\frac{m}{N\lambda}\max_{ab}\|C^\#_{ab}-C^G_{ab}\|$, which enters
$\psi$ through the Lipschitz bound of Lemma~\ref{lem:leverage2} multiplied by $\|e_i\|^2$. Taking expectations
and using Lemma~\ref{lem:apriori} for the moments of $w_-$ and $e_i$, each of the $n$ steps contributes at most
$\frac{C}{N}\sqrt{\log d/d}\,\|\varphi'\|_\infty$ to the first-order sum.

Off $\Omega_F$ the two expectations differ by at most $2\|\varphi\|_\infty\Pr(\Omega_F^c)\le C\|\varphi\|_\infty e^{-cd}$.
Adding the three contributions gives \eqref{eq:T1a}.
\end{proof}

\begin{proof}[Proof of \eqref{eq:T1b}]
Let $M=\lceil d\rceil$ and let $(t^0_l,a^0_l)_{l\le M}$ be an independent clean block, with actual features
$x^{0,A}_l=\sigma(F^\top a^0_l)$; put $c_0=\frac1M\sum_l(t^0_l)^2$, $r_0=\frac1M\sum_lt^0_lx^{0,A}_l$,
$S_0=\frac1M\sum_lx^{0,A}_l(x^{0,A}_l)^\top$ and $\widehat R(w)=c_0-2r_0^\top w+w^\top S_0w$, the empirical clean risk of
$w$ on the block with actual features. Four comparisons, each in $L^1$ after multiplication by
$\|\varphi'\|_\infty$ (with the second-order and off-$\Omega_F$ terms as before):

(1) $R_A(w_A^\star\mid F)$ against $\widehat R(w_A^\star)$: conditional on $F$ and $w_A^\star$ the block is
independent of $w_A^\star$, and $\widehat R(w)-R_A(w\mid F)$ is an average of $M$ independent centred variables
with variance at most $C(1+\|w\|^4)$ by Lemma~\ref{lem:apriori}; so the $L^1$ difference is at most
$C M^{-1/2}\le Cd^{-1/2}$.

(2) $\widehat R(w_A^\star)$ against $\widehat R(w_G^\star)$: the Lindeberg scheme of the previous proof with the
statistic $\widehat R(w^\star)$. Conditional on $\mathcal F_{-i}$ and the clean block, with $v_-=B(S_0w_--r_0)$ and
$R_i=\frac1NX_iBS_0BX_i^\top$, the second formula in \eqref{eq:ranktwo} gives the exact increment
\[
\widehat R(w^\star)-\widehat R(w_-)=\frac1N\Big[2e_i^\top(I+Q_i)^{-1}X_iv_-+e_i^\top(I+Q_i)^{-1}R_i(I+Q_i)^{-1}e_i\Big],
\]
a quadratic form in $Z'=(t_i,u_1,u_2,v_1,v_2)$ with $v_a=x_{ia}^\top v_-$, once $Q_i$ and $R_i$ are frozen at their
conditional means. On the event $\{\|S_0\|\le C\}$ (of probability $1-Ce^{-cd}$ by \cite[Theorem~4.6.1]{Vershynin}
and Lemma~\ref{lem:apriori}), $\|v_-\|\le C(1+\|w_-\|)$, the second moments of $Z'$ match exactly between the
actual and the $\#$ block (Proposition~\ref{prop:gaussianization}, applied to the $\mathcal F_{-i}$-measurable
vectors $w_-$ and $v_-$), the frozen means of $Q_i$ and $R_i$ match exactly, the fluctuations of $Q_i$ and $R_i$
are $O(1/d)$ in mean square (Lemma~\ref{lem:leverage2}), and the $\#$-to-$G$ step is controlled by
Lemma~\ref{lem:bridge} exactly as before, with $\frac1N\|BS_0B\|_{\mathrm{tr}}\le\frac{m\|S_0\|}{N\lambda^2}$. Summing
over $n$ blocks gives $C\|\varphi'\|_\infty\sqrt{\log d/d}+C\|\varphi''\|_\infty/d$.

(3) $\widehat R(w_G^\star)$ against $R_\#(w_G^\star\mid F)=\E[(t^0-(x^{0,A})^\top w_G^\star)^2\mid F,w_G^\star]$: the
empirical-to-population step (1) again, $O(d^{-1/2})$.

(4) $R_\#(w_G^\star\mid F)$ against $R_G(w_G^\star\mid F)$: by Proposition~\ref{prop:gaussianization} the first
equals the population risk under the exact clean Gaussian row, and the two Gaussian clean risks differ by
$|2(w_G^\star)^\top(c_0^\#-c_0^G)|+|(w_G^\star)^\top(C_0^\#-C_0^G)w_G^\star|$, which is $O(\sqrt{\log d/d})$ in $L^1$ by
Lemma~\ref{lem:bridge} and Lemma~\ref{lem:apriori}.

Adding the four bounds gives \eqref{eq:T1b}.
\end{proof}

\begin{remark}\label{rem:T1}
The factor $\sqrt{\log d}$ in \eqref{eq:T1a}--\eqref{eq:T1b} comes from a single place, the maximum of the
$m$ column norms $|q_r-1|$ in Lemma~\ref{lem:bridge}(i). Replacing that maximum by an average over the
coordinates of $w_-$ requires a delocalisation bound $\E|w_{-,r}|^4\le Cd^{-2}$ on the coordinates of the
leave-block minimiser, under which the rate becomes $d^{-1/2}$ \cite{PN24}; the equivalence itself does not
depend on it. Two of the objects in the proof are measured in Section~\ref{sec:t1}: the leverage matrix
$Q_i$ of Lemma~\ref{lem:ranktwo}, whose fluctuation Lemma~\ref{lem:leverage2} bounds by $C/d$, and the
Hermite coefficients \eqref{eq:hermitecoef}, which determine the surrogate and whose residual variances
$\eta(1)-\gamma(1)$ and $\gamma(1)$ are the quantities behind Proposition~\ref{prop:hermite} there.
\end{remark}


\part{Finite-size measurements}

\section{Target 3: the finite-block certificate}\label{sec:t3}

\subsection{What the proof uses}

For zero signal the excess risk of the minimum-norm interpolant is a quadratic form in the noise, and
its expectation over the noise is the leverage functional
\begin{equation}
\frac{V(X)}{\sigma^2} \;=\; \E_U\big\|k(U,X)K^{-1}\big\|^2
\;=\; \Tr\!\big(K^{-1}GK^{-1}\big),
\qquad G=\sum_s \lambda_s^2N_sP_s,
\label{eq:V}
\end{equation}
the second equality because $\E_U[k(X,U)k(U,X)]=\sum_s\lambda_s^2 N_sP_s$ by Mercer and the addition
theorem. Writing $K_s=\lambda_sN_sP_s$ and $t_s=\Tr(K^{-1}K_s)$, Cauchy--Schwarz in the Frobenius
inner product gives $t_s^2\le nN_s\lambda_sq_s$ with $q_s=\Tr(K^{-1}K_sK^{-1})$, and summing over a
window $B$ of degrees,
\begin{equation}
\frac{V(X)}{\sigma^2}\;\ge\;\frac{T_B^2}{n\,D_B},
\qquad T_B=\sum_{s\in B}t_s,\quad D_B=\sum_{s\in B}N_s .
\label{eq:window}
\end{equation}

The proof of Target~3 \cite{PN24} takes $B$ to be the \emph{cumulative} block $0,\dots,L$, writes
$D_L=\sum_{s\le L}N_{d,s}$, $\delta_L=D_L/n$, $T_L=\Tr(K^{-1}\sum_{s\le L}K_s)$, and bounds the
leverage the block captures in expectation over the design,
\begin{equation}
\E\,T_L\;\ge\;\frac{n\,D_L}{D_L+n-1+\tau_{>L}/\lambda_{d,L}},
\qquad \tau_{>L}=\sum_{s>L}\tau_s ,
\label{eq:capture}
\end{equation}
with no concentration step: only operator monotonicity of the inverse, Cauchy--Schwarz, the Mercer
decomposition, the addition theorem and Jensen's inequality enter. In the critical regime the tail ratio $\theta_L=\tau_{>L}/(n\lambda_{d,L})$ is $o(1)$,
so \eqref{eq:capture} reads $\E\,T_L/n\ge\delta_L/(1+\delta_L+o(1))$, and taking expectations in
\eqref{eq:window} and applying Jensen once more,
\begin{equation}
\E\,\frac{V(X)}{\sigma^2}\;\ge\;\frac{(\E\,T_L/n)^2}{\delta_L}\;\ge\;
\frac{\delta_L}{(1+\delta_L)^2}-o(1).
\label{eq:onewindow}
\end{equation}
For the largest $L$ with $D_L\le n$ the ratio $\delta_L$ lies in $(1/5,1]$ for all large $d$, by the
exact multiplicity ratio and Stirling's formula, and $y\mapsto y/(1+y)^2$ is increasing on $[0,1]$,
so
\[
\liminf_{d\to\infty}\E\,\frac{V(X)}{\sigma^2}\;\ge\;\frac{1/5}{(6/5)^2}=\frac{5}{36}=c_1 ,
\]
for every fixed $h>0$. This is the one-window theorem. Splitting the degrees into nested cumulative
windows $B_0\subset\dots\subset B_J$ with $y_j=D_{B_j}/n$ and $A_j=T_{B_j}/n$, the pathwise
reduction \eqref{eq:window} applied shell by shell gives
$V/\sigma^2\ge A_0^2/y_0+\sum_{j\ge1}(A_j-A_{j-1})^2/(y_j-y_{j-1})$, a convex quadratic form in the
$A_j$, and the same expectation bounds inserted into it, with the window offsets chosen against the
phase of $n=2^d$ relative to the harmonic multiplicities, give the ten-window certificate
\[
c_{10}\;=\;\frac{3882594969529}{12500000000000}\;=\;\numCtenLong\ldots ,
\]
an exact rational number, again for every fixed $h>0$ \cite{PN24}. The one-window constant is Theorem~\ref{thm:T3}; the ten-window certificate is a strengthening
obtained in \cite{PN24} by running the argument of Section~\ref{sec:proofT3} over ten blocks at once
(Remark~\ref{rem:sharper}), and is quoted here as a comparison value. No random-matrix universality
enters at any point. Both constants are independent of $h$; the measured functional is not, and the distance
between them is measured below.

\subsection{What is measured}

Every object in the last subsection is computable from the Gram matrix alone. $P_s$ evaluated at the
matrix of inner products follows the dimension-$d$ Legendre recurrence
$(s+d-2)P_{s+1}(t)=(2s+d-2)tP_s(t)-sP_{s-1}(t)$, so no spherical harmonic is ever formed; the
projection kernel $\Pi_L=\sum_{s\le L}N_sP_s$ gives $\|S_L-I\|_F^2=\|\Pi_L\|_F^2/n^2-D_L$ for the
block's orthonormal feature matrix $\Phi_L$ and $S_L=\frac1n\Phi_L^\top\Phi_L$; and $G$ in
\eqref{eq:V} is another accumulation over the same stream. For each design and each $L$ the driver
records $D_L$, $\delta_L$, $T_L/n$, the certificate $(1/\delta_L)(T_L/n)^2$, and the ratio of the
measured $\|S_L-I\|_F^2$ to its exact expectation
\begin{equation}
\E\big\|S_L-I\big\|_F^2=\frac{D_L(D_L-1)}{n},
\label{eq:second}
\end{equation}
an elementary identity that serves as a check on the Gram-only implementation. It reports
$V(X)/\sigma^2$ twice: once as a Monte-Carlo average over fresh test points, once exactly through
\eqref{eq:V}.

Carrying both is what exposes a numerical trap. At $h=1/2$ and $h=1$ the two derivations agree to a
fraction of a percent. At $h=2$ the form $\Tr(K^{-1}GK^{-1})$ evaluated as written does not agree
with the Monte-Carlo value, and the disagreement grows with $d$; the identity is not at fault, its
evaluation is. In the eigenbasis of $K$ the trace reads $\sum_k \tilde G_{kk}/w_k^2$, so rounding of
size $\varepsilon\|G\|$ in the entries of $G$ is amplified by $\varepsilon\,\kappa(K)^2$, not by
$\varepsilon\kappa$: the inverse appears twice. The Gaussian kernel at $h=2$ is the smoothest of the
three bandwidths and $K$ the worst conditioned, and Table~\ref{tab:traceform} shows
$\varepsilon\kappa^2$ climbing from $3$ at $d=9$ to $2\times10^3$ at $d=11$, where the trace form has
departed from the Monte-Carlo control by two percent while the factored evaluation described next
stays within a tenth of a percent of it. Factoring $G$ alone does not help, because the rounding
directions of $G$ carry eigenvalues of order $\varepsilon\|G\|$ that are still divided by $w_k^2$.
The evaluation that survives uses the \emph{same} factors for both matrices: with
$N_sP_s(T)=Y_sY_s^{\!\top}$ from one eigendecomposition per degree, $K=\sum_s\lambda_sY_sY_s^{\!\top}$
and $G=\sum_s\lambda_s^2Y_sY_s^{\!\top}$, so that
\[
V/\sigma^2=\sum_{s\le L}\lambda_s^2\,\|K^{-1}Y_s\|_F^2 ,
\]
a form with one inverse in which rounding enters squared, and in which each factor's own rounding
directions are separated from its true spectrum by many orders of magnitude and dropped. The exact
column of Table~\ref{tab:t3band} is evaluated this way wherever $\varepsilon\kappa^2$ exceeds
$10^{-4}$, and both forms are kept in the records with the criterion. The point is general: an exact
trace identity carrying two inverses is not an exact number once $\varepsilon\kappa^2$ is of order
one, and for a smooth kernel that happens at modest $n$.

\begin{table}[t]
\centering\small
\caption{The two evaluations of \eqref{eq:V} at $h=2$, $n=2^d$, means over the designs of
Table~\ref{tab:t3band}. $\varepsilon\kappa^2$ is machine precision times the squared condition
number of $K$; the trace form is $\Tr(K^{-1}GK^{-1})$ as written, the factor form uses the shared
per-degree factors, and the Monte-Carlo column averages $\|k(U,X)K^{-1}\|^2$ over $4{,}000$ fresh
points per design.}
\label{tab:traceform}
\begin{tabular}{rrlrrr}
\toprule
$d$ & designs & $\varepsilon\kappa^2$ & trace form & factor form & Monte Carlo\\
\midrule
% BEGIN AUTOTABLE tab:traceform
9 & 40 & 3.2 & 3.1941 & 3.1941 & 3.1925\\
10 & 24 & 84 & 3.8754 & 3.8753 & 3.8757\\
11 & 16 & $2.0\times10^{3}$ & 2.5691 & 2.5157 & 2.5138\\
% END AUTOTABLE tab:traceform
\bottomrule
\end{tabular}
\end{table}

\begin{table}[t]
\centering\small
\caption{The cumulative blocks at $d=10$, $n=2^{10}$, $h=1$, means over $\numTthreeBlockDesigns$
designs (the standard deviation across designs of $T_L/n$ and of the certificate is at most
$\numTthreeSdMax$; that of the identity ratio falls from $\numTthreeIdSdFirst$ at
$L=\numTthreeIdLfirst$ to $\numTthreeIdSdLast$ at $L=\numTthreeIdLlast$). $\delta_L=D_L/n$;
the certificate is $(1/\delta_L)(T_L/n)^2$; the two columns headed by functions of $\delta_L$ are the
theorem's asymptotic lower bounds \eqref{eq:capture} and \eqref{eq:onewindow} for $T_L/n$ and for the
certificate; the identity column is the measured $\|S_L-I\|_F^2$ divided by the exact
\eqref{eq:second}; the last column marks the block the one-window theorem uses, the largest $L$ with
$D_L\le n$.}
\label{tab:t3}
\begin{tabular}{rrrrrrrrc}
\toprule
$L$ & $D_L$ & $\delta_L$ & $T_L/n$ & $\frac{\delta_L}{1+\delta_L}$ & certificate &
$\frac{\delta_L}{(1+\delta_L)^2}$ & \eqref{eq:second} & theorem's block\\
\midrule
% BEGIN AUTOTABLE tab:t3
1 & 11 & 0.0107 & 0.0107 & 0.0106 & 0.0107 & 0.0105 & 0.982 & \\
2 & 65 & 0.0635 & 0.0634 & 0.0597 & 0.0633 & 0.0561 & 1.003 & \\
3 & 275 & 0.2686 & 0.2636 & 0.2117 & 0.2588 & 0.1669 & 1.001 & \\
4 & 935 & 0.9131 & 0.7330 & 0.4773 & 0.5884 & 0.2495 & 1.000 & yes\\
5 & 2717 & 2.6533 & 0.9639 & 0.7263 & 0.3502 & 0.1988 & 1.000 & \\
6 & 7007 & 6.8428 & 0.9964 & 0.8725 & 0.1451 & 0.1112 & 0.999 & \\
% END AUTOTABLE tab:t3
\bottomrule
\end{tabular}
\end{table}

\begin{table}[t]
\centering\small
\caption{The bounded functional $V(X)/\sigma^2$ at $n=2^d$, mean $\pm$ standard deviation over
independent designs, and its ratio to the ten-window constant $c_{10}=\numCten$. The last three
columns give the one-window theorem's block $L_-$, the largest $L$ with $D_L\le n$, its $\delta_L$,
and the mean certificate $(1/\delta_L)(T_L/n)^2$ at that block, to be read against
$c_1=5/36=\numCone$. The cells $(d,h)=(12,2)$, $(13,2)$ and all of $d=14$ were not run.}
\label{tab:t3band}
\begin{tabular}{rrlrrrrr}
\toprule
$d$ & $h$ & $V(X)/\sigma^2$ & designs & $/\,c_{10}$ & $L_-$ & $\delta_{L_-}$ & certificate at $L_-$\\
\midrule
% BEGIN AUTOTABLE tab:t3band
9 & 0.5 & $0.5371\pm0.0067$ & 40 & 1.73 & 3 & 0.410 & 0.214\\
9 & 1 & $1.7743\pm0.0468$ & 60 & 5.71 & 3 & 0.410 & 0.373\\
9 & 2 & $3.1925\pm0.1125$ & 40 & 10.28 & 3 & 0.410 & 0.403\\
10 & 0.5 & $0.5621\pm0.0032$ & 24 & 1.81 & 4 & 0.913 & 0.285\\
10 & 1 & $1.8357\pm0.0246$ & 36 & 5.91 & 4 & 0.913 & 0.588\\
10 & 2 & $3.8757\pm0.0775$ & 24 & 12.48 & 4 & 0.913 & 0.760\\
11 & 0.5 & $0.5891\pm0.0023$ & 16 & 1.90 & 4 & 0.628 & 0.279\\
11 & 1 & $1.7087\pm0.0144$ & 24 & 5.50 & 4 & 0.628 & 0.517\\
11 & 2 & $2.5138\pm0.0269$ & 16 & 8.09 & 4 & 0.628 & 0.598\\
12 & 0.5 & $0.6125\pm0.0016$ & 8 & 1.97 & 4 & 0.422 & 0.250\\
12 & 1 & $1.7047\pm0.0096$ & 12 & 5.49 & 4 & 0.422 & 0.390\\
13 & 0.5 & $0.6384\pm0.0010$ & 4 & 2.06 & 5 & 0.978 & 0.320\\
13 & 1 & $1.8797\pm0.0061$ & 6 & 6.05 & 5 & 0.978 & 0.611\\
% END AUTOTABLE tab:t3band
\bottomrule
\end{tabular}
\end{table}

Four things are visible in Tables~\ref{tab:t3} and~\ref{tab:t3band}. The identity \eqref{eq:second}
holds: the ratio sits within $2\%$ of one from $D_L=65$ on and within a tenth of a percent from
$D_L=275$ on. The leverage the block captures exceeds the expectation bound \eqref{eq:capture} at
every $L$ with $D_L\le n$, and by a wide margin at the theorem's block: at $d=10$ the block $L=4$ has
$\delta_L=0.913$ and captures $T_L/n=0.733$ against the bound's $0.477$, so at these sizes
\eqref{eq:capture} is the loose step. The certificate at the theorem's block sits above $c_1=5/36$
at every dimension and bandwidth measured, between $\numCertMin$ and $\numCertMax$; its variation
with $d$ is a sawtooth that follows $\delta_{L_-}$, which is set by where $2^d$ falls between two
block boundaries, and not a trend. And the quantity being bounded exceeds the ten-window constant
$c_{10}$ by a factor between $\numRatioMinHalf$ and $\numRatioMaxHalf$ at $h=1/2$, between
$\numRatioMinOne$ and $\numRatioMaxOne$ at $h=1$, and between $\numRatioMinTwo$ and
$\numRatioMaxTwo$ at $h=2$: the functional is monotone in the bandwidth while the constant is not,
and the gap between them is where an $h$-dependent constant would live.

The two derivations of $V/\sigma^2$ agree throughout: the largest relative discrepancy over all
designs is $\numVgapMax$, and the mean values agree to within $\numVmeanGapMax$ percent. The Mercer mass
$\sum_s\tau_s$ recovers to $1.000000$ at $h\ge1$ and to between $\numTauHalfMin$ and
$\numTauHalfMax$ at $h=1/2$, where the eigenvalue tail reaches past the harmonic truncation.

\section{The dimension gap between Targets 2 and 3}\label{sec:gap}

Targets~2 and~3 are the same construction at two ends of one number. Let $k(d,n)=\min\{s:N_{d,s}\ge n\}$
be the critical degree and
\[
\delta \;=\; \frac{\log\big(N_{d,k}/n\big)}{\log d}.
\]
Target~2's sequence has $\log N_{d,k}/\log d \to R+\delta_\star$ with
$\delta_\star = 1/\log2-1=0.442695\ldots$, so it runs at $\delta\to\delta_\star>0$, and the proof's
tail-isotropy step turns on exactly that positive gap: the harmonic dimension above the cutoff
outruns the sample by $d^{\delta_\star}$, which is what makes $\|\Kt/\tau-I\|\to0$ for the tail
block $\Kt$ of the kernel matrix and its Mercer mass $\tau$. Target~3's $n=2^d$ puts the crossing at
$k\sim c_\star d$ with $c_\star=0.2938\ldots$ solving $(1+c)\log(1+c)-c\log c=\log2$, and there the
gap closes. The same coordinate organises the large-dimensional theory of kernel ridge regression
\cite{Ghorbani,Misiakiewicz,LuLinXiaHuang}: with $n\asymp d^\gamma$ the fractional part of $\gamma$
is the gap of the last resolved harmonic block.

The natural experiment is to fix $d$, sweep $n$, and watch the noise risk against $\delta$. One
feature of the geometry has to be respected first: $\delta$ is not monotone in $n$. Each time $n$
crosses a multiplicity $N_{d,k}$ the critical degree advances and $\delta$ resets to the top of its
range, so a plain sweep in $n$ traces a sawtooth in $\delta$ while the risk climbs steadily with
sample size for reasons that have nothing to do with the gap. Comparing across that sawtooth measures
$n$, not $\delta$. The comparison has to be made \emph{within} a block, at fixed critical degree.

Done that way the ordering is clean. Table~\ref{tab:gap} shows one dimension in full: at $d=8$,
$h=1/2$, at fixed critical degree $k$, the mean noise risk falls monotonically as $\delta$ rises, in
every one of the six blocks, over four independent designs per cell. Table~\ref{tab:gapsum}
summarises the same sweep at seven dimensions from $d=8$ to $d=64$, all at $h=1/2$: over
$\numGapBlocks$ blocks with at least three values of the gap, $\numGapInv$ of the $\numGapPairs$
ordered pairs of cells are inverted, and every inversion lies either in a block with
$N_{d,k}\le\numGapInvNmax$, where the noise risk is at most $\numGapInvRiskMax$ and the two inverted
cells differ by at most $\numGapInvDiffMax$, or at $d=64$, where only two designs were run per cell.
Across blocks the risk rises with $n$, as it must. From $d=12$ on the ladder reaches $\delta=0.44$, the value of $\delta_\star$, and at
$d=24$, $32$ and $48$ it reaches $\delta=0.6$ (Table~\ref{tab:gapdeep}); the risk keeps falling
through both, with no turnover at any dimension.

The tail isotropy $\|\Kt/\tau-I\|$, the quantity Target 2's proof drives to zero, follows the same
ordering (Table~\ref{tab:gapiso}), and strictly so in all but one of the blocks of
Table~\ref{tab:gapsum}. It is an operator norm, a maximum over eigenvalues, and a maximum is a
noisier estimate than an average at the same number of designs; the one block in which it inverts
does so between the two smallest gaps, where its means differ by one percent.

What this does and does not show is worth stating carefully. The two theorems constrain the two ends
of the order parameter --- one proves a vanishing risk along a sequence with $\delta$ bounded away
from zero, the other a positive floor along a sequence where the gap closes --- but neither implies
anything about the monotonicity of a finite-size curve at fixed $d$. The measured ordering is
consistent with them and is the shape one would expect if $\delta$ really is the controlling
variable; it is evidence for that reading, not a consequence of either proof. All of it is at one
bandwidth.

\begin{table}[t]
\centering
\caption{The dimension gap at fixed critical degree, $d=8$, $h=1/2$, four designs per cell. Each row
is one harmonic block: $n$ runs over the sample sizes whose critical degree is $k$, and the entries
are the mean noise risk $\mathcal R_{\mathrm{noise}}/\sigma^2$ at each value of the gap $\delta$.
Within every row the risk falls as $\delta$ rises.}
\label{tab:gap}
\begin{tabular}{rrrrrrr}
\toprule
$k$ & $N_{d,k}$ & $\delta=0.02$ & $0.05$ & $0.10$ & $0.20$ & $0.30$\\
\midrule
% BEGIN AUTOTABLE tab:gap
5 & 672 & 0.772 & 0.754 & 0.705 & 0.697 & 0.627\\
6 & 1386 & 1.125 & 1.088 & 0.989 & 0.885 & 0.819\\
7 & 2640 & 1.482 & 1.337 & 1.286 & 1.207 & 1.126\\
8 & 4719 & 1.760 & 1.668 & 1.609 & 1.487 & \\
9 & 8008 & 2.054 & 1.983 & 1.967 & 1.802 & \\
10 & 13013 & 2.428 & 2.389 & 2.257 & 2.147 & \\
% END AUTOTABLE tab:gap
\bottomrule
\end{tabular}
\end{table}

\begin{table}[t]
\centering
\caption{The tail isotropy $\|\Kt/\tau-I\|$ in the cells of Table~\ref{tab:gap}.}
\label{tab:gapiso}
\begin{tabular}{rrrrrrr}
\toprule
$k$ & $N_{d,k}$ & $\delta=0.02$ & $0.05$ & $0.10$ & $0.20$ & $0.30$\\
\midrule
% BEGIN AUTOTABLE tab:gapiso
5 & 672 & 1.897 & 1.893 & 1.844 & 1.533 & 1.447\\
6 & 1386 & 2.130 & 2.158 & 1.913 & 1.765 & 1.615\\
7 & 2640 & 2.302 & 2.285 & 2.214 & 1.828 & 1.650\\
8 & 4719 & 2.439 & 2.421 & 2.174 & 2.023 & \\
9 & 8008 & 2.648 & 2.435 & 2.310 & 1.976 & \\
10 & 13013 & 2.783 & 2.664 & 2.388 & 2.105 & \\
% END AUTOTABLE tab:gapiso
\bottomrule
\end{tabular}
\end{table}

\begin{table}[t]
\centering\small
\caption{Within-block ordering at every dimension measured, $h=1/2$. A block is a critical degree
$k$ at which at least three values of $\delta$ were run; a pair of cells in a block is inverted when
the cell with the larger gap has the larger mean noise risk. The isotropy column counts the blocks in
which $\|\Kt/\tau-I\|$ falls strictly along the whole row; the last column names the blocks that
carry an inverted risk pair.}
\label{tab:gapsum}
\begin{tabular}{rrlrrl}
\toprule
$d$ & blocks & designs per cell & inverted pairs & isotropy monotone & blocks with an inversion\\
\midrule
% BEGIN AUTOTABLE tab:gapsum
8 & 6 & 4 & 0 / 48 & 5 / 6 & --\\
12 & 4 & 4 & 2 / 43 & 4 / 4 & $k=3$, $k=4$\\
16 & 3 & 4 & 0 / 35 & 3 / 3 & --\\
24 & 2 & 4 & 0 / 27 & 2 / 2 & --\\
32 & 2 & 4 & 2 / 27 & 2 / 2 & $k=2$\\
48 & 2 & 4 & 0 / 16 & 2 / 2 & --\\
64 & 1 & 2 & 4 / 15 & 1 / 1 & $k=2$\\
% END AUTOTABLE tab:gapsum
\bottomrule
\end{tabular}
\end{table}

\begin{table}[t]
\centering\small
\caption{The blocks whose ladder reaches $\delta\ge0.44$: mean noise risk $\mathcal
R_{\mathrm{noise}}/\sigma^2$ at $h=1/2$, four designs per cell, out to $\delta=0.6$ where it was
run.}
\label{tab:gapdeep}
\begin{tabular}{rrrrrrrrrr}
\toprule
$d$ & $k$ & $N_{d,k}$ & $\delta=0.02$ & $0.05$ & $0.10$ & $0.20$ & $0.30$ & $0.44$ & $0.60$\\
\midrule
% BEGIN AUTOTABLE tab:gapdeep
12 & 4 & 1287 & 0.393 & 0.388 & 0.392 & 0.369 & 0.318 & 0.271 & --\\
12 & 5 & 4004 & 0.619 & 0.576 & 0.568 & 0.517 & 0.487 & 0.424 & --\\
16 & 4 & 3740 & 0.347 & 0.331 & 0.316 & 0.300 & 0.277 & 0.228 & --\\
16 & 5 & 14688 & -- & 0.478 & 0.468 & 0.436 & 0.407 & 0.366 & --\\
24 & 3 & 2576 & 0.160 & 0.153 & 0.150 & 0.135 & 0.117 & 0.096 & 0.079\\
24 & 4 & 17250 & -- & -- & 0.247 & 0.230 & 0.212 & 0.188 & --\\
32 & 3 & 5952 & 0.135 & 0.135 & 0.120 & 0.117 & 0.103 & 0.084 & 0.070\\
48 & 3 & 19552 & -- & -- & -- & 0.095 & 0.087 & 0.075 & 0.067\\
64 & 2 & 2079 & 0.042 & 0.045 & 0.045 & 0.035 & 0.025 & 0.025 & --\\
% END AUTOTABLE tab:gapdeep
\bottomrule
\end{tabular}
\end{table}

\section{Target 4: the Veronese reduction and its two lemmas}\label{sec:t4}

\subsection{The reduction}

Write $x=ru$ with $u$ uniform on $\Sph$ and $r^2\sim\chi^2_d$ independent, and identify $z(x)$ with
$Q=(xx^\top-I)/\sqrt2\in\mathrm{Sym}_d$ in an orthonormal Frobenius basis whose first element is
$E_0=I/\sqrt d$. Then the trace coordinate is $c(r)=(r^2-d)/\sqrt{2d}$ and each traceless coordinate
is $a(r)h_j(u)$ with $a(r)=r^2/\sqrt{d(d+2)}$ and
$h_j(u)=\sqrt{d(d+2)/2}\;u^\top E_ju$, so that after an orthogonal change of feature coordinates
\begin{equation}
ZO=D\,[\,b,\;H\,],\qquad D=\mathrm{diag}(a_i),\quad b=D^{-1}c,
\label{eq:reduction}
\end{equation}
with $b$ \emph{exactly} independent of $H$, and every row of $H$ centred, isotropic and of fixed norm
$m=p-1$. The scalar column satisfies $b=\sqrt{(d+2)/2}\,(r^2-d)/r^2$, whence for $d>4$
\[
\E b=-\frac{\sqrt{2(d+2)}}{d-2},\qquad
\Var(b)=\frac{d^2(d+2)}{(d-4)(d-2)^2}=1+\frac{10}{d}+O(d^{-2}),
\]
and $\max_i|a_i-1|=O_p(\sqrt{\log d/d})$, so the hard edge of $Z$ and of $[b,H]$ coincide in the
limit. In the basis of an SVD of $H$ the smallest eigenvalue of $[b,H][b,H]^\top$ obeys an exact
rank-one secular equation, and the leave-one-out normal has trace weight
$\alpha^2=(1+\|H_-^{-1}b_-\|^2)^{-1}$ \cite{PN24}.

The reduction names two lemmas whose truth would materially advance the transfer:
\begin{description}
\item[(L-a)] $|u_j^\top\mathbf{1}|=o_p(\sqrt d)$ for the bottom left singular vectors $u_j$ of $H$
(this suffices to kill the mean shift of $b$; $O_p(1)$ is plausible);
\item[(L-b)] $\E\Tr(A^4)=O(1/d)$ for the feature-space normal $A$ of $p-1$ rows, normalized to
$\|A\|_F=1$.
\end{description}

\subsection{What is measured}

The identity \eqref{eq:reduction} is verifiable exactly, and is verified:
$s_{\min}(Z)$ and $s_{\min}(D[b,H])$ agree, in the median over designs, to $\numRedMedFirst$
relative at $d=\numHardDmin$ and to $\numRedMedLast$ at $d=\numHardDmax$, the discrepancy growing
with $d$ as the conditioning does; every row of $H$ has $\|h(u)\|^2=m$ to $10^{-15}$, and the mean
and variance of the scalar
column $b$ match their closed forms to within $\numBmomWorst$ percent at every dimension. The
asymptotic half --- that dropping $D$ does not move the edge --- converges slowly but visibly: the
ratio $s_{\min}(Z)/s_{\min}([b,H])$ in Table~\ref{tab:t4} rises from $\numSminRatioFirst$ at
$d=\numHardDmin$ to $\numSminRatioLast$ at $d=\numHardDmax$, with a spread across designs that shrinks
from $\numSminSdFirst$ to $\numSminSdLast$, and $d$ times its distance from one stays between
$\numDratioMin$ and $\numDratioMax$ over the whole range, so the diagonal scaling moves the edge by
about $3.9/d$ at these sizes.

Both lemmas look true, and (L-a) in its stronger form (Table~\ref{tab:t4}). The overlaps
$|u_j^\top\mathbf 1|$ show no growth with $d$ at all --- $\numOverlapFirst$ at $d=\numHardDmin$ and
$\numOverlapLast$ at $d=\numHardDmax$, with a standard deviation near $0.6$ throughout --- which is
$O_p(1)$, not merely the $o_p(\sqrt d)$ the reduction asks for.

\begin{table}[t]
\centering\small
\caption{Target 4 lemma diagnostics. $A$ is the unit-Frobenius-norm feature-space normal of $p-1$
rows; $u_1,u_2,u_3$ are the bottom left singular vectors of $H$ and $\mathbf 1\in\R^p$. Entries are
mean $\pm$ standard deviation over the designs, with the $90$th percentile of the overlap beside it;
extremes are not reported, since the width of a min--max grows with the number of designs. The
seventh column is
$s_{\min}(Z)/s_{\min}([b,H])$, the effect of the diagonal scaling $D$ in \eqref{eq:reduction}. A
third column tending to a positive limit is (L-b); a fifth column that does not grow like $\sqrt d$
is (L-a).}
\label{tab:t4}
\begin{tabular}{rrlllrlr}
\toprule
$d$ & $p$ & $d\,\Tr(A^4)$ & $\sqrt d\,\|A\|_{\mathrm{op}}$ & $|u_j^\top\mathbf 1|$ & p90 &
$s_{\min}$ ratio & designs\\
\midrule
% BEGIN AUTOTABLE tab:t4
20 & 210 & $2.37\pm0.24$ & $2.03\pm0.15$ & $0.71\pm0.54$ & 1.47 & $0.815\pm0.040$ & 1200\\
30 & 465 & $2.24\pm0.15$ & $2.02\pm0.12$ & $0.73\pm0.56$ & 1.49 & $0.873\pm0.022$ & 900\\
40 & 820 & $2.18\pm0.09$ & $2.02\pm0.10$ & $0.74\pm0.56$ & 1.54 & $0.902\pm0.014$ & 600\\
60 & 1830 & $2.12\pm0.06$ & $2.02\pm0.08$ & $0.77\pm0.59$ & 1.59 & $0.934\pm0.007$ & 420\\
80 & 3240 & $2.09\pm0.04$ & $2.01\pm0.06$ & $0.79\pm0.60$ & 1.66 & $0.951\pm0.005$ & 270\\
100 & 5050 & $2.07\pm0.03$ & $2.02\pm0.05$ & $0.79\pm0.60$ & 1.63 & $0.961\pm0.004$ & 240\\
120 & 7260 & $2.06\pm0.03$ & $2.02\pm0.05$ & $0.80\pm0.58$ & 1.53 & $0.967\pm0.002$ & 36\\
% END AUTOTABLE tab:t4
\bottomrule
\end{tabular}
\end{table}

(L-b) is sharper than a glance suggests. $d\,\Tr(A^4)$ is not constant (Table~\ref{tab:lb}): the
values at $d=20$ and $d=40$ are $\numLbEndsSe$ combined standard errors apart, so a flat fit is
rejected outright. It is converging, with a $1/d$ correction. Fitting $d\,\Tr(A^4)=a+b/d$ to the two
end points alone gives
\[
d\,\Tr(A^4)\;=\;\numLbTwoA+\frac{\numLbTwoB}{d},
\]
which reproduces the middle point, never used in the fit, to within $\numLbMidPull$ of its standard
error; the weighted three-point fit reads $a=\numLbFitA\pm\numLbFitAse$ and
$b=\numLbFitB\pm\numLbFitBse$ with a chi-squared of $\numLbChi$ on one degree of freedom. So
$\Tr(A^4)=\Theta(1/d)$, which is (L-b), with the constant identified rather than merely bounded.

The constant repays a second look, because $2$ is not an arbitrary number here. Were the $p-1$ rows
generic rather than Veronese, their normal would be a uniform direction on the Frobenius sphere of
$\mathrm{Sym}_d$; the spectrum of such a direction is semicircular \cite{Wigner,AGZ}, and
$d\,\Tr(A^4)$ then tends to the fourth moment of the semicircle law, the Catalan number $C_2=2$. At
finite $d$ the null is exact: a uniform direction is $U=G/\|G\|_F$ with $G$ a GOE matrix, whose
norm is independent of its direction, so
\begin{equation}
d\,\E\Tr(U^4)=\frac{d\,\E\Tr(G^4)}{\E\|G\|_F^4}=\frac{d\,(2d^3+5d^2+5d)}{4p(p+2)},
\qquad p=\frac{d(d+1)}{2},
\label{eq:null}
\end{equation}
using $\E\Tr(G^4)=2d^3+5d^2+5d$ for the GOE normalised to unit off-diagonal variance, by Wick's
formula, and $\|G\|_F^2\sim2\chi^2_p$. The measured intercept $\numLbFitA\pm\numLbFitAse$ sits $\numLbFitPull$
standard errors from $2$. To leading order the leave-one-out normal behaves like a generic direction,
and the limit in (L-b) is the semicircle's.

What is not generic is the approach to it. At each dimension the measured value exceeds the null
\eqref{eq:null} by forty to fifty standard errors (Table~\ref{tab:lb}), so the gap is not a sampling
artefact, and multiplying the excess by $d$ gives values that drift gently downward rather than
sitting still: a limit with its own $1/d$ correction, the two-point fit $\numLbExcA+\numLbExcB/d$
landing the middle value within $\numLbExcPull$ of a standard error. So the excess over a uniform
direction is itself $\Theta(1/d)$, and
\[
\Tr(A^4)\;=\;\frac{2}{d}+\frac{c}{d^2}+\cdots,\qquad c\approx \numLbExcC,
\]
where a uniform direction contributes under one to that second coefficient. The Veronese structure
of the rows is invisible in the leading term of $\Tr(A^4)$ and enters only at the next one. For
(L-b), which asks only for the order, that is a wide margin; for anyone setting out to prove it, the
leading calculation is one about a generic direction in $\mathrm{Sym}_d$, and the rows' structure
can be carried as a correction.

\begin{table}[t]
\centering\small
\caption{Lemma (L-b) at the three dimensions with the largest design counts. The null is the exact
generic-direction value \eqref{eq:null}; the excess is the measured mean minus the null, with the
standard error of the mean and the excess in units of it; the last column is $d$ times the excess.}
\label{tab:lb}
\begin{tabular}{rrlrlrr}
\toprule
$d$ & designs & $d\,\Tr(A^4)$ & null \eqref{eq:null} & excess & in s.e. & $d\times$excess\\
\midrule
% BEGIN AUTOTABLE tab:lb
20 & 1200 & $2.3728\pm0.0069$ & 2.0328 & $0.340\pm0.007$ & 49 & 6.80\\
30 & 900 & $2.2436\pm0.0049$ & 2.0256 & $0.218\pm0.005$ & 45 & 6.54\\
40 & 600 & $2.1813\pm0.0039$ & 2.0207 & $0.161\pm0.004$ & 41 & 6.43\\
% END AUTOTABLE tab:lb
\bottomrule
\end{tabular}
\end{table}

\subsection{Why the constant is two}\label{sec:semicircle}

The value the constant took is not a numerical accident, and following it gives a statement about
$A$ that is stronger than (L-b), simpler than (L-b), and --- unlike (L-b) --- of a shape that a proof
can attack directly. Two exact facts do the work, and both are elementary.

\begin{proposition}[Conjugation invariance]\label{prop:conj}
The law of $A$ is invariant under $A\mapsto OAO^\top$ for every $O\in O(d)$. Consequently $A=U\Lambda
U^\top$ with $U$ Haar on $O(d)$ independent of $\Lambda$, and every question about $A$ is a question
about its eigenvalue law alone.
\end{proposition}

\begin{proof}
For $v\in\R^p$ let $M_v\in\mathrm{Sym}_d$ be its un-vectorization, so that $z(x)^\top v=\langle
M_v,xx^\top-I\rangle/\sqrt2$. Then $z(Ox)^\top v=\langle M_v,Oxx^\top O^\top-I\rangle/\sqrt2
=\langle O^\top M_vO,xx^\top-I\rangle/\sqrt2$, so replacing every input $x_i$ by $Ox_i$ maps the row
space of $Z$ through the linear map $M\mapsto OMO^\top$, which is orthogonal for the Frobenius inner
product. The orthogonal complement of the row space maps the same way, and the inputs $Ox_i$ have
the law of the $x_i$.
\end{proof}

\begin{proposition}[The Veronese feature vector is isotropic]\label{prop:iso}
$\E\,z(x)z(x)^\top=I_p$ exactly, for every $d$.
\end{proposition}

\begin{proof}
The diagonal coordinates have $\E[(x_a^2-1)^2/2]=(3-1)/2=1$ and the off-diagonal ones
$\E[x_a^2x_b^2]=1$. Every cross term vanishes: $\E[(x_a^2-1)(x_b^2-1)]=0$ for $a\ne b$ by
independence, and any term carrying an odd power of some coordinate vanishes by symmetry.
\end{proof}

Both scalars in Proposition~\ref{prop:iso} are forced, and it is worth seeing why, because it says
what would have to fail for the leading order to stop being generic. The group $O(d)$ acts on
$\mathrm{Sym}_d$ preserving the splitting $\mathrm{Sym}_d=\R I\oplus\mathrm{Sym}_d^0$ into the scalars
and the traceless part, and for $d\ge3$ these are inequivalent irreducibles. Since
$z(Ox)=Oz(x)O^\top$ and $Ox$ has the law of $x$, the operator $\E\,zz^\top$ commutes with the action,
so by Schur's lemma it is scalar on each summand:
\[
\E\,zz^\top=\alpha\,P_{\R I}+\beta\,P_{\mathrm{Sym}_d^0} .
\]
No entry-by-entry computation can produce anything else, and isotropy is the single equality
$\alpha=\beta$. For the Gaussian row both are $1$: along the traceless block a pure off-diagonal
coordinate is $x_ax_b$ with variance $1$, and along $I/\sqrt d$ the row reads
$(\lvert x\rvert^2-d)/\sqrt{2d}$, whose variance is $1$ because $\chi^2_d$ has variance $2d$. So
$\alpha=\beta$ is a property of the Gaussian and not of the Veronese map, and it is the only thing
about the row law that the leading order uses. The same argument applied to the normal itself gives
its second moment: $\E\,A\otimes A$ is $a_dP_{\R I}+b_dP_{\mathrm{Sym}_d^0}$ with $a_d=O(1/d)$ and
$b_d=(1-a_d)/(p-1)$, so that $\E\langle A,B\rangle^2=(1+O(1/d))/(p-1)$ for every fixed unit
traceless $B$ \cite{PN24}; this is the second-moment shadow of the statement below, at the scale a
Haar-distributed direction would give.

Proposition~\ref{prop:iso} is what makes the leading order generic. Had the $p-1$ rows been
isotropic \emph{Gaussian} vectors in $\R^p$, their normal would have been exactly uniform on the unit
Frobenius sphere of $\mathrm{Sym}_d$; a uniform direction is a Frobenius-normalized GOE matrix, whose
spectrum is semicircular, and the semicircle's moments are the Catalan numbers. The rows are not
Gaussian --- they are quadratic in $x$ --- but they are exactly isotropic, and that is the part the
leading order sees.

Proposition~\ref{prop:conj} also says what cannot be measured. Every $\Tr(A^{2k})$, the operator
norm, the Kolmogorov--Smirnov distance to the semicircle --- every statistic of $A$ alone that is
invariant under conjugation --- is a function of the eigenvalues, and the eigenvectors carry no
information beyond being Haar. So no such statistic can distinguish the Veronese normal from any
other conjugation-invariant matrix with the same spectrum, and whatever the design leaves in $A$
beyond its spectrum can only be seen in a statistic that involves $A$ \emph{and} the data: the
alignment of an eigenvector with one of the $x_i$, or $x^\top Ax$ at a fresh point. That is where a
proof of the conjecture may not go, and where a sharper description of $A$ must.

So the statement the measurement is pointing at is this.

\begin{conjecture}\label{conj:semicircle}
$\sqrt d\,A$ converges in moments to the standard semicircle law on $[-2,2]$; equivalently, for every
$k\ge1$,
\[
d^{\,k-1}\,\E\Tr(A^{2k})\;\longrightarrow\;C_k=\frac1{k+1}\binom{2k}{k},
\]
the odd moments vanishing by symmetry. Lemma (L-b) is the case $k=2$, where $C_2=2$.
\end{conjecture}

Table~\ref{tab:moments} is the test proper --- the same normal taken to the sixth, eighth and tenth
moments and out to dimensions the full hard-edge driver cannot reach, the null vector being a linear
solve rather than a decomposition. Every moment sits above its Catalan number and falls toward it as
$d$ grows, and the edge $\sqrt d\,\|A\|_{\mathrm{op}}$ sits above $2$ throughout, between
$\numEdgeMin$ and $\numEdgeMax$ over $\numMomDmin\le d\le\numMomDmax$, with no growth. The pipeline
has a control (Table~\ref{tab:control}): replacing the Veronese rows by i.i.d.\ Gaussian or
Rademacher rows in $\R^p$ makes the null direction exactly uniform on the sphere, so its
un-vectorization is a Frobenius-normalized GOE matrix and its fourth moment must be the exact null
\eqref{eq:null}; it is, to within $\numControlWorstSe$ standard errors at every dimension run, while
the Veronese law
sits far above at the same dimensions. The whole of the Veronese excess is a property of the feature
law rather than of the measurement.

The limit is an extrapolation, and it is worth being exact about what the data support.
Table~\ref{tab:momentfit} reads off the intercept two ways over the whole range of dimensions,
weighting each dimension by the inverse squared standard error of its mean. A two-term fit $a+b/d$
sits below the Catalan number at every $k$, by under one standard error at $k=2$ and by between
two and four at $k=4$ and $k=5$, with a chi-squared that reaches $\numMomChiTwoLast$ on
$\numMomChiDof$ degrees of freedom at $k=5$: the higher moments carry a visible $1/d^2$ term that a
two-term fit cannot absorb and loads onto the intercept. Adding that term, the three-term fit
$a+b/d+e/d^2$ brings the intercept to within $\numMomFitWorst$ standard errors of the Catalan number
at each of the four moments, and its $1/d$ coefficient at $k=2$, $b=\numMomFitB\pm\numMomFitBse$,
agrees with the value the hard-edge driver measures independently in Table~\ref{tab:lb}. So the
reading is the careful one: the data are consistent with the Catalan limits and the limits are
approached, not attained --- at $d=\numMomDmax$, the largest reached, $d\,\Tr(A^4)$ is still
$\numMomLastMfour$, $\numMomLastPull$ standard errors above $2$. Above $d=100$ the ladder rests on
ten, six and six designs.

\begin{table}[t]
\centering\small
\caption{Conjecture~\ref{conj:semicircle} under test. $A$ is the same unit-Frobenius normal as in
Table~\ref{tab:t4}; each column reports $d^{\,k-1}\Tr(A^{2k})$ against the Catalan number it should
converge to, and the last column the edge against $2$. Entries are mean $\pm$ standard error over
independent designs. The null vector is obtained by a linear solve and checked by its residual
$\|Zv\|$, which is below $10^{-10}$ on every design reported here.}
\label{tab:moments}
{\setlength{\tabcolsep}{4.5pt}
\begin{tabular}{rrlllll}
\toprule
$d$ & designs & $d\,\Tr A^4$ & $d^2\,\Tr A^6$ & $d^3\,\Tr A^8$ & $d^4\,\Tr A^{10}$ &
$\sqrt d\,\|A\|_{\mathrm{op}}$\\
 & & ($C_2=2$) & ($C_3=5$) & ($C_4=14$) & ($C_5=42$) & ($2$)\\
\midrule
% BEGIN AUTOTABLE tab:moments
20 & 320 & $2.371\pm0.012$ & $7.458\pm0.098$ & $27.005\pm0.621$ & $106.621\pm3.671$ & $2.034\pm0.008$\\
30 & 270 & $2.258\pm0.008$ & $6.678\pm0.062$ & $22.665\pm0.364$ & $83.736\pm1.995$ & $2.036\pm0.007$\\
40 & 200 & $2.187\pm0.006$ & $6.198\pm0.047$ & $20.082\pm0.272$ & $70.724\pm1.485$ & $2.029\pm0.007$\\
60 & 140 & $2.118\pm0.005$ & $5.752\pm0.036$ & $17.784\pm0.204$ & $59.649\pm1.087$ & $2.021\pm0.007$\\
80 & 60 & $2.091\pm0.005$ & $5.565\pm0.029$ & $16.743\pm0.149$ & $54.266\pm0.709$ & $2.007\pm0.006$\\
100 & 40 & $2.071\pm0.004$ & $5.440\pm0.030$ & $16.151\pm0.166$ & $51.722\pm0.850$ & $2.009\pm0.010$\\
140 & 10 & $2.055\pm0.009$ & $5.332\pm0.056$ & $15.575\pm0.269$ & $48.934\pm1.212$ & $2.019\pm0.013$\\
160 & 6 & $2.054\pm0.008$ & $5.319\pm0.053$ & $15.500\pm0.265$ & $48.577\pm1.206$ & $2.009\pm0.017$\\
200 & 6 & $2.037\pm0.005$ & $5.227\pm0.029$ & $15.100\pm0.132$ & $46.912\pm0.572$ & $2.013\pm0.006$\\
% END AUTOTABLE tab:moments
\bottomrule
\end{tabular}}
\end{table}

\begin{table}[t]
\centering\small
\caption{The pipeline control. The same null-vector computation with the Veronese rows replaced by
i.i.d.\ Gaussian and by i.i.d.\ Rademacher rows in $\R^p$, whose normal is exactly a uniform
direction; its $d\,\Tr(A^4)$ must equal \eqref{eq:null}, which is $2.0329$, $2.0258$, $2.0207$ and
$2.0141$ at $d=20,30,40,60$. Mean $\pm$ standard error over the designs.}
\label{tab:control}
{\footnotesize\setlength{\tabcolsep}{3pt}
\begin{tabular}{rlrlllll}
\toprule
$d$ & rows & designs & $d\,\Tr A^4$ & $d^2\,\Tr A^6$ & $d^3\,\Tr A^8$ & $d^4\,\Tr A^{10}$ &
$\sqrt d\,\|A\|_{\mathrm{op}}$\\
\midrule
% BEGIN AUTOTABLE tab:control
20 & Veronese & 320 & $2.371\pm0.012$ & $7.458\pm0.098$ & $27.005\pm0.621$ & $106.621\pm3.671$ & $2.034\pm0.008$\\
20 & Gaussian rows & 200 & $2.041\pm0.011$ & $5.258\pm0.066$ & $15.311\pm0.334$ & $48.263\pm1.588$ & $1.861\pm0.008$\\
20 & Rademacher rows & 200 & $2.037\pm0.010$ & $5.263\pm0.065$ & $15.443\pm0.333$ & $49.207\pm1.600$ & $1.867\pm0.008$\\
30 & Veronese & 270 & $2.258\pm0.008$ & $6.678\pm0.062$ & $22.665\pm0.364$ & $83.736\pm1.995$ & $2.036\pm0.007$\\
30 & Gaussian rows & 150 & $2.039\pm0.008$ & $5.263\pm0.050$ & $15.384\pm0.254$ & $48.667\pm1.208$ & $1.916\pm0.008$\\
30 & Rademacher rows & 150 & $2.025\pm0.007$ & $5.186\pm0.045$ & $15.023\pm0.229$ & $47.080\pm1.106$ & $1.905\pm0.008$\\
40 & Veronese & 200 & $2.187\pm0.006$ & $6.198\pm0.047$ & $20.082\pm0.272$ & $70.724\pm1.485$ & $2.029\pm0.007$\\
40 & Gaussian rows & 100 & $2.018\pm0.006$ & $5.118\pm0.040$ & $14.616\pm0.208$ & $44.984\pm0.995$ & $1.916\pm0.008$\\
40 & Rademacher rows & 100 & $2.002\pm0.008$ & $5.046\pm0.051$ & $14.373\pm0.259$ & $44.323\pm1.229$ & $1.917\pm0.009$\\
60 & Veronese & 140 & $2.118\pm0.005$ & $5.752\pm0.036$ & $17.784\pm0.204$ & $59.649\pm1.087$ & $2.021\pm0.007$\\
60 & Gaussian rows & 60 & $2.022\pm0.006$ & $5.144\pm0.039$ & $14.725\pm0.184$ & $45.305\pm0.812$ & $1.944\pm0.007$\\
60 & Rademacher rows & 60 & $2.013\pm0.006$ & $5.077\pm0.034$ & $14.384\pm0.157$ & $43.792\pm0.700$ & $1.938\pm0.007$\\
% END AUTOTABLE tab:control
\bottomrule
\end{tabular}}
\end{table}

\begin{table}[t]
\centering\small
\caption{The intercept of the moment ladder of Table~\ref{tab:moments}, computed from the records
over every dimension measured, each weighted by the inverse squared standard error of its mean. For
each $k$: a two-term fit $a+b/d$, its intercept's distance from the Catalan number in standard
errors and its chi-squared over degrees of freedom; then the same for a three-term fit
$a+b/d+e/d^2$, with its $1/d$ coefficient.}
\label{tab:momentfit}
{\setlength{\tabcolsep}{4.5pt}
\begin{tabular}{rrlrrllrr}
\toprule
$k$ & $C_k$ & two-term $a$ & s.e.\ from $C_k$ & $\chi^2/\mathrm{dof}$ & three-term $a$ & $b$ &
s.e.\ from $C_k$ & $\chi^2/\mathrm{dof}$\\
\midrule
% BEGIN AUTOTABLE tab:momentfit
2 & 2 & $1.997\pm0.004$ & -0.7 & 4.4/7 & $2.000\pm0.006$ & $7.2\pm0.7$ & +0.0 & 4.1/6\\
3 & 5 & $4.959\pm0.024$ & -1.7 & 4.0/7 & $4.995\pm0.040$ & $44.8\pm4.7$ & -0.1 & 2.8/6\\
4 & 14 & $13.669\pm0.121$ & -2.7 & 5.8/7 & $13.970\pm0.201$ & $213.9\pm25.3$ & -0.2 & 2.3/6\\
5 & 42 & $39.904\pm0.581$ & -3.6 & 9.6/7 & $41.956\pm0.952$ & $917.7\pm127.7$ & -0.0 & 2.2/6\\
% END AUTOTABLE tab:momentfit
\bottomrule
\end{tabular}}
\end{table}

Moments are a confirmation, and four of them agreeing is still four numbers; a law that is not the
semicircle can match them. The observation that comes out differently if the conjecture is false is
the shape of the whole spectrum, and the statistic for that is the Kolmogorov--Smirnov distance
between the empirical distribution of $\sqrt d\,\lambda$ and the semicircle law itself. If the conjecture
holds this must fall to zero; if the limit is some other law it must flatten at a positive value,
however well the low moments behave. It falls (Table~\ref{tab:kssemi}), and what makes that readable
is the control beside it, since the absolute value of a Kolmogorov--Smirnov distance at $d$
eigenvalues means nothing on its own: Gaussian and Rademacher rows give a null direction that is
exactly uniform, hence genuinely semicircular in the limit, and at the two dimensions where both arms
were run the Veronese distance is $\numKsRatioMin$ to $\numKsRatioMax$ times the Gaussian-row
distance. A laptop-scale run of the same construction at $d=20$, with $200$ designs per arm, gives
$0.0919\pm0.0012$ against $0.0848\pm0.0012$, the same ratio. The Veronese spectrum sits a tenth
further from the semicircle than an exactly uniform direction does, at every dimension where the
comparison exists, and the two fall together. A law converging to something else would have separated
from that control rather than tracking it.

\begin{table}[t]
\centering\small
\caption{Kolmogorov--Smirnov distance between the empirical law of $\sqrt d\,\lambda$ and the
semicircle law on $[-2,2]$, mean $\pm$ standard error over the designs whose records carry the
statistic, by row law.}
\label{tab:kssemi}
\begin{tabular}{rlrl}
\toprule
$d$ & rows & designs & KS distance\\
\midrule
% BEGIN AUTOTABLE tab:kssemi
40 & Veronese & 100 & $0.0515\pm0.0009$\\
40 & Gaussian rows & 100 & $0.0484\pm0.0009$\\
40 & Rademacher rows & 100 & $0.0481\pm0.0008$\\
60 & Veronese & 100 & $0.0379\pm0.0006$\\
60 & Gaussian rows & 60 & $0.0346\pm0.0007$\\
60 & Rademacher rows & 60 & $0.0359\pm0.0007$\\
80 & Veronese & 30 & $0.0300\pm0.0009$\\
100 & Veronese & 20 & $0.0247\pm0.0009$\\
200 & Veronese & 6 & $0.0141\pm0.0005$\\
% END AUTOTABLE tab:kssemi
\bottomrule
\end{tabular}
\end{table}

Nor is the conjecture special to the Gaussian row. Take the inputs uniform on the unit sphere
instead and centre inside the traceless block, $z(x)=xx^\top-I/d$; then the design never leaves
$\mathrm{Sym}_d^0$, the scalar block of Proposition~\ref{prop:iso} is absent, isotropy is automatic
by Schur, and $A$ is the unit traceless symmetric matrix whose quadratic form kills the design,
$x_i^\top Ax_i=0$ for $i\le\dim\mathrm{Sym}_d^0-1$. In laptop-scale runs of that construction
the four moments converge to the same Catalan numbers, and the corrections are the same size: with
$7{,}373$, $2{,}650$ and $925$ designs at $d=24,32,40$ the excess over $C_k$ is $0.93$ to $1.03$
times the excess of the Gaussian-input normal at the same $d$, read off the $1/d+1/d^2$ model fitted
to the records, for every $k$ from $2$ to $5$; a direct pairing of the two constructions at the same
$d$ with a few hundred designs each gives ratios between $0.9$ and $1.2$ with standard errors of
five to ten percent. Same sign, same $1/d$ shape, the same size to within the errors. Two row laws,
two ambient spaces, one limit.

Where the non-Gaussianity sits is visible, and it has an exact formula. Write $U$ for a unit-Frobenius
symmetric matrix and $u_1,\dots,u_d$ for its eigenvalues. Then $\langle z,U\rangle=\sum_iu_i(g_i^2-1)/\sqrt2$
with $g$ standard Gaussian, a weighted sum of independent centred $\chi^2_1$ variables; since
$(g^2-1)$ has variance $2$ and fourth cumulant $48$,
\begin{equation}\label{eq:kurt}
\operatorname{kurt}\langle z,U\rangle\;=\;3+12\,\Tr(U^4).
\end{equation}
The two coordinate values are the extremes of \eqref{eq:kurt}: a pure diagonal coordinate has
$\Tr(U^4)=1$ and kurtosis $15$, a pure off-diagonal coordinate has $\Tr(U^4)=\tfrac12$ and kurtosis
$9$. Along a \emph{typical} direction $\Tr(U^4)$ is the fourth moment of a Frobenius-normalized GOE
spectrum, \eqref{eq:null} divided by $d$, so the kurtosis along a uniformly random direction is
$3+24/d$ to leading order; a random direction mixes $\Theta(d^2)$ coordinates and the mixture is
asymptotically Gaussian. Along the null direction of a design the row is slightly more leptokurtic
than along an independent direction: for the sphere construction, where the typical value falls just
below three, the pairs read $2.84$ against $3.21$ at $d=8$, $2.89$ against $3.23$ at $d=10$, $2.94$
against $3.18$ at $d=14$ and $2.97$ against $3.11$ at $d=20$, the gap closing with $d$, while the two
variances agree to three digits at every $d$ --- a direct check that the null direction is isotropic
at second order (laptop, $12{,}753$ to $97{,}056$ designs).

What the fourth moment does \emph{not} explain is the size of the correction, and the sphere
construction is the reason. Its row $xx^\top-I/d$ has no $\lvert x\rvert^2$ fluctuation, and its
kurtosis along a typical direction falls just \emph{below} three, against $3+24/d$ for the Gaussian
row. Two row laws whose kurtoses straddle the Gaussian value, one leptokurtic by a factor of nearly
two at $d=10$ and the other slightly platykurtic, produce corrections agreeing to within a tenth. A
fourth moment cannot do that. Matching directly on the kurtosis makes the point sharper still. Take
rows with \emph{independent} coordinates of variance one and kurtosis $\kappa$ --- the three-point
law $\Pr(\pm\kappa^{1/2})=1/(2\kappa)$, $\Pr(0)=1-1/\kappa$ realises every $\kappa\ge1$ --- and the
excess over $C_k$ is decided by $\kappa$ and nothing else: at $d=14$, over about $10^5$ designs per
arm, a three-point law at $\kappa=3$ and a Gaussian agree to within half a standard error at
$k=2,3,4,5$, while moving $\kappa$ from $1$ to $9$ multiplies the excess by $2.3$; at $d=24$ the same
comparison holds for three matched pairs (uniform against three-point at $\kappa=1.8$, Gaussian
against three-point at $\kappa=3$, Laplace against three-point at $\kappa=6$, all within two standard
errors over $\numKurtDtwentyfourNmin$ to $\numKurtDtwentyfourNmax$ designs per arm) and the same
move in $\kappa$ multiplies the
excess by only $1.15$, so the $\kappa$-dependence is itself a higher-order term in $1/d$ than the
excess. The Veronese row at $d=14$ has kurtosis $3+24/14\approx4.7$ along a typical direction and an
excess of $0.46$ at $k=2$ against $0.041$ for independent coordinates matched to that kurtosis, an
order of magnitude larger. With independent coordinates the correction is a fourth-moment functional;
the Veronese row does not have independent coordinates, it is rank one, and that is the structure the
two constructions share and the fourth moment does not see. A row $(XX^\top-rI)/\sqrt{2r}$ with $X$
a $d\times r$ Gaussian matrix is a sum of $r$ independent rank-one forms, has kurtosis
$3+12\,\Tr(U^4)/r$ along $U$, and its excess falls roughly as $1/r$ toward the generic value: $0.46$,
$0.23$, $0.13$, $0.085$ at $r=1,2,4,8$ ($k=2$, $d=14$). Within one dependence structure the kurtosis
tunes the correction; across structures at a fixed kurtosis it does not, and the Veronese row's $r=1$
is the extreme. A derivation of the constant $b$ has to come from the rank-one dependence, not from a
kurtosis. All of the numbers in this paragraph are laptop-scale computations at $d\le24$, recorded
with their scripts beside the run records.

\begin{remark}[(L-b) follows from a delocalization bound, and the loss is a factor two]
Since $\Tr(A^2)=1$,
\[
\Tr(A^4)=\sum_i\lambda_i^4\;\le\;\Big(\max_i\lambda_i^2\Big)\sum_i\lambda_i^2=\|A\|_{\mathrm{op}}^2 ,
\]
so (L-b) is implied by $\E\|A\|_{\mathrm{op}}^2=O(1/d)$ --- a statement about one eigenvalue rather
than about the whole spectrum, and the kind of delocalization estimate the random-matrix literature
is built to supply. It is worth knowing how much that costs, because a crude bound that loses an
unbounded factor is not a route. It loses exactly two. Under Conjecture~\ref{conj:semicircle} the
edge is $2/\sqrt d$ and the fourth moment $2/d$, so the bound gives $4/d$ against a truth of $2/d$;
measured, the ratio of $\|A\|_{\mathrm{op}}^2$ to $\Tr(A^4)$ climbs from $\numOpRatioFirst$ at
$d=\numMomDmin$ to $\numOpRatioLast$ at $d=\numMomDmax$, converging to the two that the semicircle
predicts. So an operator-norm route to (L-b) is not lossy in order and is lossy by a factor of two in
the constant, which the reader attacking it should know before choosing between it and the moment
route.
\end{remark}

\begin{remark}[The Gram matrix as a sum of independent rank-one terms]
The Gram matrix of the reduction, $\sum_{i}z(X_i)z(X_i)^\top$, is a sum of independent rank-one
matrices with heavily dependent entries within each summand, which is the setting of the matrix
concentration inequalities of Brailovskaya and van Handel \cite{BvH} and of the Gaussian theory they
rest on \cite{BBvH}, and outside that of edge universality for matrices with independent entries. In
that framework the spectrum is controlled by the free model with the same mean and covariance --- here
the Marchenko--Pastur law at aspect ratio one \cite{MarchenkoPastur}, whose hard edge is the one the
conjecture names --- provided the norm of a single summand is small against the variance parameter
of the sum. The variance parameter is exact: by the same Schur decomposition, $\E[\|z\|^2zz^\top]$
has the two eigenvalues $(d^2+13d+16)/2$ on $\R I$ and $(d^2+9d+16)/2$ on $\mathrm{Sym}_d^0$, so for
$X=\sum_{i\le p-1}(z_iz_i^\top-I)$ one has $\sigma(X)^2=\|\E X^2\|=(p-1)(d^2+13d+14)/2$. The summand
norm is $\|z(x)\|^2=(|x|^4-2|x|^2+d)/2$ with $|x|^2\sim\chi^2_d$, of mean $p$ and with $|x|^4/2$ as
its leading term, and in a laptop-scale sample at $d=\numVhDmin$ to $\numVhDmax$ its $99$th
percentile runs from $\numVhRqFirst$ down to $\numVhRqLast$ times $\sigma(X)$ while its maximum
stays between $\numVhRmaxMin$ and $\numVhRmaxMax$ times $\sigma(X)$: the boundedness parameter is of
the order of the variance parameter at every reachable $d$, not small against it. Its tail is exactly what makes the inverse
smallest singular value fail to be integrable at fixed $d$ (the remark after Table~\ref{tab:ks}); a
truncation of that tail is the natural device, and the rate at which the truncated summand becomes
negligible is the quantitative question the framework poses.
\end{remark}

\subsection{Does the reduced matrix look like a Wishart at the edge?}

The reduction is exact and the two lemmas look true, so the question that remains is whether
$[b,H]$ actually obeys the target law. That is directly testable: Table~\ref{tab:ks} compares the
empirical distribution of each statistic against $1-e^{-\sqrt t-t/2}$ by a Kolmogorov--Smirnov
distance, with the $95\%$ critical value $1.36/\sqrt N$ at $N$ designs beside it. After the
reduction the matrix is statistically indistinguishable from the square real Wishart hard edge at
every dimension from $d=20$ on, while $Z$ itself is rejected at $d=20$, $30$ and $40$ and becomes
consistent only from $d=60$, where the critical value has also grown. The test has finite
resolution and it should be quoted with it: at $1{,}200$ designs a Kolmogorov--Smirnov test sees a
discrepancy only above about $0.04$ in the cumulative distribution, so this is evidence against
distortions larger than four percentage points at $d=20$, and against progressively larger ones at
the higher dimensions with fewer designs. The two statistics differ by exactly the diagonal scaling
$D$, whose effect the theory shows vanishes and which at $d=20$ still contributes a factor
$s_{\min}(Z)/s_{\min}([b,H])=\numSminRatioFirst$; squaring that moves the whole distribution of
$p\,s_{\min}(Z)^2$ left by a third, which is what the first column registers. The finite-size
discrepancy therefore sits in the one place the proof already controls.

\begin{table}[t]
\centering\small
\caption{Kolmogorov--Smirnov distance of the empirical law of $p\,s_{\min}^2$ to the square real
Wishart hard edge $1-e^{-\sqrt t-t/2}$, for $Z$, for the reduced matrix $[b,H]$, and for the
$p\times(p-1)$ block $H$ alone (tested against the same law, to which it is not expected to
conform; see the text). The critical value is $1.36/\sqrt N$ at the $95\%$ level; the verdict is for
$Z$ and for $[b,H]$.}
\label{tab:ks}
\begin{tabular}{rrrrrrl}
\toprule
$d$ & designs & $Z$ & $[b,H]$ & $H$ & critical & verdict ($Z$ / $[b,H]$)\\
\midrule
% BEGIN AUTOTABLE tab:ks
20 & 1200 & 0.104 & 0.024 & 0.397 & 0.039 & rejected / consistent\\
30 & 900 & 0.078 & 0.024 & 0.401 & 0.045 & rejected / consistent\\
40 & 600 & 0.058 & 0.020 & 0.399 & 0.056 & rejected / consistent\\
60 & 420 & 0.054 & 0.035 & 0.431 & 0.066 & consistent / consistent\\
80 & 270 & 0.058 & 0.037 & 0.421 & 0.083 & consistent / consistent\\
100 & 240 & 0.035 & 0.047 & 0.427 & 0.088 & consistent / consistent\\
120 & 36 & 0.114 & 0.099 & 0.379 & 0.227 & consistent / consistent\\
% END AUTOTABLE tab:ks
\bottomrule
\end{tabular}
\end{table}

\begin{remark}
The same test applied to $H$ alone rejects at a distance near $0.4$ at every dimension, and that
number should not be read as evidence about anything. $H$ is $p\times(p-1)$, one column short of
square, and the hard-edge law depends on the aspect deficit through the Bessel parameter
$\alpha=|\text{rows}-\text{cols}|$ \cite{TracyWidom}; $1-e^{-\sqrt t-t/2}$ is the $\alpha=0$ law.
Testing an $\alpha=1$ matrix against it rejects by construction. The column is recorded only to say
that it is uninformative.
\end{remark}

\begin{remark}
For every fixed $d\ge2$, $\E\,s_{\min}(Z)^{-1}=\infty$ \cite{PN24}. The rows span $\mathrm{Sym}_d$
almost surely, so $\det Z\ne0$ almost surely, and the least singular value of a square matrix with
i.i.d.\ rows has a positive linear lower-tail coefficient at the origin in every fixed dimension
\cite{GregorattiMaran}, which forces the inverse first moment to diverge. This is not evidence
against the conjecture --- the Gaussian square ensemble diverges the same way \cite{Edelman}, and the
conjectured law $1-e^{-\sqrt t-t/2}$ has exactly the linear lower tail in $s_{\min}$ that causes it
--- but it forbids any argument that asks for uniform integrability of the full inverse smallest
singular value, as opposed to a truncated one or one with finitely many bottom modes removed.

It is visible in the records, and it is a warning about reporting. Table~\ref{tab:invmoment} splits
the designs at each dimension into disjoint blocks of $150$ and computes $(\sqrt p\,s_{\min}(Z))^{-1}$
in each. A sample mean of $150$ independent draws of a quantity with finite variance varies by a few
percent from block to block; these block means vary by factors of $\numInvMeanSpreadMin$ to
$\numInvMeanSpreadMax$, with coefficients of variation between $\numInvMeanCvMin$ and
$\numInvMeanCvMax$, because each is estimating the smallest draw in its block rather than any
population quantity, while the block medians spread by $\numInvMedSpreadMin$ to
$\numInvMedSpreadMax$ percent (coefficient of variation $\numInvMedCvMin$ to $\numInvMedCvMax$). That
is what a
divergent first moment looks like from inside a finite sample, and it is why the tables in this
section report medians and standard deviations rather than means of inverse quantities.
\end{remark}

\begin{table}[t]
\centering\small
\caption{Block means and block medians of $(\sqrt p\,s_{\min}(Z))^{-1}$ over disjoint blocks of
$150$ designs, with the coefficient of variation across blocks.}
\label{tab:invmoment}
\begin{tabular}{rrlrlr}
\toprule
$d$ & blocks & block means & c.v. & block medians & c.v.\\
\midrule
% BEGIN AUTOTABLE tab:invmoment
20 & 8 & 5.5--33.0 & 0.82 & 2.12--2.66 & 0.07\\
30 & 6 & 5.8--38.0 & 0.80 & 1.92--2.24 & 0.06\\
40 & 4 & 5.3--10.2 & 0.35 & 1.95--2.15 & 0.04\\
60 & 2 & 3.7--14.7 & 0.84 & 1.94--2.07 & 0.05\\
% END AUTOTABLE tab:invmoment
\bottomrule
\end{tabular}
\end{table}

\begin{remark}
The lower tail is where the conjecture lives, so it is worth looking at directly rather than through
a single Kolmogorov--Smirnov number. Table~\ref{tab:tail} takes the $2$nd, $5$th, $10$th and $20$th
percentiles of $p\,s_{\min}([b,H])^2$ and reports the mass the conjectured law assigns below each of
them; the nominal values are $0.02$, $0.05$, $0.10$, $0.20$. The law is a little heavy in the extreme
tail at $d=20$ and a little light at $d=30$, with no monotone drift in $d$ --- a pattern consistent
with sampling noise at these design counts rather than with a systematic distortion, and consistent
with the distances of Table~\ref{tab:ks}.
\end{remark}

\begin{table}[t]
\centering\small
\caption{The lower tail of $p\,s_{\min}([b,H])^2$: the mass assigned by $1-e^{-\sqrt t-t/2}$ below
the empirical $2$nd, $5$th, $10$th and $20$th percentiles, at every dimension with at least $200$
designs.}
\label{tab:tail}
\begin{tabular}{rrrrrr}
\toprule
$d$ & designs & at p2 & at p5 & at p10 & at p20\\
\midrule
% BEGIN AUTOTABLE tab:tail
20 & 1200 & 0.025 & 0.058 & 0.102 & 0.205\\
30 & 900 & 0.022 & 0.043 & 0.088 & 0.181\\
40 & 600 & 0.021 & 0.044 & 0.111 & 0.201\\
60 & 420 & 0.025 & 0.067 & 0.102 & 0.194\\
80 & 270 & 0.021 & 0.051 & 0.088 & 0.187\\
100 & 240 & 0.025 & 0.042 & 0.127 & 0.227\\
% END AUTOTABLE tab:tail
\bottomrule
\end{tabular}
\end{table}

\begin{remark}
A measurement cannot prove either lemma, and neither can it prove the hard edge. What it can do is
tell the next attempt where the difficulty is likely to be. On this evidence neither lemma is the
obstruction, and the reduced matrix already passes the test the conjecture is about at the smallest
dimension measured; what is missing is a hard-edge theorem for the harmonic block, not a reason to
doubt the target.
\end{remark}

\section{Target 1: the surrogate's residuals, and why no Stein comparison is needed}\label{sec:t1}

\subsection{A Hermite identification}

The proof of Theorem~\ref{thm:T1} matches a Gaussian surrogate whose same-view variance and cross-view
covariance agree with the activation's, with the shared and private residual coefficients of
Dhifallah and Lu \cite{DhifallahLu}. The identification of those coefficients is a Hermite
expansion, and their nonnegativity is immediate from it.

\begin{proposition}\label{prop:hermite}
Let $s=\sqrt{1+\Delta^2}$ and $\rho=1/(1+\Delta^2)$, and expand $\varphi(\alpha)=\sigma(s\alpha)
=\sum_{k\ge0}c_k\mathrm{He}_k(\alpha)/\sqrt{k!}$ for $\alpha\sim N(0,1)$. Write $V=\E\varphi^2$ and
$C_{12}=\E[\varphi(A_1)\varphi(A_2)]$ for a standard bivariate normal pair of correlation $\rho$.
Then the shared and private residual variances of the matched surrogate are
\[
c_{\mathrm{sh}}=C_{12}-c_1^2\rho=\sum_{k\ge2}c_k^2\rho^k,
\qquad
c_{\mathrm{pr}}=V-c_1^2-c_{\mathrm{sh}}=\sum_{k\ge2}c_k^2\,(1-\rho^k),
\]
and both are nonnegative for every $\rho\in[0,1]$.
\end{proposition}

\begin{proof}
Orthonormality of the Hermite basis gives $V=\sum_k c_k^2$ and, since a bivariate normal pair of
correlation $\rho$ satisfies $\E[\mathrm{He}_j(A_1)\mathrm{He}_k(A_2)]=\delta_{jk}k!\rho^k$, also
$C_{12}=\sum_k c_k^2\rho^k$. Subtracting the $k=1$ terms gives the two displayed series, and
$c_0=0$ for odd $\sigma$. Each series has nonnegative terms for $0\le\rho\le1$.
\end{proof}

Numerically, with $\sigma=\tanh$ and $\Delta=1/2$ (so $\rho=0.8$): $c_1=\numHermCone$,
$V=\numHermV$, $C_{12}=\numHermCtwelve$, $c_{\mathrm{sh}}=\numHermSh$, $c_{\mathrm{pr}}=\numHermPr$,
and the two definitions of each agree to machine precision.

\subsection{The leverage matrix}

The reason the proof of Theorem~\ref{thm:T1} does not need a three-dimensional correlated Stein comparison is that
the exact leave-one-block objective increment depends on six scalars, not three --- the teacher, the
two logits, and the three entries of the $2\times2$ leverage matrix
$Q_{ab}=N^{-1}x_a^\top H_-^{-1}x_b$ --- and $Q$ concentrates. The Gaussian Poincar\'e inequality
\cite{BLM} gives $\Var(Q_{ab}\mid F,H_-)\le C/d$, after which the block functional is quadratic in
the three logits and second-moment matching suffices \cite{PN24}; this is the step at which the
argument departs from the Stein-method route of \cite{HuLu}.

That is the step measured here. Table~\ref{tab:t1} reports $d\,\Var(q_{11})$ over fresh blocks,
together with the relative gap between the optimized objectives of the actual and surrogate models,
the gap between their clean risks, and the gap between the two models' risks at the same $w$, at
$d=64$ for three augmentation strengths and at $d=128$ for one, and for two normalizations of the
feature matrix: columns of $F$ normalized to unit length, and raw Gaussian columns, whose
$O(\sqrt{\log d/d})$ norm fluctuation is one of the two sources the rate $\varepsilon_d$ is made of.
The product $d\,\Var(q_{11})$ is of order one and depends on $\Delta$, falling from $0.49$ at
$\Delta=1/4$ to $0.29$ at $\Delta=1$; the two column normalizations agree to within a hundredth on
it, on the objective gap and on the end-to-end risk gap, and differ only in the same-$w$ risk gap,
where the raw columns are $\numSwRatioMinSixtyFour$ to $\numSwRatioMaxSixtyFour$ times larger at
$d=64$ and $\numSwRatioOneTwoEight$ times at $d=128$; and at $\Delta=1/4$ the value at $d=128$ equals
the value at $d=64$. One doubling of $d$ is consistent with a constant and cannot distinguish it from
a slow drift, and nothing more is claimed for the behaviour in $d$.

\begin{table}[t]
\centering\small
\caption{Target 1 diagnostics, $\sigma=\tanh$, $m=2d$, $n=3d$, $\lambda=10^{-2}$; ``unit'' rows
normalize the columns of $F$, ``raw'' rows leave them Gaussian. Entries are mean $\pm$ standard
deviation over the designs; the two risk columns are the gap between the two models' clean risks at
their own minimizers and at the same $w$. A fourth column constant in $d$ is the $\Theta(1/d)$
leverage variance the proof needs.}
\label{tab:t1}
{\setlength{\tabcolsep}{4.5pt}
\begin{tabular}{rrlllrrr}
\toprule
$d$ & $\Delta$ & columns & $d\,\Var(q_{11})$ & $|L^*_{\mathrm{act}}-L^*_{\mathrm{sur}}|/L^*$ &
$|R_{\mathrm{act}}-R_{\mathrm{sur}}|$ & same-$w$ & designs\\
\midrule
% BEGIN AUTOTABLE tab:t1
64 & 0.25 & unit & $0.49\pm0.04$ & $0.07\pm0.05$ & $1.3\times10^{-2}$ & $1.4\times10^{-3}$ & 360\\
64 & 0.25 & raw & $0.49\pm0.04$ & $0.08\pm0.06$ & $1.4\times10^{-2}$ & $3.8\times10^{-3}$ & 360\\
64 & 0.5 & unit & $0.37\pm0.03$ & $0.06\pm0.05$ & $1.9\times10^{-2}$ & $1.8\times10^{-3}$ & 360\\
64 & 0.5 & raw & $0.37\pm0.03$ & $0.06\pm0.05$ & $2.0\times10^{-2}$ & $4.8\times10^{-3}$ & 360\\
64 & 1 & unit & $0.29\pm0.02$ & $0.05\pm0.04$ & $3.0\times10^{-2}$ & $2.6\times10^{-3}$ & 360\\
64 & 1 & raw & $0.29\pm0.02$ & $0.05\pm0.04$ & $3.0\times10^{-2}$ & $6.5\times10^{-3}$ & 360\\
128 & 0.25 & unit & $0.49\pm0.04$ & $0.06\pm0.04$ & $8.5\times10^{-3}$ & $1.3\times10^{-3}$ & 300\\
128 & 0.25 & raw & $0.49\pm0.04$ & $0.05\pm0.05$ & $7.7\times10^{-3}$ & $2.4\times10^{-3}$ & 100\\
% END AUTOTABLE tab:t1
\bottomrule
\end{tabular}}
\end{table}

\section{An exact same-sample identity, on two emulation corpora}\label{sec:corr}

\subsection{The identity}

A residual kernel correction fitted on top of a data-dependent mean is a construction used in
emulation pipelines for radiative transfer \cite{Shmalo26a,Shmalo26b}, following the kernel-flow
emulators of \cite{Lamminpaa} built on the kernel-flow method of \cite{OwhadiYoo}, and the question of when it helps is a question about the same sample
being used twice. Fix a realized training sample $S$; let $m_S$ and $k_S$ be a mean and a kernel that
may depend on $S$ in any way whatsoever, $r_S=G-m_S$ the residual function, $X$ the training inputs,
$A=K+n\lambda I$, $R=r_S(X)$, $\alpha=A^{-1}R$ and $q(u)=k_S(X,u)$. With $U\sim\mu$ independent of
$S$, put $Q_S=\E_U q(U)q(U)^\top$ and $B_S=\E_U q(U)r_S(U)^\top$. Then \cite{PN24}
\begin{equation}
\mathcal R_1(S)-\mathcal R_0(S)\;=\;-2\Tr(\alpha^\top B_S)+\Tr(\alpha^\top Q_S\alpha),
\label{eq:identity}
\end{equation}
exactly, with no independence assumed anywhere; the correction improves the conditional population
risk if and only if $2\Tr(\alpha^\top B_S)>\Tr(\alpha^\top Q_S\alpha)$. A two-point counterexample
in the same source shows that a same-sample residual correction can strictly worsen population risk
over a whole range of positive ridge, so the sign in \eqref{eq:identity} is not automatic. A
companion bound $\eta\le\kappa\,\delta_{\mathrm{rms}}/(2\sqrt\lambda)$ controls the $L^2$ distance
$\eta$ between the correction fitted against out-of-fold training-row predictions and the one fitted
against the deployed predictions, in terms of the root-mean-square difference $\delta_{\mathrm{rms}}$
between the two sets of training-row residuals.

\subsection{On the data}

The identity is evaluated on the two corpora of the emulation pipelines \cite{Shmalo26a}: the four
radiance components $Y_1,\dots,Y_4$ of the EMIT imaging spectrometer \cite{Green} and the three
spectral bands of the OCO-2 instrument \cite{Eldering,Susiluoto} --- the O$_2$ A-band and the weak and strong
CO$_2$ bands. In each case a neural mean (width $384$, depth $3$, $120$ epochs) is trained and two
residual vectors are formed on the training rows: the \emph{consistent} residuals, against the
deployed mean's own predictions, and the \emph{out-of-fold} residuals, against predictions from four
folds of the same architecture retrained without the row, which are the labels the emulation
pipeline fits its correction to. A Mat\'ern-$5/2$ correction is fitted to each, with one length
scale for both, chosen on the validation split for the correction fitted to the consistent
residuals, and both are scored on the test split, which supplies the sample from $\mu$. The kernel stage is fitted on a random subsample of
$8{,}000$ of the mean's own training rows: the identity is exact for whatever realized sample it is
evaluated on, and the cap is what keeps an exact Gram inside memory. The identity, the margin and
every ``gain'' below refer to the consistent arm; the out-of-fold arm is reported beside it in
Table~\ref{tab:corr} and is the subject of the mismatch bound. Neither $Q_S$ nor $B_S$ is ever
formed: with $K_{Xt}$ the train-by-test kernel block,
$\Tr(\alpha^\top Q_S\alpha)=\|K_{Xt}^\top\alpha\|_F^2/n_{\mathrm{te}}$ and
$\Tr(\alpha^\top B_S)=\langle K_{Xt}^\top\alpha,R_{\mathrm{te}}\rangle/n_{\mathrm{te}}$. Eight ridge
values from $10^{-8}$ to $10^{-1}$ are swept on every lane; there are five seeds per component,
except for the strong CO$_2$ band, which has one. Two corpora are two corpora: the remaining
emulation data sets of the same pipelines were not run, and nothing here generalises beyond EMIT and
OCO-2.

The identity holds to machine precision: across the $\numCorrLanes$ lanes and $\numCorrPairs$
(lane, ridge) pairs the largest relative discrepancy between the two sides of \eqref{eq:identity} is
$\numIdentWorst$ (Table~\ref{tab:corr}). That is a check on the arithmetic, and it passes.

The sign is where the content is, and it is decided by the component rather than by the ridge alone
(Tables~\ref{tab:corr} and~\ref{tab:signs}). On $Y_1$ and $Y_2$ the margin
$2\Tr(\alpha^\top B)-\Tr(\alpha^\top Q\alpha)$ is positive at every one of the eight ridges on every
seed, and the correction reduces the mean's conditional test risk by $\numGainemitYOne$ and
$\numGainemitYTwo$ percent at the best ridge. On $Y_3$ the margin is positive at every ridge from
$10^{-5}$ up, but it is negative at $10^{-8}$ on $\numYthreeNegA$ seeds, at $10^{-7}$ on
$\numYthreeNegB$ and at $10^{-6}$ on $\numYthreeNegC$ (Table~\ref{tab:signs}),
$\numHurtemitYThree$ of the forty pairs in all, with a crossing that is monotone in $\lambda$ and
interior to the swept grid; $Y_3$ is also the component the correction helps least when it does help,
$\numGainemitYThree$ percent. On $Y_4$ the margin is negative at every ridge on four of the five
seeds --- $\numHurtemitYFour$ of forty pairs --- and the correction changes the risk by
$\numGainemitYFour$ percent at the best ridge, which is to say by nothing. The mean's own test risk
on that component is an order of magnitude larger than on the other three, at the same rows, epochs
and folds: the network fits $Y_4$ poorly, and the correction spends its capacity on the part of the
residual that does not transfer to fresh points, which is the mechanism of the counterexample on the
production pipeline itself. On the OCO-2 bands the correction is close to neutral: gains of
$\numGainocoTwooTwo$, $\numGainocoTwoscoTwo$ and $\numGainocoTwowcoTwo$ percent on the O$_2$, strong
CO$_2$ and weak CO$_2$ bands at the best ridge, with the margin negative at every ridge up to
$10^{-6}$ on every lane, negative at $10^{-5}$ on $\numOcoNegFive$ of the $\numOcoLanes$ lanes,
and of either sign and within $\numOcoMarginMaxFour$ percent of $\mathcal R_0$ from $10^{-4}$ up.
So the harmful regime that the counterexample
describes is real and reachable on real data, and the identity locates it per component and per
ridge, computed rather than assumed.

\begin{table}[t]
\centering\small
\caption{The same-sample identity on the EMIT components and the OCO-2 bands. ``Gain'' is the
reduction in conditional test risk of the consistent arm --- the correction fitted to the residuals
against the deployed mean's own training-row predictions --- at its best ridge $\lambda^*$, as a
percentage of the mean's own risk, mean $\pm$ standard deviation over the lanes; ``phys.'' is the
same gain in the pipeline's physical relative error; ``hurt'' counts the (lane, ridge) pairs, out of
eight per lane, where the margin $2\Tr(\alpha^\top B)-\Tr(\alpha^\top Q\alpha)$ is negative. The two
``o.o.f.'' columns are the out-of-fold arm --- the correction fitted to the residuals against
out-of-fold predictions, the labels the pipeline fits to --- scored on the same test block: its gain
at the consistent arm's $\lambda^*$, and its best gain over the eight ridges. The last column is the
worst relative discrepancy between the two sides of \eqref{eq:identity}.}
\label{tab:corr}
{\small\setlength{\tabcolsep}{4pt}
\begin{tabular}{lrllrllr}
\toprule
corpus & lanes & gain \% & phys.\ gain \% & hurt & o.o.f.\ at $\lambda^*$ \% & o.o.f.\ best \% &
identity\\
\midrule
% BEGIN AUTOTABLE tab:corr
emit Y1 & 5 & $25.5\pm4.2$ & $25.9\pm3.0$ & 0 / 40 & $-45.7\pm17.6$ & $0.4\pm0.3$ & $2\times10^{-13}$\\
emit Y2 & 5 & $22.1\pm2.1$ & $27.1\pm5.1$ & 0 / 40 & $-41.3\pm20.4$ & $0.0\pm0.1$ & $1\times10^{-13}$\\
emit Y3 & 5 & $15.3\pm3.9$ & $22.2\pm5.3$ & 6 / 40 & $-34.7\pm21.4$ & $0.4\pm0.4$ & $2\times10^{-13}$\\
emit Y4 & 5 & $0.5\pm3.0$ & $0.5\pm1.5$ & 32 / 40 & $-2.1\pm3.1$ & $0.0\pm1.8$ & $1\times10^{-13}$\\
oco2 o2 & 5 & $0.3\pm0.1$ & $0.2\pm0.1$ & 31 / 40 & $-5.5\pm1.6$ & $-0.2\pm0.2$ & $4\times10^{-12}$\\
oco2 sco2 & 1 & $1.3\pm0.0$ & $0.6\pm0.0$ & 4 / 8 & $-15.1\pm0.0$ & $-0.2\pm0.0$ & $8\times10^{-14}$\\
oco2 wco2 & 5 & $0.3\pm0.3$ & $0.2\pm0.2$ & 28 / 40 & $-3.4\pm4.3$ & $0.1\pm0.1$ & $7\times10^{-13}$\\
% END AUTOTABLE tab:corr
\bottomrule
\end{tabular}}
\end{table}

\begin{table}[t]
\centering\small
\caption{Lanes on which the margin $2\Tr(\alpha^\top B)-\Tr(\alpha^\top Q\alpha)$ is positive, by
ridge value, out of the lanes in the second column.}
\label{tab:signs}
\begin{tabular}{lrrrrrrrrr}
\toprule
corpus & lanes & $10^{-8}$ & $10^{-7}$ & $10^{-6}$ & $10^{-5}$ & $10^{-4}$ & $10^{-3}$ &
$10^{-2}$ & $10^{-1}$\\
\midrule
% BEGIN AUTOTABLE tab:signs
emit Y1 & 5 & 5 & 5 & 5 & 5 & 5 & 5 & 5 & 5\\
emit Y2 & 5 & 5 & 5 & 5 & 5 & 5 & 5 & 5 & 5\\
emit Y3 & 5 & 2 & 3 & 4 & 5 & 5 & 5 & 5 & 5\\
emit Y4 & 5 & 1 & 1 & 1 & 1 & 1 & 1 & 1 & 1\\
oco2 o2 & 5 & 0 & 0 & 0 & 1 & 1 & 4 & 1 & 2\\
oco2 sco2 & 1 & 0 & 0 & 0 & 0 & 1 & 1 & 1 & 1\\
oco2 wco2 & 5 & 0 & 0 & 0 & 0 & 3 & 3 & 3 & 3\\
% END AUTOTABLE tab:signs
\bottomrule
\end{tabular}
\end{table}

One transfer is worth checking rather than assuming. The identity is a statement about squared
$L^2$ risk, while the emulator pipelines report a physical relative error, and an improvement in the
first need not survive into the second --- a gain concentrated on high-magnitude components can wash
out entirely once errors are taken relative. It survives here. Over the $\numEmitGainLanes$ EMIT
lanes on $Y_1$, $Y_2$ and $Y_3$ the squared-risk gain averages $\numEmitSqGain$ percent and the gain
in the physical metric $\numEmitPhysGain$ percent, and across all $\numSignTot$ lanes of
Table~\ref{tab:corr} the two gains agree in sign on $\numSignAgree$. The physical metric inverts the
principal-component and standardisation maps and returns the mean relative $L^2$ error in physical
units, the pipeline's own reported quantity. On these two corpora the quantity the theory governs
and the reported quantity move together.

The mismatch bound holds and is loose: at the best ridge of every lane, $\eta$ measures between
$\numEtaMin$ and $\numEtaMax$ against a bound that exceeds it by a factor between $\numSlackMin$ and
$\numSlackMax$. It certifies nothing useful at this operating point. The quantity that does the
certifying is the margin in \eqref{eq:identity}, computed directly, per corpus and per seed --- and
it certifies the consistent arm. The pipeline's own arm is a different estimator, and the bound's
transfer condition $\gamma>\eta$, with $\gamma=\sqrt{\mathcal R_0}-\sqrt{\mathcal R_1}$ the
consistent arm's root-risk margin, is what would carry the certificate across. It holds on
$\numTransferPairs$ of the $\numCorrPairs$ (lane, ridge) pairs and on $\numTransferLanes$ of the
$\numCorrLanes$ lanes, which is to say nowhere: wherever the consistent arm has a positive margin at
all, $\gamma$ is at most $\numGammaMax$ and $\eta$ exceeds it by a factor of at least
$\numTransferSlackMin$. The out-of-fold
arm behaves accordingly (Table~\ref{tab:corr}). At the ridge
where the consistent arm is best it raises the test risk on every EMIT lane, by
$\numOofAtEmitRiseMin$ to $\numOofAtEmitRiseMax$ percent of $\mathcal R_0$ on $Y_1$, $Y_2$ and
$Y_3$, and on $\numOofAtOcoNeg$ of the $\numOofAtOcoLanes$ OCO-2 lanes; its best gain over the eight
ridges is at most $\numOofBestEmitMax$ percent on $Y_1$--$Y_3$, $\numOofBestYfourMax$ percent on one
$Y_4$ seed and negative on the other $\numOofBestYfourNeg$, and between $\numOofBestOcoMin$ and
$\numOofBestOcoMax$ percent on OCO-2. Its physical error at the consistent arm's best ridge is worse
than the mean's on $\numOofPhysNeg$ of the $\numCorrLanes$ lanes, by up to $\numOofPhysWorseMax$
percent of the mean's own error, and the sign of that change agrees with the sign of its squared-risk
change on all $\numCorrLanes$ lanes. So at this
operating point --- an $8{,}000$-row kernel stage, four folds, the length scale chosen on the
consistent labels --- the correction the pipeline fits does nothing on these lanes, while the
consistent arm on the same rows gains $\numGainemitYThree$ to $\numGainemitYOne$ percent on the
components where its margin is positive. The identity locates the gain; the mismatch is what stands
between it and the deployed estimator, and neither the bound nor the data close that gap here.

\begin{remark}
The identity assumes nothing; the leave-one-out gate built on top of it in Section~\ref{sec:reg}
carries the hypotheses. That gate can be stated under uniform sup-norm deletion stability of the mean
and of the correction, with bounded-difference constants, or under the weaker integrated $L^2$
stability by way of the Efron--Stein inequality \cite{EfronStein}; the second is the form used here.
The fifth problem of Section~\ref{sec:prune} needs the same same-sample machinery for a quite
different reason.
\end{remark}

\section{The fifth problem: growing-rank spikes and what pruning may assume}\label{sec:prune}

The source is a manuscript on growing-rank spiked singular values and vectors with applications to
pruning deep networks, and the theorem about it established in \cite{PN24}. The model is
$W_n=X_n+P_n$ with $P_n=U_n\Theta_nV_n^*$ of rank $r_n\to\infty$, $r_n=o(n)$, and $n/m\to c\in(0,1]$,
a growing-rank form of the spiked models of \cite{BBP,BGN}. Everything the source proves about
cluster projectors and mixed dyads --- and therefore everything it says about pruning --- rests on a
compression hypothesis for the resolvent compressed to the spike frame: writing
\[
A_n(z)=U_n^*z(z^2-X_nX_n^*)^{-1}U_n,\qquad
D_n(z)=V_n^*z(z^2-X_n^*X_n)^{-1}V_n,
\]
\[
B_n(z)=U_n^*X_n(z^2-X_n^*X_n)^{-1}V_n,
\]
the hypothesis asks, in operator norm and uniformly on a contour $\Gamma$ off the spectrum, that
$A_n\approx\phi_nI$, $D_n\approx\widetilde\phi_nI$ and $B_n\approx0$, where $\phi_n$ and
$\widetilde\phi_n$ are the corresponding normalized traces. The source verifies the hypothesis in the
Haar class. The theorem extends it to entries that are i.i.d., symmetric and obey a dimension-free
Euclidean concentration inequality, throughout $r_n=o(n)$ (claim P19-IID-GC-BLOCK-02 in
\cite{PN24}), with
\begin{equation}\label{eq:prunerate}
\beta_n(\Gamma)\;:=\;\sup_{z\in\Gamma}\Big[\|A_n-\phi_nI\|_{\mathrm{op}}+\|D_n-\widetilde\phi_nI\|_{\mathrm{op}}
+\|B_n\|_{\mathrm{op}}\Big]
\;=\;O_{\mathbb P}\!\left(\sqrt{\frac{r}{n}}+\sqrt{\frac{\log n}{n}}\right);
\end{equation}
up to constants and the contour net; for independent entries with bounded moments the anisotropic
local law of Knowles and Yin \cite{KnowlesYin} supplies the compression for deterministic frames at
$r=o(\sqrt n)$. An extension to centred laws without the symmetry is recorded in the same track's
ledger without an archived proof. The concentration hypothesis remains, and nothing in the theorem
applies to a pruning frame selected from the same matrix it is applied to.

\subsection{What is measured}

Three things. The entry law is the dial. A centred Bernoulli$(0.2)$ variable is bounded, so it
satisfies the concentration hypothesis by Talagrand's inequality \cite{Talagrand}, and it is as
asymmetric as a two-point law with unit variance can be, with mass $0.8$ at $-1/2$ and mass $0.2$ at
$2$; a Student $t_5$ variable is exactly symmetric and has no exponential moment, so it fails the
concentration hypothesis; a centred exponential variable fails both. Rademacher and Gaussian entries
satisfy both. If the $t_5$ rows depart from the Rademacher rows, the concentration hypothesis is
load-bearing at these sizes; if the Bernoulli rows depart, the symmetry is.

Everything is read off one thin singular value decomposition per design. With
$X=\sum_k s_kw_k\tilde v_k^*$ and $\bar U=W^\top U$, $\bar V=\tilde VV$,
the three blocks are $\bar U^*\mathrm{diag}\big(z/(z^2-s_k^2)\big)\bar U$,
$\bar V^*\mathrm{diag}\big(z/(z^2-s_k^2)\big)\bar V+z^{-1}(I-\bar V^*\bar V)$ and
$\bar U^*\mathrm{diag}\big(s_k/(z^2-s_k^2)\big)\bar V$, so a whole contour costs nothing once the
decomposition is in hand. The contour is a circle in the right half plane standing off the observed
edge by fifteen percent of it, sampled at $64$ points, and $\phi_n,\widetilde\phi_n$ are the
empirical traces the source itself centres at. Table~\ref{tab:prunerate} reports $\beta_n$ and,
beside it, $\beta_n$ divided by the rate in~\eqref{eq:prunerate}: if that rate is right the ratio
does not run away with $r$ or with $n$. The three blocks contribute to $\beta_n$ in fixed
proportion at aspect $1/2$, the $A$ block $\numShareAMin$ to $\numShareAMax$, the $D$ block
$\numShareDMin$ to $\numShareDMax$ and the cross block $B$ $\numShareBMin$ to $\numShareBMax$ of the
sum, for every law, $n$ and $r$ run; the off-diagonal block, which the hypothesis wants to vanish, is
the largest term at every size.

\begin{table}[htbp]
\centering\small
\caption{The compression block on a contour off the spectrum, at aspect $c=n/m=0.5$ unless the row
says otherwise. $\beta_n$ is the supremum in~\eqref{eq:prunerate} over $64$ contour points;
the ratio divides it by $\sqrt{r/n}+\sqrt{\log n/n}$. Entries are mean $\pm$ standard deviation over
independent designs. Rows marked $\dagger$ lie outside the hypotheses of the theorem:
Bernoulli$(0.2)$ is bounded but not symmetric, $t_5$ is symmetric but not concentrated, and the
exponential law is neither.}
\label{tab:prunerate}
\begin{tabular}{lrrllr}
\toprule
entry law & $n$ & $r$ & $\beta_n$ & $\beta_n/(\sqrt{r/n}+\sqrt{\log n/n})$ & designs\\
\midrule
% BEGIN AUTOTABLE tab:prunerate
Gaussian & 1000 & 4 & $0.06\pm0.01$ & $0.40\pm0.05$ & 40\\
Gaussian & 1000 & 16 & $0.14\pm0.01$ & $0.65\pm0.04$ & 40\\
Gaussian & 1000 & 64 & $0.32\pm0.01$ & $0.96\pm0.04$ & 40\\
Gaussian & 1000 & 256 & $0.82\pm0.02$ & $1.39\pm0.04$ & 40\\
Rademacher & 500 & 16 & $0.20\pm0.02$ & $0.69\pm0.06$ & 20\\
Rademacher & 1000 & 4 & $0.06\pm0.01$ & $0.39\pm0.05$ & 40\\
Rademacher & 1000 & 16 & $0.14\pm0.01$ & $0.65\pm0.05$ & 60\\
Rademacher & 1000 & 64 & $0.32\pm0.01$ & $0.94\pm0.03$ & 60\\
Rademacher & 1000 & 256 & $0.82\pm0.02$ & $1.38\pm0.04$ & 40\\
Rademacher & 2000 & 16 & $0.09\pm0.01$ & $0.62\pm0.04$ & 20\\
Rademacher & 2000 & 32 & $0.14\pm0.00$ & $0.77\pm0.02$ & 16\\
Rademacher & 2000 & 128 & $0.33\pm0.01$ & $1.03\pm0.03$ & 16\\
Rademacher & 4000 & 16 & $0.07\pm0.00$ & $0.60\pm0.02$ & 8\\
Rademacher & 4000 & 64 & $0.14\pm0.00$ & $0.84\pm0.02$ & 8\\
Rademacher & 4000 & 256 & $0.33\pm0.00$ & $1.10\pm0.01$ & 8\\
Bernoulli$(0.2)$$^{\dagger}$ & 500 & 16 & $0.20\pm0.02$ & $0.69\pm0.05$ & 20\\
Bernoulli$(0.2)$$^{\dagger}$ & 1000 & 4 & $0.06\pm0.01$ & $0.39\pm0.06$ & 20\\
Bernoulli$(0.2)$$^{\dagger}$ & 1000 & 16 & $0.14\pm0.01$ & $0.65\pm0.03$ & 20\\
Bernoulli$(0.2)$$^{\dagger}$ & 1000 & 64 & $0.32\pm0.01$ & $0.94\pm0.04$ & 20\\
Bernoulli$(0.2)$$^{\dagger}$ & 1000 & 256 & $0.81\pm0.02$ & $1.38\pm0.04$ & 20\\
Bernoulli$(0.2)$$^{\dagger}$ & 2000 & 16 & $0.10\pm0.00$ & $0.65\pm0.03$ & 20\\
Bernoulli$(0.2)$$^{\dagger}$ & 4000 & 16 & $0.07\pm0.00$ & $0.62\pm0.02$ & 8\\
Bernoulli$(0.2)$$^{\dagger}$, aspect 1 & 1000 & 16 & $0.20\pm0.01$ & $0.93\pm0.06$ & 20\\
exponential$^{\dagger}$ & 1000 & 4 & $0.06\pm0.01$ & $0.40\pm0.05$ & 20\\
exponential$^{\dagger}$ & 1000 & 16 & $0.13\pm0.01$ & $0.64\pm0.04$ & 20\\
exponential$^{\dagger}$ & 1000 & 64 & $0.32\pm0.01$ & $0.94\pm0.03$ & 20\\
exponential$^{\dagger}$ & 1000 & 256 & $0.81\pm0.02$ & $1.38\pm0.03$ & 20\\
Student $t_5$$^{\dagger}$ & 500 & 16 & $0.19\pm0.03$ & $0.64\pm0.09$ & 20\\
Student $t_5$$^{\dagger}$ & 1000 & 4 & $0.06\pm0.01$ & $0.41\pm0.06$ & 20\\
Student $t_5$$^{\dagger}$ & 1000 & 16 & $0.14\pm0.01$ & $0.65\pm0.04$ & 20\\
Student $t_5$$^{\dagger}$ & 1000 & 64 & $0.32\pm0.01$ & $0.94\pm0.04$ & 20\\
Student $t_5$$^{\dagger}$ & 1000 & 256 & $0.76\pm0.09$ & $1.29\pm0.15$ & 20\\
Student $t_5$$^{\dagger}$ & 2000 & 16 & $0.09\pm0.01$ & $0.61\pm0.08$ & 20\\
Student $t_5$$^{\dagger}$ & 4000 & 16 & $0.07\pm0.00$ & $0.61\pm0.03$ & 8\\
Student $t_5$$^{\dagger}$, aspect 1 & 1000 & 16 & $0.20\pm0.01$ & $0.93\pm0.06$ & 20\\
% END AUTOTABLE tab:prunerate
\bottomrule
\end{tabular}
\end{table}

The third measurement is the adaptive spike. The theorem, like the source, assumes the spike frame
is deterministic or independent of the noise. A rank-one perturbation built from the \emph{observed}
top singular pair of the same $X$ leaves the empirical spectral distribution where it was, because a
rank-one change cannot move it by more than $1/n$, while behaving nothing like an independent spike.
The two arms are run side by side at matched $\theta$ --- same matrix, same rank, same strength, and
the only difference is whether the spike frame was drawn before looking at $X$ or after.

\subsection{What the entry law does, and what it does not}

At $n=1000$ the five entry laws agree to the precision of the measurement at every rank
(Table~\ref{tab:prunerate}). The asymmetric bounded law and the Rademacher law give
$\beta_n=\numBetaBernFour$ against $\numBetaRadFour$ at $r=4$, $\numBetaBernSixteen$ against
$\numBetaRadSixteen$ at $r=16$ and $\numBetaBernTwoFiveSix$ against $\numBetaRadTwoFiveSix$ at
$r=256$, standard errors over $\numNBernSixteen$ and $\numNRadSixteen$ designs at $r=16$. Symmetry
enters the proof at one identifiable place: the expectation of the clipped self-resolvent is made an
exact scalar matrix by invariance of the law under independent signed row and column permutations,
whose commutant is the scalars; without the signs the commutant gains the all-ones matrix, and the
expectation acquires a rank-one part whose contribution to the compressed block is of order $r$
times its coefficient. What is measured is the hypothesis as the source states it, on the unclipped
resolvent centred at the empirical trace, which absorbs any scalar part and leaves precisely the
non-scalar bias visible; and that bias does not reach the size of the fluctuation at these $r$ and
$n$. The bias that a proof without the symmetry would have to control is not visible here. The
heavy-tailed $t_5$ rows, outside the theorem, are indistinguishable from
the Rademacher rows except at $r=256$, where their spread across designs is four times larger, and the
exponential rows, outside it on both counts, are indistinguishable throughout. At these sizes neither
hypothesis of the theorem is visible in the compression error.

\subsection{The rate}

The theorem is an upper bound with an unspecified constant, and only its shape can be tested. At
fixed $n=1000$ the exponent of $\beta_n$ in $r$ reads $\numExpRadA$ from $r=4$ to $16$,
$\numExpRadB$ from $16$ to $64$ and $\numExpRadC$ from $64$ to $256$ for the Rademacher rows, and
$\numExpBernA$, $\numExpBernB$, $\numExpBernC$ for the Bernoulli rows, against the $1/2$ the rate
carries. Part of the excess at small $r$ is the finite-$r$ approach of an $r\times r$ random matrix's
operator norm to its $2\sqrt r$ asymptote from below: measured directly on the three shapes that
occur --- real symmetric, complex symmetric, complex Ginibre --- the operator norm divided by
$2\sqrt r$ runs $0.77$--$0.85$ at $r=4$, $0.91$--$0.94$ at $r=16$, $0.97$--$0.98$ at $r=64$ and
$0.99$ at $r=256$, so a quantity growing exactly like $\sqrt r$ would be measured with an apparent
exponent between $0.59$ and $0.62$ from $r=4$ to $16$, near $0.53$ from $16$ to $64$, and near
$0.51$ from $64$ to $256$. The first step is consistent with that; the second and third are not, and
the third rises where the finite-$r$ correction says it should fall. One confound has to be separated
before that is read as a statement about the bound: at $r=256$ and $n=1000$ the ratio $r/n$ is
$0.256$, outside the regime $r_n=o(n)$ the theorem is about, and there $\sqrt{r/n}=0.51$ is itself of
order one.

The experiment that separates the two is a ray of fixed $r/n$, along which the rate is constant up to
the $\sqrt{\log n/n}$ term (Table~\ref{tab:ray}). Along each ray $\beta_n$ does not fall: it rises
slightly and then flattens, while the rate itself falls by a fifth and an eighth over the same span
because its $\sqrt{\log n/n}$ term is shrinking. So that term is not describing anything at these
sizes: at fixed ratio the compression error converges to a constant rather than decaying with $n$.
What it converges to is worth writing down. Subtracting the leading term $\sqrt{r/n}$ from the
converged values at $n=4000$ leaves an excess of $1.12$ and $1.17$ times $r/n$ on the two rays, so
\begin{equation}\label{eq:betalaw}
\beta_n(\Gamma)\;\longrightarrow\;\sqrt{\frac rn}\;+\;c\,\frac rn,\qquad c\approx\numRayC ,
\end{equation}
with a leading coefficient consistent with exactly one. Fitted on those two points alone,
\eqref{eq:betalaw} predicts the lanes it did not see: $\numRayPredA$ against $\numRayMeasA$
measured at $(r,n)=(32,2000)$, $\numRayPredB$ against $\numRayMeasB$ at $(128,2000)$, and
$\numRayPredC$ against $\numRayMeasC$ at $(16,4000)$ --- within a percent on the first two and
$\numRayOffC$ percent on the third, the smaller-$n$ points sitting below the law because they have
not finished rising.

In the regime the theorem is about, $r_n=o(n)$, the second term of \eqref{eq:betalaw} is
$o(\sqrt{r/n})$ and the bound's rate is the truth --- with, on this evidence, constant one rather
than an unspecified $C$. The exponent above $1/2$ measured at fixed $n$ is the sum of the two terms of
\eqref{eq:betalaw} read over a range of $r$ where $r/n$ climbed to $0.256$ and the second term was
no longer small. The sharper statement the measurements support is not that the compression error is
$O(\sqrt{r/n}+\sqrt{\log n/n})$ but that it approaches $\sqrt{r/n}+c\,r/n$, with no visible $\log n$
at all.

\begin{table}[t]
\centering\small
\caption{Rademacher rows along two rays of fixed $r/n$, at aspect $c=n/m=0.5$. $\beta_n$ is mean
$\pm$ standard error over the designs; the excess is $\beta_n-\sqrt{r/n}$ and the last ratio divides
it by $r/n$.}
\label{tab:ray}
\begin{tabular}{rrrlrrrr}
\toprule
$r$ & $n$ & $r/n$ & $\beta_n$ & $\sqrt{r/n}$ & excess & excess$/(r/n)$ & designs\\
\midrule
% BEGIN AUTOTABLE tab:ray
16 & 1000 & 0.016 & $0.1356\pm0.0015$ & 0.1265 & +0.0092 & 0.57 & 60\\
32 & 2000 & 0.016 & $0.1441\pm0.0011$ & 0.1265 & +0.0176 & 1.10 & 16\\
64 & 4000 & 0.016 & $0.1444\pm0.0012$ & 0.1265 & +0.0179 & 1.12 & 8\\
64 & 1000 & 0.064 & $0.3176\pm0.0013$ & 0.2530 & +0.0646 & 1.01 & 60\\
128 & 2000 & 0.064 & $0.3254\pm0.0022$ & 0.2530 & +0.0724 & 1.13 & 16\\
256 & 4000 & 0.064 & $0.3276\pm0.0014$ & 0.2530 & +0.0746 & 1.17 & 8\\
% END AUTOTABLE tab:ray
\bottomrule
\end{tabular}
\end{table}

\subsection{The frame you chose after looking}

Table~\ref{tab:pruneadapt} is the counterexample with numbers on it. The two arms differ in one
respect only: whether the rank-one spike's frame was drawn before or after seeing $X$. Below the
transition the independent spike does what the independent-spike theory \cite{BGN} says it does ---
at $\theta=0.6$ and $0.8$ the overlap with the planted direction is $0.003$ and $0.006$, which is to
say nothing is recovered --- while the frame read off $X$ has overlap $1.000$ to three decimals at
every $\theta$, with an outlier sitting at the unperturbed edge plus $\theta$ exactly. And the bulk
cannot tell the two apart: the Kolmogorov--Smirnov distance between the perturbed and unperturbed
spectra is $1.00/n$ in both arms at every $\theta$, which is the smallest value the statistic can
take on a sample of this size. Both arms are rank one, both leave a textbook Marchenko--Pastur bulk
\cite{MarchenkoPastur}, and one of them recovers a direction that was never planted.

\begin{table}[htbp]
\centering\small
\caption{The same rank-one spike, drawn independently of $X$ and drawn from the observed top
singular pair of $X$. Outlier is the leading singular value of $X+\theta uv^*$; overlap is
$|\langle\hat u_1,u\rangle|^2$; the last two columns are the Kolmogorov--Smirnov distance between
the perturbed and unperturbed bulk --- everything except the leading singular value --- in units of
$1/n$, which is the resolution of the statistic on a sample of this size. Entries
are mean $\pm$ standard deviation over independent designs; Rademacher entries throughout.}
\label{tab:pruneadapt}
{\setlength{\tabcolsep}{4.5pt}
\begin{tabular}{rrllllrrr}
\toprule
& & \multicolumn{2}{c}{independent frame} & \multicolumn{2}{c}{frame read off $X$} &
\multicolumn{2}{c}{$n\times$ bulk KS} & \\
\cmidrule(lr){3-4}\cmidrule(lr){5-6}\cmidrule(lr){7-8}
$n$ & $\theta$ & outlier & overlap & outlier & overlap & indep. & adapt. & designs\\
\midrule
% BEGIN AUTOTABLE tab:pruneadapt
1000 & 0.6 & $2.41\pm0.01$ & $0.003\pm0.004$ & $3.01\pm0.01$ & $1.000\pm0.000$ & 1.00 & 1.00 & 20\\
1000 & 0.8 & $2.41\pm0.01$ & $0.006\pm0.009$ & $3.21\pm0.01$ & $1.000\pm0.000$ & 1.00 & 1.00 & 20\\
1000 & 1 & $2.41\pm0.01$ & $0.020\pm0.022$ & $3.41\pm0.01$ & $1.000\pm0.000$ & 1.00 & 1.00 & 20\\
1000 & 1.2 & $2.42\pm0.01$ & $0.078\pm0.050$ & $3.61\pm0.01$ & $1.000\pm0.000$ & 1.00 & 1.00 & 20\\
1000 & 1.4 & $2.44\pm0.02$ & $0.269\pm0.085$ & $3.81\pm0.01$ & $1.000\pm0.000$ & 1.00 & 1.00 & 20\\
1000 & 1.6 & $2.51\pm0.02$ & $0.479\pm0.032$ & $4.01\pm0.01$ & $1.000\pm0.000$ & 1.00 & 1.00 & 20\\
1000 & 2 & $2.73\pm0.03$ & $0.695\pm0.014$ & $4.41\pm0.01$ & $1.000\pm0.000$ & 1.00 & 1.00 & 20\\
1000 & 2.5 & $3.08\pm0.03$ & $0.817\pm0.007$ & $4.91\pm0.01$ & $1.000\pm0.000$ & 1.00 & 1.00 & 20\\
1000 & 3 & $3.49\pm0.03$ & $0.877\pm0.004$ & $5.41\pm0.01$ & $1.000\pm0.000$ & 1.00 & 1.00 & 20\\
2000 & 1 & $2.41\pm0.01$ & $0.006\pm0.011$ & $3.41\pm0.01$ & $1.000\pm0.000$ & 1.00 & 1.00 & 12\\
2000 & 1.4 & $2.45\pm0.01$ & $0.338\pm0.030$ & $3.81\pm0.01$ & $1.000\pm0.000$ & 1.00 & 1.00 & 12\\
2000 & 2 & $2.75\pm0.01$ & $0.702\pm0.013$ & $4.41\pm0.01$ & $1.000\pm0.000$ & 1.00 & 1.00 & 12\\
% END AUTOTABLE tab:pruneadapt
\bottomrule
\end{tabular}}
\end{table}

The practical reading follows. A pruning frame taken from the singular vectors of the matrix being
pruned is the second arm, not the first, and no amount of agreement between the observed
singular-value histogram and its Marchenko--Pastur prediction licenses the independent calculation;
the bulk is exactly the statistic that cannot see the difference. What is needed is a frame built
without looking, or a stability argument connecting the learned frame to one --- the same
same-sample problem that Section~\ref{sec:corr} handles for a mean and a correction. One more thing
the frame is not: a substitute for the ridge. For kernel ridge regression truncated to the top $r$
eigenvalues of the Gram matrix, Amini \cite{Amini} gives the worst-case risk over the unit ball of
the RKHS exactly,
$\max\{H_r(\lambda),\mu_{r+1}\}+\frac{\sigma^2}{n}\sum_{i\le r}\big(\frac{\mu_i}{\mu_i+\lambda}\big)^2$
with $H_r(\lambda)=\max_{i\le r}\lambda^2\mu_i/(\mu_i+\lambda)^2$, and shows that truncating at any
$r$ at or above the level the ridge itself resolves never raises the minimax risk and strictly
lowers it whenever $\mu_{r+1}>0$; the implicit regularisation of the cut and the explicit one of
$\lambda$ are both needed. A pruning theory that assumes the ridge away, or a ridge path that assumes
the frame away, is describing half of the estimator.

\begin{remark}
The object the fifth problem describes is easy to see in a trained network. For five weight matrices
of a vision transformer trained on ImageNet (timm's \texttt{vit\_base\_patch16\_224}, the first,
seventh and last MLP input layers, $3072\times768$, one attention projection $2304\times768$ and one
output projection $768\times768$), the number of singular values above the Marchenko--Pastur edge
computed from each matrix's own entry standard deviation is $67$, $78$, $36$, $78$ and $2$ of $768$,
against $0$ of $768$ for i.i.d.\ Gaussian matrices of the same shapes. This is one checkpoint, the
edge is estimated from a matrix that contains the spikes and is therefore conservative, and the frame
those singular vectors would supply is the second arm of Table~\ref{tab:pruneadapt} by construction:
it shows the object, not the theorem.
\end{remark}



\section{Regularisation, overfitting, and a learned mean, on two corpora}\label{sec:reg}

Targets~2 and~3 ask when an interpolating kernel estimator is harmless. The same question can be put
to data along the whole regularisation path rather than at its ridgeless end, with the kernel placed
on top of a network rather than standing alone. It is put here to two corpora: the one-dimensional
Burgers problem of PDEBench \cite{PDEBench} and the advection problem of the operator-learning
benchmark of de Hoop, Huang, Qian and Stuart \cite{deHoop}, a setting in which kernel methods are
competitive with neural operators \cite{BatlleDarcyHosseiniOwhadi} and Gaussian processes solve the
underlying equations directly \cite{ChenHosseiniOwhadiStuart}. The evidence is thin and is reported as
such: one seed of the regularisation path and one of the network-mean gate on Burgers, three seeds
of the constant-mean gate there, two seeds of the dimension-gap diagnostic on Burgers and one on
advection. None of it is a population statement; each table is one or two draws.

\subsection{The path as a learning curve}

For each corpus and sample size the ridge estimator is fitted at thirteen values of $\lambda$ from
$10^{-12}$ to $1$ and its test risk recorded; the interpolant is the $\lambda\to0$ end of the path.
Two exponents are read off the same Gram matrix and labels the path was fitted on: the capacity
$\hat\alpha$, the slope of log eigenvalue against log index, and the source $\hat r$, from the decay
of the target's coefficients in the Gram eigenbasis --- the source and capacity conditions of
\cite{CaponnettoDeVito}. In the parametrisation of Cui, Loureiro, Krzakala and Zdeborov\'a
\cite{Cui}, ridge regression has a noiseless regime and a noisy one: the noiseless decay
$n^{-2\alpha\min(r,1)}$ of the excess risk holds for $n$ well below
$n_1^*\asymp\sigma^{-1/(\alpha\min(r,1))}$, and the noisy optimal rate takes over for $n$ well
above $n_2^*\asymp\sigma^{-\max(2,\,1/(\alpha\min(r,1)))}$. The statistic that says whether
interpolating was harmless is the ratio of the interpolant's risk to the best risk on the path.

On Burgers at $n=4000$ with no label noise (Table~\ref{tab:regpath}) the interpolant is the optimum
of the path --- the ratio is $1.00$ and the best ridge is the smallest one tried --- with
$\hat\alpha=\numBurgAlpha$ and $\hat r=\numBurgR$, so that $1/(\hat\alpha\min(\hat r,1))=\numBurgExp$:
in the theory of \cite{Cui} a tenfold increase in the noise moves the end of the noiseless decay by a
factor of about $\numBurgFactor$ in $n$, and the takeover of the noisy rate, governed by
$\max(2,\numBurgExp)=\numBurgTakeoverExp$, by a factor of $\numBurgTakeover$. Past the takeover the
same theory fixes the shape of the ratio itself: with
$\lambda\to0$ the excess risk plateaus at the noise variance while the optimally regularised risk
keeps falling like $n^{-2\alpha\min(r,1)/(1+2\alpha\min(r,1))}$, so the ratio should grow with $n$
at that exponent, $\numBurgRate$ for Burgers. The lanes at the other noise levels, which would test
both predictions, were not run (Table~\ref{tab:crossover} carries the one noiseless cell), so on
this corpus the predictions stand as predictions.

\begin{table}[t]
\centering\small
\caption{The regularisation path with no label noise, one seed. $\hat\alpha$ and $\hat r$ are the
capacity and source exponents read off the Gram spectrum; the fourth column is
$1/(\hat\alpha\min(\hat r,1))$, the exponent of the noise level at which the noiseless decay ends in
\cite{Cui}; $\lambda_\ast$ is the best ridge; the ratio is the
interpolant's test risk over the best risk on the path, so $1.00$ means interpolation was the
optimum.}
\label{tab:regpath}
\begin{tabular}{@{}llrcccccr@{}}
\toprule
corpus & comp. & $n$ & $\hat\alpha$ & $\hat r$ & exponent & $\lambda_\ast$ & interp./best & seeds\\
\midrule
% BEGIN AUTOTABLE tab:regpath
burgers &  & 4000 & 2.02 & 0.63 & 0.78 & $10^{-12}$ & 1.00 & 1\\
% END AUTOTABLE tab:regpath
\bottomrule
\end{tabular}
\end{table}

\begin{table}[t]
\centering\small
\caption{The ratio interpolant over best ridge as the label noise is turned up. Only the noiseless
cell was run.}
\label{tab:crossover}
\begin{tabular}{@{}llrcccc@{}}
\toprule
corpus & comp. & $n$ & noise $0$ & $0.03$ & $0.1$ & $0.3$\\
\midrule
% BEGIN AUTOTABLE tab:crossover
burgers &  & 4000 & 1.00 & -- & -- & --\\
% END AUTOTABLE tab:crossover
\bottomrule
\end{tabular}
\end{table}

\subsection{The gate with a learned mean}

The identity of Section~\ref{sec:corr} gates a kernel correction on top of a mean by leave-one-out
risk. Here the mean is either a constant or a multilayer perceptron of width $256$ trained for sixty
epochs, so that the network's own fit and the correction's ridge $\lambda$ enter together;
leave-one-out, which stays consistent along iterative paths where generalised cross-validation
\cite{GHW} does not \cite{Patil} and which is the selection principle of kernel flows
\cite{OwhadiYoo,HamziOwhadi}, is used throughout. The question is whether the gate keeps doing its
two jobs once the mean has itself been fitted to the training labels: taking the gain when the
correction helps, and refusing the step when it hurts. The correction the gate acts on is fitted to
the residuals against the trained mean's own training-row predictions --- the consistent arm of
Section~\ref{sec:corr} --- so the gains below, and the share of them the gate captures, are that
arm's; the out-of-fold arm is not measured here.

On Burgers (Table~\ref{tab:gatemlp}) the network mean reaches a training error of
$\numGateTrainMse$ and the gate still takes $\numGateCaptured$ percent of the remaining gain with
$\hat\theta=\numGateTheta$; the mean's own risk of $\numGateRiskMean$ falls to $\numGateRiskGated$
against a best of $\numGateRiskBest$. On the advection suite with a constant mean the correction is
the whole model, from a risk of $\numAdvRiskMean$ to $\numAdvRiskGated$ gated against a best of
$\numAdvRiskBest$, and the gate takes it. On none of the five designs does the correction, gated or
not, leave the test risk above the mean's own, so the gate's second job --- refusing a harmful step
--- is not exercised on these corpora.

\begin{table}[t]
\centering\small
\caption{The leave-one-out gate on top of a constant mean and on top of a network mean. $w$ is the
network's weight decay; $\hat\theta$ the step the gate chooses; ``captured'' the share of the gain
between the mean's risk and the best achievable risk that the gated step takes; ``worse'' the
number of designs, gated and ungated, on which the correction left the test risk above the mean's
own.}
\label{tab:gatemlp}
\begin{tabular}{@{}lllcccr@{}}
\toprule
corpus & mean & $w$ & $\hat\theta$ & captured & worse (gated / full) & designs\\
\midrule
% BEGIN AUTOTABLE tab:gatemlp
burgers  & const & -- & $0.95\pm0.00$ & $1.00\pm0.00$ & 0 / 0 & 3\\
burgers  & mlp & 0 & 0.99 & 0.99 & 0 / 0 & 1\\
caltech\_adv  & const & -- & 0.95 & 1.00 & 0 / 0 & 1\\
% END AUTOTABLE tab:gatemlp
\bottomrule
\end{tabular}
\end{table}

\subsection{The order parameter on data that has none}

The diagnostic of Section~\ref{sec:gap} maps the dimension gap onto a corpus through the Gram
spectrum: $d_{\rm eff}$ is the number of directions resolved above the numerical floor,
$p_{\rm eff}$ the participation ratio of the spectrum, standing in for the ambient dimension, and
$\delta_{\rm eff}=\log(d_{\rm eff}/n)/\log p_{\rm eff}$. Beside it stands the thing this is supposed
to predict: the advantage of ridgeless interpolation over the best ridge on a swept grid, as a
fraction of the best risk, negative when interpolating loses (Table~\ref{tab:newcorp}).

\begin{table}[t]
\centering\small
\caption{The empirical dimension gap and the interpolation advantage
$(R_{\rm best}-R_{\rm interp})/R_{\rm best}$ on the two corpora, ordered by $\delta_{\rm eff}$;
mean $\pm$ s.d.\ over the seeds where there are two.}
\label{tab:newcorp}
\begin{tabular}{@{}llrrrccr@{}}
\toprule
corpus & comp. & $n$ & $d_{\rm eff}$ & $p_{\rm eff}$ & $\delta_{\rm eff}$ & advantage & seeds\\
\midrule
% BEGIN AUTOTABLE tab:newcorp
caltech\_adv &  & 6000 & 5164 & 4.7 & -0.096 & -0.88 & 1\\
burgers &  & 6000 & 5382 & 5.1 & $-0.067\pm0.000$ & $0.25\pm0.01$ & 2\\
% END AUTOTABLE tab:newcorp
\bottomrule
\end{tabular}
\end{table}

On the advection suite interpolating loses outright: at $n=6000$ the best ridge,
$\lambda=\numAdvBestLam$, has risk $\numAdvRiskRidge$ and the interpolant $\numAdvRiskInterp$, an
advantage of $\numAdvAdvantage$, with $\delta_{\rm eff}=\numAdvDeltaEff$. On Burgers at the same $n$
the interpolant wins by a quarter of the best risk on both seeds, at a $\delta_{\rm eff}$ that is
less negative. Two corpora are two points, and a rank correlation between $\delta_{\rm eff}$ and the
advantage across corpora --- the statistic that would say whether the order parameter carries over
to data --- needs many more. The mapping from the sphere to a corpus is a modelling choice, and this
is where it would be tested.

The theory side of Table~\ref{tab:newcorp} has been written down in the same coordinate. Lu, Lin,
Xia and Huang \cite{LuLinXiaHuang} take $n\asymp d^{\gamma}$ on the sphere, a source parameter
$s\ge0$ for the smoothness of the target, and the ridge $\lambda\asymp d^{-u}$, and prove that
overfitting is benign --- in the under-regularised regime and at interpolation alike --- exactly
when $0<s\le\Gamma(\gamma)$, for a threshold function recalled in \cite{LuLinXiaHuang} from earlier
work, where with $\ell=\lfloor\gamma\rfloor$
\[
\Gamma(\gamma)=\frac{\gamma-\ell}{\ell}\quad(\ell<\gamma\le\ell+\tfrac12),\qquad
\Gamma(\gamma)=\frac{\ell+1-\gamma}{\ell+1}\quad(\ell+\tfrac12<\gamma<\ell+1),
\]
$\Gamma=\infty$ for $\gamma\le\frac12$ and $\Gamma=0$ at every integer $\gamma$. Since the
harmonic block of degree $k$ has dimension $N_{d,k}\asymp d^k$, the fractional part of $\gamma$ is the
dimension gap of Section~\ref{sec:gap} for the last resolved block, $\delta=k-\gamma$ at level $k$: the
threshold closes when $n$ sits at a block boundary, where Target~3's noise floor is exact, and it is
widest when $n$ sits at the geometric middle of two blocks. At the sizes of Section~\ref{sec:t3},
$n=2^d$ gives $\gamma=d\log2/\log d$, which is $3.35$ at $d=12$ and $3.51$ at $d=13$, so
$\Gamma(\gamma)$ is $0.12$ at both: in that asymptotic only targets of source below about $0.12$
overfit benignly on those designs, with $d$ far from the regime the theorem is stated in. In that
reading $\delta_{\rm eff}$ is an estimate of the coordinate on which the benignity threshold is a
known function, and the taxonomy of benign, tempered and catastrophic overfitting \cite{Mallinar} is
a description of the same axis in the spectral coordinates of \cite{Cui}.

\section{What the measurements say}

\begin{enumerate}
\item The proved lower bound for Target~3 is conservative at reachable sizes by a factor between
$\numRatioMinHalf$ and $\numRatioMaxHalf$ at $h=1/2$, between $\numRatioMinOne$ and
$\numRatioMaxOne$ at $h=1$, and between $\numRatioMinTwo$ and $\numRatioMaxTwo$ at $h=2$, against
the ten-window constant $c_{10}=\numCten$, the strengthening of Theorem~\ref{thm:T3} quoted in
Remark~\ref{rem:sharper}. Its constants are independent of bandwidth by construction while the
measured functional is not; the gap between them is where a sharp,
$h$-dependent constant would live. The single-block certificate the one-window argument bounds sits
above its constant $c_1=5/36$ at every dimension and bandwidth measured, following the sawtooth of
the block dimension rather than a trend, and the expectation bound on the captured leverage,
\eqref{eq:capture}, is the loose step at these sizes.
\item Within a harmonic block, at fixed critical degree, the noise risk falls as the dimension gap
widens, at seven dimensions and one bandwidth and out to $\delta=0.6$, with the few inversions
confined to blocks of small harmonic dimension or to cells with two designs; the tail isotropy that
Target~2's proof
drives to zero follows the same ordering.
\item Both lemmas that Target~4's reduction names look true at finite size, one of them in a
stronger form than the reduction requires, and the second with its constant identified rather than
merely bounded: $d\,\Tr(A^4)\to$ a limit statistically indistinguishable from $2$, which is the
fourth moment of the semicircle law and therefore the value a uniformly random direction in
$\mathrm{Sym}_d$ would give. The leave-one-out normal is generic at leading order; the Veronese
structure of the rows appears only in the $1/d$ correction, where it is an order of magnitude larger
than the generic one and is set by the rank-one dependence of the row, not by its kurtosis. After the
reduction the matrix is consistent with the square Wishart hard edge at every dimension measured.
The remaining difficulty is the rectangular hard edge for the traceless harmonic block itself, and
one thing that is known not to be available on the way there is uniform integrability of the full
inverse smallest singular value, which diverges at every fixed $d$.
\item The leverage concentration that lets Target~1's argument replace a correlated Stein comparison
with second-moment matching shows $d\,\Var(q)$ of order one at $d=64$ and $d=128$, with a clear
dependence on the augmentation strength, no dependence on the column normalization, and nothing that
one doubling of $d$ can say about the limit; the residual coefficients that argument needs have a
one-line Hermite identification with nonnegativity built in (Proposition~\ref{prop:hermite}).
\item On the EMIT and OCO-2 corpora the same-sample correction identity is exact to
$\numIdentWorst$, and its sign is computable per component and per ridge. On two of the four EMIT
components the margin is positive at every ridge on every seed; on the third it changes sign at
small ridge on some seeds; on the fourth, where the mean is poor, it is negative almost everywhere
and the correction does nothing; on the three OCO-2 bands the correction is close to neutral. Where
the consistent arm helps, its gain in squared risk transfers to the pipeline's physical error one for
one. The mismatch bound holds and is loose by two to three orders of magnitude, its transfer
condition fails on every (lane, ridge) pair, and the arm the pipeline fits, to out-of-fold
residuals, gains nothing on these lanes.
\item For the fifth problem, the compression error of the spike-frame resolvent is indistinguishable
across five entry laws inside and outside the hypotheses of the theorem, at every rank and size run.
Along rays of fixed $r/n$ it approaches $\sqrt{r/n}+\numRayC\,r/n$, with no visible $\log n$ term
and a leading coefficient of one; the exponent above $1/2$ seen at fixed $n$ is the second term.
Two rank-one spikes with bulk spectra that differ by the resolution of the statistic and nothing
more have overlaps of $0.003$ and $1.000$ with the planted direction, according to whether the frame
was drawn before or after looking at the matrix.
\item On the two corpora of the regularisation block, one seed at a time, the noiseless Burgers
interpolant is the optimum of the path at $n=4000$, and on the advection suite interpolating loses to
the best ridge by $88$ percent of the best risk, at a more negative empirical dimension gap; the
leave-one-out gate keeps its first property when the mean it corrects is a trained network, taking
$\numGateCaptured$ percent of the gain on Burgers, and its second property is not exercised.
\end{enumerate}

\section*{Funding}
The research presented in this paper was supported by the European Research Council (ERC) under the
European Union's Horizon Europe research and innovation programme (grant agreement No.~101041711), by
the Simons Foundation as part of the Collaboration on the Mathematical and Scientific Foundations of
Deep Learning, by Heights Labs, by Convex Nexus Capital, by the Israel Science Foundation (grant
number 2258/19), and by the Israel Science Foundation (ISF Grant 4101/25).

\section*{Declaration of competing interest}
The author holds equity in Heights Labs and in Convex Nexus Capital and has a management role at
Convex Nexus Capital. These entities are included in the funding declaration above.

\section*{Declaration of generative AI and AI-assisted technologies}
This paper is a product of an ongoing project on the use of AI systems for research mathematics,
provisionally called Math Brain, which studies how to make such systems better at solving open
problems; the project and its tooling will be released publicly in due course. The proofs of
Part~I were found by the automated programme ProofNet24 \cite{PN24}, whose workers run on OpenAI's
Codex, and were rewritten, checked and typeset by Claude (Anthropic; Claude Fable 5.1). The
experiments were designed, the code written and run, the records audited and this manuscript
drafted by Claude (Claude Opus 5 and Claude Fable 5.1) and Codex (OpenAI) sessions between 14 and 16
September 2026. The author set the goals of the project, commissioned and directed the work and
approved the release; the author did not carry out the mathematics, the computations or the writing,
and takes responsibility for the decision to publish.

\section*{Data and code availability}
The drivers, the \numRecords{} run records, the scripts of the laptop-scale computations quoted in
Sections~\ref{sec:semicircle} and~\ref{sec:prune}, and the build script that regenerates every table
and quoted number from the run records and checks the identities \eqref{eq:second}, \eqref{eq:V},
\eqref{eq:reduction}, \eqref{eq:identity} and \eqref{eq:null} against them are in the repository
\texttt{yspennstate/math-open-problems}, under \texttt{paper9\_kernel\_open\_problems}, together
with the ProofNet24 run reports under \texttt{proofnet24}, one directory per track with every run
preserved under its own timestamp. The repository is private and is available from the author on
request.

\begin{thebibliography}{99}

\bibitem{Amini} A. A. Amini, Spectrally-truncated kernel ridge regression and its free lunch,
\emph{Electron. J. Stat.} 15 (2021), 3743--3761. arXiv:1906.06276.

\bibitem{AGZ} G. W. Anderson, A. Guionnet and O. Zeitouni, \emph{An Introduction to Random
Matrices}, Cambridge Studies in Advanced Mathematics 118, Cambridge University Press, 2010.

\bibitem{AtkinsonHan} K. Atkinson and W. Han, \emph{Spherical Harmonics and Approximations on the
Unit Sphere: An Introduction}, Lecture Notes in Mathematics 2044, Springer, 2012.

\bibitem{BBP} J. Baik, G. Ben Arous and S. P\'ech\'e, Phase transition of the largest eigenvalue
for nonnull complex sample covariance matrices, \emph{Ann. Probab.} 33 (2005), 1643--1697.

\bibitem{BBvH} A. S. Bandeira, M. T. Boedihardjo and R. van Handel, Matrix concentration
inequalities and free probability, \emph{Invent. Math.} 234 (2023), 419--487.

\bibitem{Bartlett} P. L. Bartlett, P. M. Long, G. Lugosi and A. Tsigler, Benign overfitting in
linear regression, \emph{Proc. Natl. Acad. Sci. USA} 117 (2020), 30063--30070.


\bibitem{BatlleDarcyHosseiniOwhadi} P. Batlle, M. Darcy, B. Hosseini and H. Owhadi, Kernel methods are competitive for operator learning, arXiv:2304.13202 (2023).
\bibitem{BGN} F. Benaych-Georges and R. R. Nadakuditi, The singular values and vectors of low rank
perturbations of large rectangular random matrices, \emph{J. Multivariate Anal.} 111 (2012),
120--135.

\bibitem{BLM} S. Boucheron, G. Lugosi and P. Massart, \emph{Concentration Inequalities: A
Nonasymptotic Theory of Independence}, Oxford University Press, 2013.

\bibitem{BvH} T. Brailovskaya and R. van Handel, Universality and sharp matrix concentration
inequalities, \emph{Geom. Funct. Anal.} 34 (2024), 1734--1838. arXiv:2201.05142.

\bibitem{CaponnettoDeVito} A. Caponnetto and E. De Vito, Optimal rates for the regularized
least-squares algorithm, \emph{Found. Comput. Math.} 7 (2007), 331--368.


\bibitem{ChenHosseiniOwhadiStuart} Y. Chen, B. Hosseini, H. Owhadi and A. M. Stuart, Solving and learning nonlinear PDEs with Gaussian processes, arXiv:2103.12959 (2021).
\bibitem{Cui} H. Cui, B. Loureiro, F. Krzakala and L. Zdeborov\'a, Generalization error rates in
kernel regression: the crossover from the noiseless to noisy regime, in \emph{Advances in Neural
Information Processing Systems} 34 (2021), 10131--10143. arXiv:2105.15004.

\bibitem{deHoop} M. V. de Hoop, D. Z. Huang, E. Qian and A. M. Stuart, The cost-accuracy trade-off
in operator learning with neural networks, \emph{J. Mach. Learn.} 1 (2022), 299--341.
arXiv:2203.13181.


\bibitem{dlPG} V. H. de la Pe\~na and E. Gin\'e, \emph{Decoupling: From Dependence to Independence}, Springer, 1999.

\bibitem{dlPMS} V. H. de la Pe\~na and S. J. Montgomery-Smith, Decoupling inequalities for the tail probabilities of multivariate $U$-statistics, \emph{Ann. Probab.} 23 (1995), 806--816.
\bibitem{DhifallahLu} O. Dhifallah and Y. M. Lu, A precise performance analysis of learning with
random features, arXiv:2008.11904 (2020).


\bibitem{DLMF} NIST Digital Library of Mathematical Functions, \url{https://dlmf.nist.gov/}, F. W. J. Olver et al., eds.
\bibitem{Edelman} A. Edelman, Eigenvalues and condition numbers of random matrices, \emph{SIAM J.
Matrix Anal. Appl.} 9 (1988), 543--560.

\bibitem{EfronStein} B. Efron and C. Stein, The jackknife estimate of variance, \emph{Ann.
Statist.} 9 (1981), 586--596.

\bibitem{Eldering} A. Eldering et al., The Orbiting Carbon Observatory-2: first 18 months of
science data products, \emph{Atmos. Meas. Tech.} 10 (2017), 549--563.

\bibitem{Ghorbani} B. Ghorbani, S. Mei, T. Misiakiewicz and A. Montanari, Linearized two-layers
neural networks in high dimension, \emph{Ann. Statist.} 49 (2021), 1029--1054.

\bibitem{Goldt} S. Goldt, B. Loureiro, G. Reeves, F. Krzakala, M. M\'ezard and L. Zdeborov\'a, The
Gaussian equivalence of generative models for learning with shallow neural networks, in
\emph{Proceedings of the 2nd Mathematical and Scientific Machine Learning Conference}, PMLR 145
(2022), 426--471.

\bibitem{GHW} G. H. Golub, M. Heath and G. Wahba, Generalized cross-validation as a method for
choosing a good ridge parameter, \emph{Technometrics} 21 (1979), 215--223.

\bibitem{Green} R. O. Green et al., The Earth Surface Mineral Dust Source Investigation: an Earth
science imaging spectroscopy mission, in \emph{2020 IEEE Aerospace Conference}, IEEE, 2020.
doi:10.1109/AERO47225.2020.9172731.

\bibitem{GregorattiMaran} G. Gregoratti and D. Maran, Least singular value and condition number of a
square random matrix with i.i.d.\ rows, \emph{Statist. Probab. Lett.} 173 (2021), 109070.


\bibitem{HamziOwhadi} B. Hamzi and H. Owhadi, Learning dynamical systems from data: a simple cross-validation perspective, part I: parametric kernel flows, \emph{Physica D} 421 (2021), 132817.
\bibitem{HuLu} H. Hu and Y. M. Lu, Universality laws for high-dimensional learning with random
features, \emph{IEEE Trans. Inform. Theory} 69 (2023), 1932--1964. arXiv:2009.07669.


\bibitem{Janson} S. Janson, \emph{Gaussian Hilbert Spaces}, Cambridge Tracts in Mathematics 129, Cambridge University Press, 1997.

\bibitem{JSZ} W. B. Johnson, G. Schechtman and J. Zinn, Best constants in moment inequalities for linear combinations of independent and exchangeable random variables, \emph{Ann. Probab.} 13 (1985), 234--253.

\bibitem{KRM} C. Kaushik, J. Romberg and V. Muthukumar, A general technique for approximating high-dimensional empirical kernel matrices, arXiv:2511.03892 (2025).
\bibitem{KnowlesYin} A. Knowles and J. Yin, Anisotropic local laws for random matrices,
\emph{Probab. Theory Related Fields} 169 (2017), 257--352.

\bibitem{Lamminpaa} O. Lamminp\"a\"a, J. Susiluoto, J. Hobbs, J. McDuffie, A. Braverman and
H. Owhadi, Forward model emulator for atmospheric radiative transfer using Gaussian processes and
cross validation, \emph{Atmos. Meas. Tech.} 18 (2025), 673--694.


\bibitem{LedouxTalagrand} M. Ledoux and M. Talagrand, \emph{Probability in Banach Spaces: Isoperimetry and Processes}, Springer, 1991.
\bibitem{LiangRakhlin} T. Liang and A. Rakhlin, Just interpolate: kernel ``ridgeless'' regression
can generalize, \emph{Ann. Statist.} 48 (2020), 1329--1347.

\bibitem{LuLinXiaHuang} W. Lu, Q. Lin, Y. Xia and D. Huang, Learning curves and benign overfitting
of spectral algorithms in large dimensions, arXiv:2604.23212 (2026).

\bibitem{Mallinar} N. Mallinar, J. B. Simon, A. Abedsoltan, P. Pandit, M. Belkin and P. Nakkiran,
Benign, tempered, or catastrophic: toward a refined taxonomy of overfitting, in \emph{Advances in
Neural Information Processing Systems} 35 (2022). arXiv:2207.06569.

\bibitem{MarchenkoPastur} V. A. Mar\v{c}enko and L. A. Pastur, Distribution of eigenvalues for some
sets of random matrices, \emph{Math. USSR-Sb.} 1 (1967), 457--483.

\bibitem{MeiMontanari} S. Mei and A. Montanari, The generalization error of random features
regression: precise asymptotics and the double descent curve, \emph{Comm. Pure Appl. Math.} 75
(2022), 667--766.

\bibitem{MinhNiyogiYao} H. Q. Minh, P. Niyogi and Y. Yao, Mercer's theorem, feature maps, and
smoothing, in \emph{Learning Theory (COLT 2006)}, Lecture Notes in Computer Science 4005, Springer,
2006, 154--168.

\bibitem{Misiakiewicz} T. Misiakiewicz, Spectrum of inner-product kernel matrices in the polynomial
regime and multiple descent phenomenon in kernel ridge regression, arXiv:2204.10425 (2022).


\bibitem{Nelson} E. Nelson, The free Markoff field, \emph{J. Functional Analysis} 12 (1973), 211--227.

\bibitem{OwhadiYoo} H. Owhadi and G. R. Yoo, Kernel Flows: from learning kernels from data into the abyss, \emph{J. Comput. Phys.} 389 (2019), 22--47.
\bibitem{Patil} P. Patil, Y. Wu and R. J. Tibshirani, Failures and successes of cross-validation
for early-stopped gradient descent, in \emph{Proceedings of the 27th International Conference on
Artificial Intelligence and Statistics}, PMLR 238 (2024), 2260--2268. arXiv:2402.16793.

\bibitem{PN24} ProofNet24 run reports, tracks P15 (kernel mean correction), P19 (kernel spectral
pruning), P20 (kernel hard edge), P21 (correlated Gaussian equivalence), P22 (benign overfitting) and
P23 (kernel obstruction), repository \texttt{yspennstate/math-open-problems}, directory
\texttt{proofnet24}, 2026. Private; available from the author on request.

\bibitem{RakhlinZhai} A. Rakhlin and X. Zhai, Consistency of interpolation with Laplace kernels is a
high-dimensional phenomenon, in \emph{Proceedings of the 32nd Conference on Learning Theory}, PMLR
99 (2019), 2595--2623.


\bibitem{Rosenthal} H. P. Rosenthal, On the subspaces of $L^p$ ($p>2$) spanned by sequences of independent random variables, \emph{Israel J. Math.} 8 (1970), 273--303.
\bibitem{Shmalo26a} Y. Shmalo, Neural means and kernel corrections for radiative-transfer
emulators, manuscript, code and records at
\url{https://github.com/yspennstate/emit-oco2-kernel-emulators}, 2026.

\bibitem{Shmalo26b} Y. Shmalo, Neural means and kernel corrections for operator learning,
arXiv:2609.00389 (2026); code at \url{https://github.com/yspennstate/neural-means-kernel-corrections}.


\bibitem{Susiluoto} J. Susiluoto, A. Spantini, H. Haario, T. H\"ark\"onen and Y. Marzouk, Efficient multi-scale Gaussian process regression for massive remote sensing data with satGP v0.1.2, \emph{Geosci. Model Dev.} 13 (2020), 3439--3463.
\bibitem{PDEBench} M. Takamoto, T. Praditia, R. Leiteritz, D. MacKinlay, F. Alesiani, D. Pfl\"uger
and M. Niepert, PDEBench: an extensive benchmark for scientific machine learning, in \emph{Advances
in Neural Information Processing Systems} 35, Datasets and Benchmarks Track (2022).
arXiv:2210.07182.

\bibitem{Talagrand} M. Talagrand, Concentration of measure and isoperimetric inequalities in
product spaces, \emph{Publ. Math. Inst. Hautes \'Etudes Sci.} 81 (1995), 73--205.

\bibitem{TracyWidom} C. A. Tracy and H. Widom, Level spacing distributions and the Bessel kernel,
\emph{Comm. Math. Phys.} 161 (1994), 289--309.


\bibitem{Tropp12} J. A. Tropp, User-friendly tail bounds for sums of random matrices, \emph{Found. Comput. Math.} 12 (2012), 389--434.

\bibitem{Tropp15} J. A. Tropp, \emph{An Introduction to Matrix Concentration Inequalities}, Foundations and Trends in Machine Learning 8, 2015.

\bibitem{Vershynin} R. Vershynin, \emph{High-Dimensional Probability: An Introduction with Applications in Data Science}, Cambridge University Press, 2018.
\bibitem{Wigner} E. P. Wigner, On the distribution of the roots of certain symmetric matrices,
\emph{Ann. of Math.} 67 (1958), 325--327.

\end{thebibliography}

\end{document}
